% a01_H0rel: proof of prop:k01:H0rel ((H0rel)_C on King's cube), appendix
% local macros (providecommand: harmless if already defined upstream)
\providecommand{\Hrel}{\ensuremath{(\mathrm{H0rel})_{C}}}
\providecommand{\Rzero}{\ensuremath{(\mathrm{R}_{0})}}
\providecommand{\BALint}{\ensuremath{(\mathrm{BAL}_{5})[\mathrm{int}]}}
\providecommand{\RENorb}{\ensuremath{(\mathrm{REN}_{\mathrm{orb}})}}
\providecommand{\Zk}{\ensuremath{Z_{k}}}
\providecommand{\Ik}[1]{\ensuremath{I^{(#1)}_{k}}}
\providecommand{\Xl}{\ensuremath{X_{\mathrm{live}}}}       % the live scalar sites of the cube at a scale
\providecommand{\Kill}{\ensuremath{\mathsf K}}             % the killed sites
\providecommand{\Wt}{\ensuremath{W}}                       % the f_D sup weight (5/4)^3
\providecommand{\Mop}{\ensuremath{M}}                      % the cube operator -Delta^D_A+V+aP
\providecommand{\vth}{\ensuremath{\vartheta}}              % the fine spacing in block units, 5^{-j}
\providecommand{\DD}{\ensuremath{\partial D}}              % the Dirichlet (killed) planes
\providecommand{\dN}{\ensuremath{\partial_{N}}}            % the Neumann part of a sub-domain's boundary
\providecommand{\dep}{\operatorname{depth}}
\providecommand{\Ebar}{\ensuremath{\mathsf E}}             % even extension
\providecommand{\Lone}{\ensuremath{L_{1}}}                 % the face layer of blocks

\section{Proof of Proposition~\ref{prop:k01:H0rel}: the small-field representation on King's cube}\label{app:a01_H0rel}

This appendix proves \cref{prop:k01:H0rel}, including its clause \Rzero. The proof is King's proof of \cite[Thm.~3.1]{K1}, which is Ba{\l}aban's lower and upper bounds \cite{Ba1,Ba2} with the modifications of \cite[Appendix]{K1}, run on King's RG cube $\CK$ (\cref{def:s02:Kcube}) with King's block-spin flow of \cref{ssec:k01:setting}. Two things change relative to the torus. First, the scalar field is killed on the Dirichlet planes of the cube and averaged with King's Dirichlet weight; every place where Ba{\l}aban's or King's arguments use that the block average is uniform ($QQ^*=1$, $P$ a projector, the constant is reproduced by the block average) is replaced by the corresponding fact for the weighted average (\cref{app:a01:weights}), and every place where the torus has no boundary acquires a Dirichlet face, at which the scalar background vanishes linearly (\cref{app:a01:ops}). Second, the photon has Neumann conditions; every photon object of the cube is the even restriction of the same object on a doubled torus, so every photon estimate of the torus analysis holds on the cube with its constant multiplied by at most $8$ (\cref{app:a01:photon}). The large-field analysis near the Dirichlet faces uses King's face test \cite[(5.41)]{K2} as a large-field family at every scale, paid by the Dirichlet-face positivity of \cite[Prop.~5.4]{K2} in an elementary form valid for every gauge field (\cref{app:a01:face}). The lower bound (\cref{app:a01:lower}) carries King's Appendix class with a background-gradient clause and a block-mismatch clause in place of King's Laplacian clause, and the upper bound (\cref{app:a01:upper}) runs Ba{\l}aban's large-field decomposition on the cube; the assembly, the normalisation-free form of the remainder and the clause \Rzero\ are in \cref{app:a01:assembly}.

Throughout, $d=3$, $L=5$, the couplings $(m^2,\lmu,e,\lambda)$ with $m^2,\lmu,\lambda>0$ are fixed, as are $J\ge J_0$ and the source $s\in\{u_x,0\}$ (so $h=0$: the scalar source is treated in \cref{sec:k02}). Constants $C,c,\delta$ change from line to line and depend only on the quantities listed where they are introduced; they never depend on the spacing index $k$, the step $j$, the field or the domain unless this is said.

% a01_H0rel part 1: setting on the cube, the weight class, the photon sector
\subsection{The cube in block units, and the weighted block average}\label{app:a01:weights}

\paragraph{Block units.} Fix the spacing index $k$ and a step $j$, $0\le j\le k$, and write everything in \emph{unit-$j$ units}: the $j$-blocks (cubes of side $5^{j}\eps_k=\eta$ in physical units) have side one, and the fine lattice has spacing $\vth\defeq5^{-j}$, with $1/\vth\in5\N$ for $j\ge1$. The number of blocks per side is $N\defeq2MJ/\eta\in\N$ (the side $2MJ$ of $\CJ$ is a multiple of every block side, \cref{ssec:k01:setting}). After a translation the picture is: fine sites $\Lambda\defeq\vth\Z^3$, the \emph{killed set} $\Kill\defeq\{x\in\Lambda\cap[0,N]^3:\ x_\mu\in\{0,N\}\text{ for some }\mu\}$ (the two Dirichlet planes of $\CK$ in each direction, $2MJ$ apart), the \emph{live set} $\Xl\defeq\{x\in\Lambda:\ 0<x_\mu<N\ \forall\mu\}$, which is $\alive^K_J$ of \cref{def:s02:Kcube} ($N/\vth-1$ sites per side), the blocks $B(y)\defeq(y+[0,1)^3)\cap\Lambda$, $y\in\Z^3$, $0\le y_\mu\le N-1$, and the \emph{face layer} $\Lone$ of blocks with some $y_\mu\in\{0,N-1\}$. The bottom Dirichlet plane $x_\mu=0$ is the first site plane of the blocks with $y_\mu=0$ (inside those blocks); the top plane $x_\mu=N$ is the first site plane of the blocks $y_\mu=N$, which lie outside the cube. We write $\DD$ for the union of the Dirichlet planes, $t(x)\defeq\min_\mu\min(x_\mu,N-x_\mu)$ for the face distance of a live site ($t=\vth$ on the first live layer $F_1$, the live sites with a neighbour in $\Kill$), and a \emph{face bond} is a bond from a live site to a killed site. By \cref{def:s02:acube}, a face bond carries only the diagonal term $\frac12\vth^{d-2}|\phi(x)|^2$ at its live end: it has no link variable, and $D_A\phi(b)=-\vth^{-1}\phi(x)$ on it whatever $A$. The photon is a site field on $\XJ$ with Neumann conditions on the bond planes $x_\mu=-\vth/2$, $N-\vth/2$ (\cref{def:s02:Kcube}).

\paragraph{The weight.} King's Dirichlet weight \cite[(3.8)--(3.9)]{K2}, in the form in which the averaging blocks have their top layers on the Dirichlet planes (\cref{ssec:k01:setting}), is the function
\[
w(x)\ \defeq\ \Wt\prod_{\mu=1}^{3}\mathbf 1\big[x_\mu/\vth\not\equiv0\ (\mathrm{mod}\ 5)\big],\qquad \Wt\defeq(5/4)^3=125/64,
\]
on the fine sites of a block: zero on the three site planes of the block that lie in the residue class of the Dirichlet planes (the \emph{zero-weight layers}), the constant $\Wt$ elsewhere. Since $1/\vth\equiv0\pmod5$, the block corners and all of $\Kill$ lie in the zero-weight class, and $\sum_{x\in B(y)}\vth^3w(x)=1$ for every block of the cube. The weight has period $5\vth$, so it acts at the finest step only: the coarse averages above scale one are uniform, $Q_{j+1}=Q\,Q_j$ with $Q$ the uniform one-step average with $QQ^*=1$, and the constants $a_j$ of \cite[(2.13)]{K1} are those of the uniform flow. The block average and its adjoint are
\[
Q(A)\phi(y)\ \defeq\ \sum_{x\in B(y)\cap\Xl}\vth^3\,w(x)\,U(A(\Gamma_{y,x}))\,\phi(x),\qquad P(A)\defeq Q(A)^*Q(A),
\]
with King's hierarchical contours $\Gamma_{y,x}$ \cite[(2.10)--(2.12)]{K1}, chosen coordinate-monotone between the block centre and $x$; such a path between two live sites of one block meets no killed site, since the killed sites of a block sit at the extreme value of a coordinate. The adjoint is taken for the measure $\vth^3$ on fine sites and counting measure on blocks.

\begin{definition}[weight class]\label{def:a01:wclass}
A block weight is a function $w\ge0$ on the fine sites of a unit block with (W1) $\sum_{x\in B}\vth^3w(x)=1$; (W2) $w\le\Wt$; (W3) $w=0$ on every killed site of the block; (W4) the same $w$ on every block up to translation. The uniform weight ($w\equiv1$, $\Wt=1$) and King's weight above ($\Wt=125/64$) are in the class.
\end{definition}

\begin{lemma}[the weighted average]\label{lem:a01:WQ}
For every weight in the class and every gauge field $A$, with $W_2\defeq\sum_{x\in B}\vth^3w(x)^2\in[1,\Wt]$ (for King's weight $W_2=\Wt$):
\begin{enumerate}[label=(Q\arabic*),nosep]
\item $Q(0)\mathbf 1=\mathbf 1$ on every block, and $Q(A)$ maps the transported constant to the constant;
\item $Q(A)Q(A)^*=W_2\,\mathbf 1$;
\item $P=W_2\Pi$ with $\Pi$ an orthogonal projector; $0\le P\le W_2$ and $\norm{1-P}=\max(1,|1-W_2|)\le1$ since $W_2\le\Wt<2$;
\item $|(Qu)(y)|\le W_2^{1/2}\norm{\mathbf 1_{B(y)}u}_2$ and $\norm{Q^*e_y}_2=W_2^{1/2}$;
\item $|(Qu)(y)|\le\sup_{B(y)}|u|$;
\item $|Q^*b|_\infty\le\Wt|b|_\infty$, $\norm{P}_{\infty\to\infty}\le\Wt$, and the Schur norm $\sup_x\sum_{x'}\vth^3|P(x,x')|\le\Wt$;
\item $\sum_x\vth^3|(Pu)(x)|=\sum_y|(Qu)(y)|$;
\item $P(x,x')=0$ unless $x,x'$ lie in one block, and $P$ vanishes on killed sites: $P$ never couples a live site to a killed one, and every union of blocks is invariant;
\item under $\phi\mapsto e^{eq\lambda}\phi$, $A\mapsto A+d\lambda$, the operators $Q(A)$, $P(A)$ transform covariantly.
\end{enumerate}
\end{lemma}
\begin{proof}
(Q1), (Q2): (W1) and the orthogonality of the transports $U$. (Q3): $P^2=Q^*(QQ^*)Q=W_2P$, and $1-P$ has the eigenvalues $1$ on $\Pi^\perp$ and $1-W_2$ on $\Pi$. (Q4): Cauchy--Schwarz with the weight $w$ and (W1); $\norm{Q^*e_y}^2=\sum\vth^3w^2$. (Q5): $|Qu|\le\sum\vth^3w|u|\le\sup|u|$ by (W1). (Q6): $(Q^*b)(x)=w(x)U^*b(y)$ and $(Pu)(x)=w(x)\sum_{x'}\vth^3w(x')U\,u(x')$. (Q7): $|(Pu)(x)|=w(x)|(Qu)(y)|$ and (W1). (Q8): the definitions and (W3). (Q9): the contour transports transform with the endpoint gauge functions, and $w$ is a scalar.
\end{proof}

Where a step of Ba{\l}aban's or King's torus analysis uses $QQ^*=1$, $P^2=P$, or that the block average of a constant is that constant, the cube version below uses the corresponding item of \cref{lem:a01:WQ} and acquires a factor depending on $\Wt$ only. Three identities change in substance and are proved below: the background of a constant datum (\cref{lem:a01:RC}), the block-mismatch identity (\cref{lem:a01:EF}), and the minimiser equation with $QQ^*=W_2$ (used in \cref{lem:a01:EB}).

\paragraph{The operators.} In unit-$j$ units let $V_j\defeq m^2\eta^2\le c_V$ be the mass term, $A$ a real gauge field on the live bonds of $\Xl$ with $U(A_b)=\exp(\vth\,e_jqA_b)$, $e_j\defeq e\eta^{1/2}$, and
\[
\Mop(\Omega,A)\ \defeq\ -\Delta^{N,D}_{A,\Omega}+V_j+a_jP(A)\qquad\text{on }\ell^2(\Omega\cap\Xl),
\]
for a \emph{sub-domain} $\Omega$ (a union of $j$-blocks of the cube; $\Omega=\Xl$ is the whole cube): the Laplacian keeps every bond with both ends in $\Omega\cap\Xl$, adds $\vth^{-2}$ to the diagonal for each face bond (Dirichlet, no link), and on bonds from $\Omega\cap\Xl$ to $\Xl\setminus\Omega$ either drops them (Neumann cut) or adds the diagonal (Dirichlet cut). The \emph{Neumann boundary} $\dN\Omega$ is the set of block midplanes between a block of $\Omega$ and a block of the cube not in $\Omega$; the \emph{depth} of a point is $\dep(x)\defeq\dist_\infty(x,\dN\Omega)$ ($+\infty$ on the whole cube). The Dirichlet faces of the cube are part of the operator, never part of $\dN\Omega$: every block touching a Dirichlet plane contains that plane's killed sites (bottom) or ends one site before it (top), so no sub-domain has a Neumann face on a Dirichlet plane. The \emph{leg} of block data $b$ is $u\defeq a_j\Mop(\Omega,A)^{-1}Q(A)^*b$; the \emph{background} of the block field $\phi_j$ is $\phi^{(j)}\defeq a_j\Mop^{-1}Q^*\phi_j$ \cite[(2.14)--(2.16)]{K1}; the block form is $\Delta^{(j)}\defeq a_j-a_j^2Q\Mop^{-1}Q^*$ and the fluctuation covariance $C^{(j)}\defeq(\Delta^{(j)}+aL^{-2}Q_c^*Q_c)^{-1}$ with $Q_c$ the uniform coarse average. A field $A$ is \emph{regular} on $\Omega$ at scale $j$ if it obeys \cite[(3.5)]{K1}: $|\partial_\mu A_\nu|\le K_j$ with $\theta_j\defeq e_jK_j\le\theta_0$ small; King's small-field class gives $\theta_j\le Ce_jp_j$, small for $\eta\le\eta_1$. Gauge covariance on the cube: for any real $\lambda$ on the sites, $\Mop(\Omega,A+d\lambda)=U_\lambda\Mop(\Omega,A)U_\lambda^*$ with $U_\lambda$ the multiplication by $e^{eq\lambda}$, $Q(A+d\lambda)=U^c_\lambda Q(A)U_\lambda^*$, and on a face bond $|D_{A+d\lambda}(U_\lambda u)|=|D_Au|=\vth^{-1}|u(x)|$; so every bound at a constant field $B=d(B\cdot x)$ on a box equals the bound at $0$ for a datum of the same modulus (the contours lie in the block, the Dirichlet condition is multiplication-invariant, the weight is a scalar).

\paragraph{The small-field set and the face test.} The small-field indicator $\chi_j$ of \cref{ssec:k01:setting} is \cite[(3.2)]{K1} on the blocks of the cube, with the gradient clause $|D_{A_j}\phi_j(\beta)|\le p_{j-1}$ on every block bond $\beta$ of the cube (all of which are cube-internal), times King's face test \cite[(5.41)]{K2}
\[
\lambda_j\ \defeq\ \prod_{y'\in\Lone}\chi\big(|\phi_j(y')|\le\sigma_j\big),\qquad \sigma_j\defeq\Big(\frac{2}{\theta_2\gamma_D}\Big)^{1/2}p_j\ \le\ C_8\,p_j,\quad \gamma_D\defeq\frac{a_j}{1+a_j\Wt^2},
\]
with $\theta_2\in(0,1)$ a fixed share and $\gamma_D$ the Dirichlet-face positivity constant of \cref{lem:a01:FPC}; $p_j\defeq p(\eta)=b_0(1+\ln\eta^{-1})^{p}$ and $M_j\defeq p_{j-1}(\lambda\ell_{j-1})^{-1/4}$ is the block-sup threshold of \cite[(3.2)]{K1} ($\ell_j\defeq5^j\eps_k$ the physical block side, $\lambda(\ell_j)=\lambda\ell_j$ the running coupling in $d=3$). The constant $C_8$ depends on $(a,\Wt,\theta_2,L)$ only.

\subsection{The photon sector: even restriction of the doubled torus}\label{app:a01:photon}

Let $2X$ be the torus of period $2N$ in each direction (physical period $4MJ$), $\XJ$ embedded as $\{0,\dots,N-\vth\}^3$ after the translation above, and $\mathcal W\cong(\Z_2)^3$ the group generated by the reflections $r_\mu$ in the photon planes $x_\mu=-\vth/2$ (on $2X$ the reflection in $x_\mu=N-\vth/2$ is the same element). Since $4MJ\equiv0$ modulo every block side, the blocks of $\XJ$ are blocks of the periodic tiling of $2X$, and $\mathcal W$ maps blocks to blocks. For a photon component $A_\mu$ (a site field) let $\Ebar$ be the even extension to $2X$, and $\Ebar_c$ the same on the coarse lattices.

\begin{lemma}[photon objects are even restrictions]\label{lem:a01:PE}
For every photon object $O$ of Ba{\l}aban's and King's construction built from the Neumann form $-\Delta^N_{\XJ}+\lmu^2\eta^2$ and the uniform block averages on $\XJ$, and its counterpart $O_{2X}$ on the periodic torus $2X$:
\begin{enumerate}[label=(\alph*),nosep]
\item the propagators $G_j$, the block forms $\Delta^{(j)}$, the fluctuation covariances $C^{(j)}$, the background map $A_j\mapsto A^{(j)}=a_jG_jQ_j^*A_j$ and the slices satisfy $O_{2X}\Ebar=\Ebar\,O$;
\item quadratic forms: $\ip{\Ebar A}{\Delta^{(j)}_{2X}\Ebar A}=8\ip{A}{\Delta^{(j)}A}$, and likewise for every positive form above;
\item for every coarse $\Lambda\subset X^{(j)}$ (resp.\ every fine union of blocks $\Omega\subset\XJ$) the Dirichlet-conditioned compressed operators on $\Lambda$ and the Neumann operators on $\Omega$ intertwine with $\Ebar$ and the $\mathcal W$-invariant sets $\mathcal W\Lambda$, $\mathcal W\Omega$;
\item if Ba{\l}aban's cut-off kernel $\zeta$ of \cite[(2.44)]{Ba2} is applied image-wise on $2X$ (the cube kernel being $K_X(x,y)\defeq\sum_{g\in\mathcal W}\zeta(x,gy)K_{2X}(x,gy)$), the cut-off objects of the cube are even restrictions of those of $2X$ and the support and size properties of $\zeta$ hold for them with the torus constants;
\item every positivity bound of the photon sector (\cite[(3.29), Prop.~3.1]{Ba2} photon part, the last-step mass bound $G'_k\ge\gamma_1\lmu^2\ell_k^2$ of \cite[(3.11)]{Ba2}, the photon exponents of \cite[(3.26)]{Ba2}) holds on $\XJ$ with the constant of $2X$.
\end{enumerate}
\end{lemma}
\begin{proof}
(a) The Neumann Laplacian on $\XJ$ is the periodic Laplacian of $2X$ on $\mathcal W$-even fields with reflection planes between sites: every plane-crossing bond of $2X$ joins a site to its mirror image, where an even field has zero difference, and every bond of $\XJ$ appears $8$ times. Block averages commute with $\Ebar$ (blocks go to blocks) and $Q^*_{2X}\Ebar_c=\Ebar Q^*_X$ (the adjoint of a $\mathcal W$-covariant uniform average is the uniform block extension). Hence $(M_{2X}+aQ^*Q)\Ebar=\Ebar(M_X+aQ^*Q)$; the torus operator commutes with $\mathcal W$, preserves the even subspace and is invertible on it, so $G_{2X}\Ebar=\Ebar G_X$, and $\Delta^{(j)}$, $C^{(j)}$, the background map and the slices follow from the algebra of \cite[(2.13)--(2.17)]{K1}, which is the same on both sides. (b) $\Ebar^*\Ebar=8$ on $\XJ$. (c) Compression to a $\mathcal W$-invariant set commutes with $\mathcal W$, and $P_{\mathcal W\Lambda}\Ebar=\Ebar P_\Lambda$; then (a) on the invariant subspace. For Neumann conditions on $\Omega$, the bonds of $\mathcal W\Omega$ are the images of those of $\Omega$ plus plane-crossing bonds, on which even fields have zero difference. (d) $\zeta(x,gy)$ depends on $|x-gy|$ only; the image-wise kernel commutes with $\mathcal W$ by construction and restricts to the cube kernel by definition; the support and size statements are statements about $\zeta(x,\cdot)$ on $2X$, and $|x-gy|\ge|x-y|$ for $x,y\in\XJ$ keeps every support statement. The cut is a device inside the definition of Ba{\l}aban's background, so the choice is free. (e) By (b): the cube form of an even field is $1/8$ of the torus form, and the right-hand sides (gradient sums, mass sums, block sums) scale by the same factor since plane-crossing gradients of even fields vanish.
\end{proof}

\begin{lemma}[regularity is preserved]\label{lem:a01:REG}
(a) A block photon configuration $A_j$ on $\XJ$ satisfies the photon clauses of \cite[(3.2)]{K1} if and only if $\Ebar_cA_j$ does on $2X$. (b) If $A$ on $\XJ$ is regular in the sense of \cite[(3.5)]{K1}, or obeys a H\"older bound with constant $C$, so does $\Ebar A$ on $2X$ with the same constant.
\end{lemma}
\begin{proof}
(a) Plane-crossing coarse gradients of $\Ebar_cA_j$ vanish; the other clauses are copies. (b) $\Ebar A=A\circ F$ with $F$ the coordinatewise folding map $2X\to\XJ$, which is $1$-Lipschitz; lattice differences across a plane vanish, the others are copies.
\end{proof}

\begin{lemma}[Neumann photon lines are image sums]\label{lem:a01:NL}
Every photon line of the cube (the fluctuation covariances and slices of \cref{lem:a01:PE}(a) and the full covariance $(-\Delta^N+\lmu^2\eta^2)^{-1}$) satisfies $K_X(x,y)=\sum_{g\in\mathcal W}K_{2X}(x,gy)$ for $x,y\in\XJ$. Hence every line bound of the torus analysis (decay, pointwise and derivative clauses, H\"older clauses in either variable, moment forms) holds for $K_X$ with the constant multiplied by $8$ and the same rates.
\end{lemma}
\begin{proof}
Kernel form of \cref{lem:a01:PE}(a): $(K_Xf)(x)=(K_{2X}\Ebar f)(x)=\sum_y\sum_gK_{2X}(x,gy)f(y)$. Each image term obeys the torus bound at the pair $(x,gy)$, and $|x-gy|\ge|x-y|$; derivatives in $x$ act on $K_{2X}(\cdot,gy)$, derivatives in $y$ through the lattice isometry $y\mapsto gy$.
\end{proof}

Consequently every photon statement that the torus analysis takes from \cite{Ba1,Ba2,Ba3} holds on $\XJ$: apply it on the torus $2X$ (whose period $4MJ$ is a multiple of every large-block side, and whose sides exceed the diameter of every box used) to the even data and restrict. For $x\in\XJ$, $\dist_{2X}(x,\partial(\mathcal W\Omega))\ge\dist(x,\partial\Omega\setminus\partial\XJ)$, so the restricted-form clauses transfer, and at the faces of $\XJ$ the photon has no boundary at all in the doubled picture. The photon tadpole and exchange graphs of the counterterm that contain a Neumann line $C_N-C_{\mathrm{free}}=\sum_{g\ne1}C_{\mathrm{free}}(\cdot,g\,\cdot)$ are bounded by the same image argument: with $\varrho(x)$ the distance of $x$ to the nearest photon plane, $|x-gx|\ge2\varrho(x)$ for $g\ne1$ and $|C_{\mathrm{free}}(x,z)|\le C\min(\vth^{-1},|x-z|^{-1})e^{-\lmu|x-z|}$, so the photon-image part of the mass counterterm is $O\big((e^2+e^4)(\min(\vth^{-1},\varrho^{-1})+1)\big)$ at order $\le4$, the same size as the scalar boundary remainders near the face, and it is absorbed in the zone-B bounds of \cref{prop:a01:TD} (it is the photon-image part of the boundary remainder families there).

% a01_H0rel part 2: the operator layer on the cube and its sub-domains
\subsection{Operator bounds on the cube and its sub-domains}\label{app:a01:ops}

The statements of this subsection are \cite[Thm.~3.3]{K1} and Ba{\l}aban's propagator bounds \cite[Sects.~1--2]{Ba4} for the cube operator $\Mop(\Omega,A)$ of \cref{app:a01:weights}, on every sub-domain $\Omega$ (any union of $j$-blocks of the cube, meeting the Dirichlet planes or not, re-entrant edges allowed, Neumann or Dirichlet cuts on $\dN\Omega$), at every scale $1\le j\le K_f$, uniformly in $\vth$, $J$, $k$, the shape of $\Omega$ and the regular field $A$. At $j=0$ (fixed spacing $5^{-1}$) every clause holds with $L$-dependent constants by finite dimensionality. The constants depend on $(d,L,a,\Wt,c_V)$ and King's constants only.

\begin{lemma}[coercivity]\label{lem:a01:A1}
There are $\gamma=\gamma(d,a,\Wt)>0$ and $\theta_0>0$ such that $\Mop(\Omega,A)\ge\gamma$ for every $\Omega$, $\vth\le1/5$, $V_j\ge0$ and $A$ regular with $\theta_j\le\theta_0$.
\end{lemma}
\begin{proof}
Dropping the bonds between distinct blocks lowers the form (each dropped term is $\ge0$; Dirichlet terms only add), and $P(A)$ is block diagonal (Q8), so $\Mop\ge\bigoplus_B\Mop_B$ with $\Mop_B\defeq-\Delta^N_{A,B}+a_jP_B(A)$ on $\ell^2(B\cap\Xl)$, the killed sites of a face block entering as a Dirichlet term. At $A=0$ write $\phi=c+\psi$ on $B$ with $\psi$ orthogonal to constants (after extending $\phi$ by $0$ to the killed sites, which does not change $P\phi$ by (W3) and only raises the gradient form). The Neumann gap of the unit cube gives $\ip\psi{-\Delta^N\psi}\ge g_P\norm\psi^2$ with $g_P=4\vth^{-2}\sin^2(\pi\vth/2)\ge8$ for $\vth\le\frac12$; by (Q1), (Q4), $|Q\phi|=|c+Q\psi|\ge|c|-\Wt^{1/2}\norm\psi$. Hence $\ip\phi{\Mop_B\phi}\ge g_P\norm\psi^2+a(1-L^{-2})(|c|-\Wt^{1/2}\norm\psi)_+^2\ge\gamma_*(|c|^2+\norm\psi^2)$ with $\gamma_*=\gamma_*(a,\Wt)>0$ by the springs-in-series minimisation $\min_{s\ge0}[\Wt^{-1}s^2+a'(F_0-s)_+^2]=F_0^2/(a'^{-1}+\Wt)$; explicitly $\gamma\ge(1-L^{-2})\min\{a/2,8/(1+2\Wt)\}$. For regular $A$: on the simply connected block, $A=A(x_B)+A''$ with $A(x_B)$ constant, pure gauge, removed by a unitary gauge transformation that preserves both terms, and $|A''|\le\sqrt3K_j$; then $\norm{D_{A''}\phi-D_0\phi}\le2\theta_j\norm\phi$ and $|Q(A'')\phi-Q(0)\phi|\le3\sqrt3\,\Wt^{1/2}\theta_j\norm\phi$ (contours of length $\le3$), so $\ip\phi{\Mop_B\phi}\ge(\gamma_*/2-C\theta_j)\norm\phi^2$, which is the claim for $\theta_j\le\theta_0$.
\end{proof}

\begin{lemma}[Combes--Thomas; unit-lattice kernels]\label{lem:a01:A23}
(a) For every $1$-Lipschitz $g$, $\norm{e^{\delta_0g}\Mop^{-1}e^{-\delta_0g}}_{2\to2}\le2/\gamma$, hence $\norm{\mathbf 1_X\Mop^{-1}\mathbf 1_Y}\le(2/\gamma)e^{-\delta_0\dist(X,Y)}$, with $\delta_0=\delta_0(\gamma,a,\Wt,d)$. (b) $|\Delta^{(j)}(\Omega,A)(u,v)|\le Ce^{-\delta_0|u-v|}$. (c) If $\Omega$ is a union of $(j+1)$-blocks, then $\gamma_0''\le\Delta^{(j)}+aL^{-2}Q_c^*Q_c\le a_j+a$ on the blocks of $\Omega$, with $\gamma_0''=\gamma_0''(d,L,a,\Wt)>0$ for $A=0$ and, at regular $A$ with $\theta_j\le\theta_0'$, half of it; consequently $|C^{(j)}(u,v)|\le Ce^{-\delta_1|u-v|}$.
\end{lemma}
\begin{proof}
(a) Conjugation multiplies the hopping entry $(x,x')$ by $e^{\delta(g(x)-g(x'))}$, $|g(x)-g(x')|\le\vth$; the first-order part is anti-Hermitian and drops from $\Re\ip\phi{\Mop_\delta\phi}$, the remainder has norm $\le3d\delta^2$; the block term changes by $\le(e^{\delta\sqrt3}-1)\Wt\le2\sqrt3\,\Wt\delta$ (its Schur norm, (Q6)); Dirichlet terms are diagonal and commute. So $\Re\Mop_\delta\ge\gamma-3d\delta^2-2\sqrt3\Wt a\delta\ge\gamma/2$ for $\delta\le\delta_0$. (b) $\Delta^{(j)}(u,v)=a_j\delta_{uv}-a_j^2\ip{Q^*e_u}{\Mop^{-1}Q^*e_v}$, $\norm{Q^*e_u}=W_2^{1/2}$ (Q4), and (a). (c) The upper bound is $\Delta^{(j)}\le a_j$, from the minimiser identity $\ip\psi{\Delta^{(j)}\psi}=\min_\Phi[\ip\Phi{H\Phi}+a_j\norm{\psi-Q\Phi}^2]$, $H\defeq-\Delta_A+V_j$ (take $\Phi=0$). Lower bound: dropping the fine bonds between distinct $(j+1)$-blocks $Y$ lowers $H$ and hence $\Delta^{(j)}$ (monotonicity of the minimiser identity), and $Q_c^*Q_c$ is $(j+1)$-block diagonal, so it suffices to bound $F_Y(\psi)\defeq\ip\psi{(\Delta^{(j)}_Y+aL^{-2}Q_c^*Q_c)\psi}$ from below on the $L^3$ coarse sites of one $(j+1)$-block $Y$ (Neumann on its faces, Dirichlet at killed sites). $Q_c^*Q_c$ is the projector onto constants on $Y$; write $\psi=c+\psi_\perp$. Since $Q\mathbf 1=\mathbf 1$ (Q1), $\norm{\psi-Q\Phi}^2\ge\norm{\psi_\perp-(Q\Phi)_\perp}^2\ge(\norm{\psi_\perp}-\norm{Q(\Phi-\av\Phi)})_+^2\ge(\norm{\psi_\perp}-(\Wt C_Ps)^{1/2})_+^2$ when $\ip\Phi{-\Delta^N_Y\Phi}\le s$, by (Q4) and the Neumann Poincar\'e inequality on the live sites of $Y$ (a cube of side $L$, possibly minus a killed face layer, whose Dirichlet bonds only add; constant $C_P(L)$ uniform in $\vth$). Springs in series: $\Delta^{(j)}_Y(\psi)\ge\norm{\psi_\perp}^2/(\Wt C_P+1/a_j)$, and the coarse term is $aL^{-2}|c|^2L^3$; so $F_Y\ge\min\{(\Wt C_P+1/a)^{-1},aL^{-2}\}\norm\psi^2$ at $A=0$. At regular $A$ gauge away the constant part on $Y$ (diameter $\sqrt3L$) at cost $C(L)\theta_j$. The decay of $C^{(j)}$ then follows from the lower bound and (b) by Combes--Thomas for exponentially decaying unit-lattice kernels (the conjugation argument of (a) with the Schur bound $\norm{e^{tg}Xe^{-tg}-X}\le Ct\sum_{z\in\Z^3}|z|e^{-\delta_0|z|/2}$), which is \cite[(5.6)--(5.7)]{Ba4}.
\end{proof}

\begin{lemma}[Kato, and sup bounds with decay]\label{lem:a01:A4}
Let $u=\Mop(\Omega,A)^{-1}f$. (a) $(-\Delta^{N,D}_{0,\Omega}+V_j)|u|\le|f-a_jP(A)u|$ pointwise on $\Omega\cap\Xl$, for every $A$ and $V_j\ge0$ (the free Laplacian with the same cuts and the same Dirichlet terms). (b) For block data $b$ on the blocks of $\Omega$ and the leg $u=a_j\Mop^{-1}Q^*b$,
\[
|u(x)|\le C\sum_{v}e^{-\delta_2\dist(x,B(v))}|b(v)|\le C|b|_\infty,\qquad |(Qu)(v)|\le C\sum_{v'}e^{-\delta_0|v-v'|}|b(v')|,
\]
at every point of every $\Omega$ and every $A$ (given \cref{lem:a01:A1}). (c) The same with $Q^*b$ replaced by any $f$ with block-$L^2$ data $\norm{f\mathbf 1_{B(v)}}_2\le|b(v)|$; in particular $\norm{\Mop^{-1}f}_\infty\le C\norm f_\infty$ and $|\Mop^{-1}f(x)|\le Ce^{-\delta_2\dist(x,\supp f)/2}\norm f_\infty$ for bounded $f$.
\end{lemma}
\begin{proof}
(a) At $u(x)=0$ the left side is $\le0$. For $u(x)\ne0$, $\Re[\bar u(x)(u(x)-U_{xx'}u(x'))]/|u(x)|\ge|u(x)|-|u(x')|$ since $U_{xx'}$ is orthogonal; a face bond contributes $\vth^{-2}|u(x)|$ on both sides; sum over the kept bonds and use $(-\Delta_A+V_j)u=f-a_jPu$. (b), (c) Block averages first: $\norm{\mathbf 1_{B(v)}u}_2\le(2a/\gamma)\sum_{v'}e^{-\delta_0\dist}\norm{\mathbf 1_{B(v')}f}_2$ by \cref{lem:a01:A23}(a), and $|(Qu)(v)|\le\Wt^{1/2}\norm{\mathbf 1_{B(v)}u}_2$ (Q4). Then with $\Mop_1\defeq-\Delta^{N,D}_{0,\Omega}+V_j+a_j$ write $u=\Mop_1^{-1}[f+a_j(1-P)u]$ in the sense of (a): $|u|\le\Mop_1^{-1}|f+a_j(1-P)u|$ by inverse positivity of the $M$-matrix $\Mop_1$, with $\norm{\mathbf 1_B(1-P)u}_2\le\norm{\mathbf 1_Bu}_2$ by (Q3). The kernel of $\Mop_1^{-1}$ is dominated entrywise by that of the Neumann Yukawa operator $(-\Delta^N_{\Omega'}+a_j)^{-1}$ on the union $\Omega'$ of the blocks of $\Omega$ and of the blocks adjacent to its top face blocks (so that the killed sites adjacent to $\Omega\cap\Xl$ are sites of $\Omega'$): $\Mop_1$ is a principal compression of that $M$-matrix plus a non-negative diagonal, and for non-singular $M$-matrices the inverse of a principal compression is dominated entrywise by the compression of the inverse (Schur complement, $(M^{-1})_{SS}=(M_{SS}-M_{SR}M_{RR}^{-1}M_{RS})^{-1}$ with $M_{SR}M_{RR}^{-1}M_{RS}\ge0$ entrywise). \Cref{lem:a01:A4N} gives the block summability $\sup_x\sum_Be^{\delta_2'\dist(x,B)}\norm{K^N_a(x,\cdot)}_{L^2(B)}\le C$ of the Neumann Yukawa kernel on every union of blocks, uniformly in $\vth$; combining the two decays through $\dist(x,B')\le\dist(x,B)+\sqrt3+\dist(B,B')$ gives (b) and (c) with $\delta_2\le\delta_0$. For bounded $f$ only the blocks meeting $\supp f$ carry data.
\end{proof}

\begin{lemma}[block summability of the Neumann Yukawa kernel]\label{lem:a01:A4N}
Let $\Omega'$ be a union of unit blocks of $\vth\Z^3$ (any shape) with Neumann conditions and $K^N_a\defeq(-\Delta^N_{\Omega'}+a_j)^{-1}$, with kernel for the measure $\vth^3$. There are $C_Y',\delta_2'>0$ depending on $a$ only such that, uniformly in $\vth\le1$ and $\Omega'$, $\sup_x\sum_Be^{\delta_2'\dist(x,B)}\norm{K^N_a(x,\cdot)}_{L^2(B)}\le C_Y'$ and $K^N_a(x,y)\le C_Y'e^{-2\delta_2'|x-y|}$ for $|x-y|\ge\frac12$.
\end{lemma}
\begin{proof}
$-\Delta^N_{\Omega'}$ generates the reflecting continuous-time simple random walk on $\Omega'$, so $K^N_a(x,y)=\int_0^\infty e^{-a_jt}p_t(x,y)\,dt$ with $p_t$ symmetric and $\norm{K^N_a(x,\cdot)}_2^2=\int_0^\infty e^{-a_ju}u\,p_u(x,x)\,du$. \emph{On-diagonal bound.} On one unit cube $Q$ the discrete Nash inequality $\norm f_{2,Q}^{2+4/d}\le C_N(E_Q(f)+\norm f_{2,Q}^2)\norm f_{1,Q}^{4/d}$ holds with $C_N$ independent of $\vth$ (the discretisations of one fixed convex domain); it globalises to $\Omega'$ with no geometric hypothesis: with $a_Q\defeq\norm f_{2,Q}^2$, $e_Q\defeq E_Q(f)+a_Q$, $b_Q\defeq\norm f_{1,Q}$ the cube inequality reads $a_Q\le C_N^{d/(d+2)}e_Q^{d/(d+2)}b_Q^{4/(d+2)}$, and summing over the cubes with H\"older's inequality at the exponents $\frac{d+2}d,\frac{d+2}2$, then $\sum b_Q^2\le(\sum b_Q)^2$ and $\sum_QE_Q\le E_{\Omega'}$, gives the Nash inequality on $\ell^2(\Omega')$. By \cite[Thm.~2.1]{CKS87} this is equivalent to $p_u(x,x)\le Cu^{-d/2}$ for $u\le1$; trivially $p_u(x,x)\le\vth^{-d}$ and $p_u\le p_1$ for $u\ge1$. In $d=3$ the integral $\int_0^{\vth^2}\vth^{-3}u\,du+\int_{\vth^2}^1Cu^{-1/2}du+\int_1^\infty e^{-a_ju}Cu\,du$ is bounded, which bounds the near blocks. \emph{Far points}, $D\defeq|x-y|\ge\frac12$: for a $1$-Lipschitz $g$ the conjugated form satisfies $\Re\,e^{\theta g}(-\Delta^N_{\Omega'})e^{-\theta g}\ge-\lambda(\theta)$, $\lambda(\theta)\defeq2d\vth^{-2}(\cosh(\theta\vth)-1)\le d\theta^2\cosh(\theta\vth)$ (bond by bond, each site in $\le2d$ bonds), so $\norm{\mathbf 1_Fe^{t\Delta^N}\mathbf 1_E}\le\exp(-\theta D+d\theta^2t\cosh(\theta\vth))$ for sets at distance $D$, which is $\le e^{-D^2/(5dt)}$ for $t\ge\vth D$ and $\le e^{-D/(8d\vth)}$ for $t\le\vth D$. For $t\le\vth D$ this gives $p_t(x,y)\le\vth^{-d}e^{-D/(8d\vth)}$ and $\int_0^{\vth D}e^{-a_jt}p_t\,dt\le Ce^{-D/(16d)}$. For $\vth D\le t\le1$, Davies' method \cite{Dav93} (the twisted $L^p$ energy inequalities for the conjugated semigroup, with the Nash inequality above and $\lambda_p(\theta)\le C_p\theta^2\cosh(\theta\vth)$, give the $2\to\infty$ bound and duality the $1\to2$ bound) gives $p_t(x,y)\le Ct^{-d/2}e^{-\theta D+C_0dt\theta^2\cosh(\theta\vth)}$, and $\theta=D/(2C_0dt)$ gives $p_t\le Ct^{-d/2}e^{-cD^2/t}$, whose time integral is $\le Ce^{-cD^2/2}$. For $t\ge1$ write $p_t(x,y)=\ip{p_{1/2}(y,\cdot)}{e^{(t-1)\Delta^N}p_{1/2}(x,\cdot)}$, split each $p_{1/2}$ into its part within $D/3$ of its base point and the rest; the main term is $\le Ce^{-D^2/(45dt)}$ by the Davies--Gaffney bound with $\norm{p_{1/2}(x,\cdot)}_2^2=p_1(x,x)\le C$, and the other terms are $\le Ce^{-cD^2}$; then $\int_1^\infty e^{-a_jt}e^{-c_1D^2/t}dt\le(2C/a_j)e^{-\sqrt{2a_jc_1}D}$. Altogether $K^N_a(x,y)\le C_Y'e^{-2\delta_2'D}$ with $2\delta_2'=\min\{\frac1{48d},\frac c2,\sqrt{\frac32ac_1}\}$, and the block sum follows (at most $3^d$ near blocks; far blocks by the pointwise bound and unit volume).
\end{proof}

\begin{lemma}[linear vanishing at the Dirichlet face]\label{lem:a01:DV}
For every $\Omega$, every $A$ (given \cref{lem:a01:A1}), block data $b$ on the blocks of $\Omega$ and every $x\in\Omega\cap\Xl$: $|u(x)|\le C_{\mathrm{DV}}|b|_\infty\min(1,t(x))$ for the leg $u$; the same for $u=\Mop^{-1}f$ with bounded $f$ ($C_{\mathrm{DV}}\norm f_\infty$).
\end{lemma}
\begin{proof}
Fix the face $\mu$ with $t_\mu(x)=t(x)\le1$ (otherwise \cref{lem:a01:A4}). Let $F_0\defeq a_j\Wt(1+C)|b|_\infty$ bound $|a_jQ^*(b-Qu)|$ ((Q6), \cref{lem:a01:A4} for $|Qu|$) and $S_0$ the sup bound of \cref{lem:a01:A4}. Put $\rho(s)\defeq(S_0+F_0)s-F_0s^2/2$ on $[0,1+\vth]$: increasing there, $\rho(0)=0$, $\rho(1)\ge S_0$, $\rho\ge0$, and $-\Delta_s\rho=F_0$ exactly (second difference of a quadratic). Comparison region $R\defeq$ the first-layer slab $\Omega\cap\Xl\cap\{0<t_\mu\le1\}$, with $\Phi(x)\defeq\rho(t_\mu(x))$. On $R$: $(-\Delta^{N,D}_{0,\Omega}+V_j)\Phi\ge F_0\ge$ the right-hand side of Kato (\cref{lem:a01:A4}(a)), because tangential bonds and tangential Neumann faces see $\Phi$ constant, a Neumann face of $\Omega$ at the far side of the first-layer block drops a neighbour with larger $\Phi$ (which raises $-\Delta^N\Phi$), and on face bonds the killed neighbour carries $\Phi=\rho(0)=0=|u|$. The boundary of $R$ inside $\Omega\cap\Xl$ consists of sites at $t_\mu\in\{1,1+\vth\}$, where $|u|\le S_0\le\Phi$. The operator $-\Delta^{N,D}_{0,R}+V_j$ with these data is a non-singular $M$-matrix ($V_j>0$), so the maximum principle gives $|u|\le\Phi\le(S_0+F_0)t_\mu$ on $R$. Other Dirichlet faces met by $R$ carry $u=0\le\Phi$. For bounded $f$ take $F_0\defeq\norm f_\infty+a_j\Wt C\norm f_\infty$.
\end{proof}

No hypothesis on the faces of $\Omega$ is used: on the cube every sub-domain carries the Dirichlet condition of the cube at the Dirichlet planes, and the barrier sees the Neumann cuts only through the zero drop of a function constant along them.

\begin{lemma}[tangentially localised amplitude]\label{lem:a01:DVloc}
There are $\delta_5>0$ and $C$ such that for every $\Omega$, every $A$, every block $v$ of $\Omega$, every face $\mu$ and every $x$ in the first-layer slab of $\mu$, the leg $l(\cdot,v)$ of the unit datum on $v$ satisfies $|l(x,v)|\le C(\vth+t_\mu(x))e^{-\delta_5\dist(x,B(v))}$; consequently $|u(x)|\le C|b|_\infty(\vth+t(x))e^{-\delta_5\dist(x,\supp b)/2}$ there, and the same for bounded sources supported in a union of blocks.
\end{lemma}
\begin{proof}
The barrier of \cref{lem:a01:DV} is multiplied by tangential weights: for each $\nu\ne\mu$ put $f_\nu(s)\defeq\exp(-\delta|\tilde\psi(s)-\tilde\psi(v_\nu)|)$ with $\tilde\psi(s)\defeq s+(2\pi)^{-1}\sin(2\pi(s+\frac12+\frac\vth2))$, whose derivative vanishes exactly on the block midplanes $s\in\Z-\vth/2$, so that $f_\nu$ takes equal values at the two sites adjacent to any block midplane and every Neumann cut of $\Omega$ (which lies on a block midplane) drops a bond of zero difference. The tangential directions are intervals $[0,N]$ with killed ends, on which $u=0\le$ barrier, so no periodisation is needed. Term by term, $-\Delta_\nu f_\nu\ge-c_\delta f_\nu$ with $c_\delta=(2\pi\delta+4\delta^2)e^{3\delta}$ (the kink at $v_\nu$ is concave). With $\Phi\defeq\rho(t_\mu)\prod_{\nu\ne\mu}f_\nu$ and $F_0,S_0$ enlarged by a fixed factor, $(-\Delta^{N,D}+V_j)\Phi\ge\prod f_\nu[F_0-2(c_\delta+8\delta)\sup\rho]\ge\frac12F_0\prod f_\nu$ for $\delta\le\delta_5$ small; the source $|a_jQ^*(\delta_v-Ql)|\le F_0e^{-\delta_2\dist(\cdot,B(v))}$ (\cref{lem:a01:A4}) is dominated by $\frac12F_0\prod f_\nu$ for $\delta\le\delta_2/\sqrt3$, and so are the data at $t_\mu\in\{1,1+\vth\}$. The maximum principle gives $|l|\le\Phi$, and $\prod f_\nu\le e^{2\delta/\pi}e^{-\delta\sum_\nu|x_\nu-v_\nu|}$ with $\dist(x,B(v))\le\sum_\nu|x_\nu-v_\nu|+(t_\mu(v)-1)_++3$. For bounded sources decompose $f=\sum_vf\mathbf 1_{B(v)}$, apply the barrier block by block and sum, halving the rate.
\end{proof}

\paragraph{Interior estimates.} For regular $A$, $r\in[2\vth,1]$ and a ball $B_{2r}(x_0)$ of sites on which $-\Delta_Au=g$ holds as a full-lattice equation, $\sup_{B_r}|D_Au|\le C(r^{-1}\sup_{B_{2r}}|u|+r\sup_{B_{2r}}|g|)$, together with the $\alpha$-H\"older bound of $D_Au$ on $B_r$; this is \cite[(2.31)--(2.33)]{Ba4}, which gauges away the constant part of $A$ on the ball and perturbs from the free lattice Green's function. On the cube the full-lattice equation holds at every live site whose neighbours are live or killed (extend $u$ by $0$ to $\Kill$: a face bond carries no link), so the estimate applies to every ball containing no killed site and no site within one bond of $\dN\Omega$; the open ball $B_{t(x)}(x)$ contains no killed site.

\begin{lemma}[first differences with decay, up to the face]\label{lem:a01:A5}
For every $\Omega$ and regular $A$, for the leg $u$ of data $b$ supported in a set $S$ of blocks, and for $u=\Mop^{-1}f$ with bounded $f$ supported in $S$: $|D_Au(b')|\le C\norm b_\infty\,e^{-\delta_3\dist(b',S)}$ (resp.\ $C\norm f_\infty$) at every bond $b'$ of $\Omega\cap\Xl$ whose endpoints have depth $\ge\frac14$, face bonds included; on the whole cube at every bond. The $\alpha$-H\"older quotients of $D_Au$ on the balls used in the proof obey the same bounds with the same decay.
\end{lemma}
\begin{proof}
Let $x$ be the base of $b'$ and $D\defeq\dist(x,S)$. For $\vth>\frac1{32}$ every clause follows from $|D_Au(b')|\le2\vth^{-1}\max_{x'\sim x}|u(x')|$ and \cref{lem:a01:A4}; let $\vth\le\frac1{32}$. (1) Face bonds: $|D_Au|=|u(x)|/\vth\le C_{\mathrm{DV}}\norm b$ by \cref{lem:a01:DV}, with decay by \cref{lem:a01:DVloc}. (2) $t(x)\le4\vth$: $|D_Au(b')|\le2\vth^{-1}\max_{x'\sim x}|u(x')|\le Ce^{-\delta_5D/2}$ by \cref{lem:a01:DVloc} (amplitude $\vth+t\le6\vth$). (3) $4\vth<t(x)\le\frac18$: $r\defeq t(x)/2\ge2\vth$; $B_{2r}(x)=B_{t(x)}(x)$ has no killed site and, as $\dep(x)\ge\frac14\ge t(x)+\frac18$, no site within one bond of $\dN\Omega$; the source $g\defeq-V_ju-a_jPu+a_jQ^*b$ (resp.\ $f-V_ju-a_jPu$) obeys $|g|\le Ce^{-\delta_5(D-t)/2}$ on the ball, and $\sup_{B_t}|u|\le Ct\,e^{-\delta_5(D-t)/2}$ by \cref{lem:a01:DVloc} (the sites of $B_t(x)$ have $t_\mu\le2t\le\frac14$); the interior estimate gives $|D_Au(b')|\le C(r^{-1}t+r)e^{-\delta_5D/2}\le Ce^{-\delta_5D/2}$ and $[D_Au]_{\alpha;B_r}\le Ct^{-\alpha}e^{-\delta_5D/2}$. (4) $t(x)>\frac18$: $r\defeq\frac1{16}$, $B_{1/8}(x)$ has no killed site and no site within one bond of $\dN\Omega$ (depth $\ge\frac14\ge\frac18+2\vth$); the interior estimate with \cref{lem:a01:A4} gives $|D_Au|,[D_Au]_\alpha\le Ce^{-\delta_2(D-1/8)}$. Take $\delta_3\defeq\delta_5/4$.
\end{proof}

\begin{corollary}[external legs]\label{lem:a01:legs}
With the amplitude profile $\rho_j(x)\defeq\min\{1,\vth+t(x)\}$ and Ba{\l}aban's external-field norm \cite[(1.32)]{Ba3} (sup with the amplitude, first differences, $\alpha_0$-H\"older quotients of the first differences on unit balls), for base points of depth $\ge2$ on $\Omega$ (every base point on the whole cube): (a) $\norm{a_j\Mop^{-1}Q^*b}\le C\norm b_\infty$ for the leg of block data; (b) $\norm{G_{(j)}(\cdot,z)}\le C$ for the slice lines $G_{(j)}=G_j-G_{j-1}$ on unions of $(j+1)$-blocks, $j\ge1$.
\end{corollary}
\begin{proof}
(a) The sup clause with the amplitude $\rho_j^{1/2}$ from \cref{lem:a01:DV,lem:a01:A4} by the geometric mean $\min(P,Q)\le P^{1/2}Q^{1/2}$; the first-difference clause from \cref{lem:a01:A5}; the H\"older clause: for $|x-y|\ge r(x)$ the quotient is $\le2\sup_{B_1(x)}|D_Au|r^{-\alpha_0}$ and $\rho^{1/2}r^{-\alpha_0}\le C$ in each case of the proof of \cref{lem:a01:A5} ($\rho\le2t$, $r=t/2$, $\alpha_0\le\frac12$; or $r=\frac18$; or $\rho\le5\vth$, $r\asymp\vth$); for $|x-y|<r(x)$ the H\"older bounds of that proof give $\rho^{1/2}t^{-\alpha_0}\le C$. (b) The slice identity $G_{(j)}=a_jG_jQ^*C^{(j)}QG_j$ of \cite[(2.17)]{K1}, which is algebra valid for every weight, with \cref{lem:a01:A23,lem:a01:A4}.
\end{proof}

% a01_H0rel part 3: face-localised differences, the constant datum, the derivative clause for small-field data
\subsection{The face-localised difference and the background of small-field data}\label{app:a01:fld}

\begin{lemma}[reproduction of the constant]\label{lem:a01:RC}
On the doubled torus $2X$ with the periodic extension of the weight (the killed sites of the cube being ordinary zero-weight sites there), and on the cube, $a_jG_jQ^*\mathbf 1=\mathbf 1-V_jG_j\mathbf 1$ at $A=0$, where on the cube $G_j\defeq\Mop(\Xl,0)^{-1}$ and the identity holds up to the face term of \cref{lem:a01:POI}; and $\norm{G_j\mathbf 1}_\infty\le C$.
\end{lemma}
\begin{proof}
$Q\mathbf 1=\mathbf 1$ (Q1), so $P\mathbf 1=Q^*\mathbf 1$. With $v\defeq\mathbf 1-V_jG_j\mathbf 1$: $(-\Delta+V_j+a_jP)v=V_j+a_jP\mathbf 1-V_j=a_jQ^*\mathbf 1$ on the torus, where $-\Delta\mathbf 1=0$. On the cube $-\Delta^D\mathbf 1=\mathsf B$, the face term of \cref{lem:a01:POI}. The sup bound is \cref{lem:a01:A4}(c).
\end{proof}

\begin{lemma}[the Dirichlet Poisson term decays]\label{lem:a01:POI}
Let $\mathsf B(x)\defeq\vth^{-2}\#\{\text{killed neighbours of }x\}$ (supported on $F_1$) and $v_0\defeq\Mop(\Xl,0)^{-1}\mathsf B$. Then $|v_0(x)|\le Ce^{-(\delta_2/2)(t(x)-4)}$ for every $x$, and $|D_0v_0(b')|\le Ce^{-(\delta_2/2)(t(b')-5)}$ at every bond with $t\ge6$.
\end{lemma}
\begin{proof}
(1) $\Mop(\Xl,0)\mathbf 1=\mathsf B+V_j+a_jQ^*\mathbf 1$ on $\Xl$, so $h^0\defeq a_j\Mop^{-1}Q^*\mathbf 1=\mathbf 1-\Mop^{-1}V_j-v_0$ and $v_0$ is bounded by $1+Cc_V+C$ (\cref{lem:a01:A4}). (2) Let $\chi(x)\defeq\prod_\mu\psi(x_\mu)\psi(N-x_\mu)$ with $\psi$ smooth, $0$ on $(-\infty,3.25]$, $1$ on $[3.75,\infty)$, constant on a $\vth$-neighbourhood of every block plane. Then $\Mop(\chi v_0)=[\Mop,\chi]v_0$ (since $\mathsf B=0$ where $\chi\ne0$), supported on the blocks at face distance in $[3,4]$; the Caccioppoli inequality on a two-block neighbourhood (where $\Mop v_0=0$; the bond sum over $\Xl$ has no boundary term, the face bonds entering with $v_0=0$ at $\Kill$) bounds $\norm{\mathbf 1_B[\Mop,\chi]v_0}_2\le C\sup|v_0|\le C$, and \cref{lem:a01:A4}(c) with block-$L^2$ data gives $|\chi v_0(x)|\le C\sum_{B\in\text{shell}}e^{-\delta_2\dist(x,B)}\le Ce^{-(\delta_2/2)(t(x)-4)}$. (3) The interior estimate on $B_{1/4}(x)$ at $t\ge6$ with $g=-V_jv_0-a_jPv_0$.
\end{proof}

\begin{lemma}[face-localised difference, whole cube, regular $A$]\label{lem:a01:FLD}
Let $T\defeq2X$ carry the periodic extension of the weight and the even extension $A_T\defeq\Ebar A$ of a regular field $A$ (regular with the same constant by \cref{lem:a01:REG}(b)), and $\Mop_T\defeq-\Delta_{A_T}+V_j+a_jP^T$. For bounded $F$ supported in $\Xl$ put $v\defeq\Mop_T^{-1}F$, $u\defeq\Mop(\Xl,A)^{-1}F$ and $w\defeq v|_{\Xl}-u$. Then for every site $x$ and fine bond $b$ of the cube,
\[
|w(x)|+|D_Aw(b)|\ \le\ C\exp\big[-\delta_F\big(\dist(\cdot,\DD)+\dist(\supp F,\DD)+\dist(\cdot,\supp F)\big)\big]\,\norm F_\infty,
\]
with $C,\delta_F$ depending on $(d,L,a,\Wt,c_V)$ and King's constants only.
\end{lemma}
\begin{proof}
(1) For $x\in\Xl$, $(\Mop_Tv)(x)$ and $(\Mop(\Xl,A)v|_X)(x)$ differ only by the hopping to killed neighbours, since $P^T(x,\cdot)$ is supported in the block of $x$, whose weighted sites are live (Q8). So $(\Mop(\Xl,A)v|_X)(x)=F(x)+h(x)$ with $h(x)\defeq\vth^{-2}\sum_{x'\in\Kill,x'\sim x}U_{xx'}v(x')$ supported on $F_1$, and $w=\Mop^{-1}h$. (2) Near the face, $\dist(x,\DD)\le D_0\defeq4+2\sqrt3$: $|w|\le|v|+|u|$ and $|D_Aw|\le|D_Av|+|D_Au|$; on the torus $|v(x)|\le Ce^{-\delta\dist(x,\supp F)}\norm F$ (\cref{lem:a01:A4}(c) on $T$) and $|D_Av(b)|\le Ce^{-\delta\dist(b,\supp F)/2}\norm F$ (interior estimate on unit balls); on the cube the same for $u$ up to the face (\cref{lem:a01:A4}(c), \cref{lem:a01:A5} in the bounded-source form). As $\dist(\supp F,\DD)\le\dist(x,\supp F)+D_0$ near the face, the claim holds there with $\delta_F\le\delta/4$. (3) Away from the face: let $\chi$ be smooth, $\chi=0$ on $\{\dist(\cdot,\DD)\le1+\sqrt3\}$, $\chi=1$ on $\{\dist\ge D_0-1\}$, $|\nabla\chi|\le1$, $|\nabla^2\chi|\le C$. Then $\chi h=0$ and $\Mop(\chi w)=[\Mop,\chi]w=:g$ with $g=[-\Delta_A,\chi]w+a_j[P,\chi]w$; the lattice Leibniz rule gives $|[-\Delta_A,\chi]w|\le C(|w|+|D_Aw|)$ on the shell $S\defeq\{\sqrt3\le\dist(\cdot,\DD)\le D_0\}$, and $|[P,\chi]w(x)|\le\sum_{x'}\vth^3|P(x,x')||\chi(x)-\chi(x')||w(x')|\le2\sqrt3\,\Wt\sup_{B(x)}|w|$ by (Q6), (Q8), supported on the blocks meeting the shell. So $g$ is bounded, supported in $S'$ (the union of $S$ with the blocks meeting it, at distance $\le D_0+\sqrt3$ from $\DD$ and $\ge1$ from $\Kill$), and by (2) $|g(z)|\le Ce^{-\delta\dist(z,\supp F)/2}\norm F$. (4) By \cref{lem:a01:A4}(c) in the decaying form, $|\chi w(x)|\le C\sup_{z\in S'}e^{-\delta_2|x-z|/2}|g(z)|$, hence for $\dist(x,\DD)\ge D_0$: $|w(x)|\le C\sup_{z\in S'}e^{-\delta_2|x-z|/2-\delta\dist(z,\supp F)/2}\norm F$; with $|x-z|\ge\dist(x,\DD)-D_0-\sqrt3$, $\dist(z,\supp F)\ge\dist(\supp F,\DD)-D_0-\sqrt3$ and $|x-z|+\dist(z,\supp F)\ge\dist(x,\supp F)$, averaging the exponents gives the claim for $|w|$ with $\delta_F\defeq\min(\delta,\delta_2)/6$. (5) For $\dist(x,\DD)\ge D_0$, $w$ solves $-\Delta_Aw=-V_jw-a_jPw$ on $B_{1/4}(x)$ ($h=0$ there), so the interior estimate gives $|D_Aw|\le C\Wt\sup_{B_{1/2}(x)+\text{one block}}|w|$, which obeys the bound of (4) with $D_0$ replaced by $D_0+1$.
\end{proof}

The comparison operator is the torus operator with an arbitrary regular extension of $A$; the scalar is never reflected, so no regularity of a reflected field is needed. No flat-face condition is used: the cut-off sees only $\dist(\cdot,\DD)$, so edges and corners of the cube are included.

\begin{lemma}[the Dirichlet image effect on a sub-domain, sup form]\label{lem:a01:LOC}
Let $\Omega$ be a sub-domain and $\Omega^\sharp$ a Neumann domain without Dirichlet face that contains $\Omega\cap\Xl$ and the killed sites adjacent to it, adds no bond from $\Omega\cap\Xl$ except the bonds to those sites, and carries a regular extension $A^\sharp$ of $A$ (for instance $\Omega$ together with its mirror blocks across the Dirichlet planes it meets, with the even extension of every component of $A$). For data $b$ on the blocks of $\Omega$ (zero on the added blocks) let $u\defeq a_j\Mop(\Omega,A)^{-1}Q^*b$ and $u^\sharp\defeq a_j\Mop(\Omega^\sharp,A^\sharp)^{-1}Q^*b$. Then for every $x\in\Omega\cap\Xl$, $|u(x)-u^\sharp(x)|\le C_U|b|_\infty e^{-(\delta_2/2)(t(x)-4)}$ with $C_U=C_U(d,a,\Wt,\gamma)$, free of Ba{\l}aban's large-block size $M$, of $\vth$, $j$, $\Omega$ and $A$; and $|D_A(u-u^\sharp)(b')|\le C|b|_\infty e^{-(\delta_2/2)(t(b')-5)}$ at bonds with $t\ge6$ and depth $\ge1$.
\end{lemma}
\begin{proof}
$w\defeq u-u^\sharp|_\Omega$ solves $\Mop(\Omega)w=0$ at every site of $\Omega\cap\Xl$ off $F_1$ (same bonds, same blocks; the killed sites have zero weight). The product cut-off of \cref{lem:a01:POI}, the Caccioppoli inequality in the region $t\ge2.25$, $\sup|w|\le|u|+|u^\sharp|\le C|b|$ (\cref{lem:a01:A4} on both domains) and \cref{lem:a01:A4}(c) with block-$L^2$ data give the sup bound exactly as in \cref{lem:a01:POI}; the derivative bound is the interior estimate.
\end{proof}

\begin{lemma}[transported constant up to the face]\label{lem:a01:H}
Let $A$ be regular on $\Xl$, $b$ a fine bond with base $x_0$, $\lambda_0(z)\defeq\sum_\mu A_\mu(x_0)z_\mu$, $A'\defeq A-d\lambda_0$ (gauge equivalent to $A$ on the simply connected cube, with $A'=0$ on the bonds at $x_0$ and $|A'_{b''}|\le K_j|b''-x_0|_1$), and $h\defeq a_j\Mop(\Xl,A')^{-1}Q(A')^*\mathbf 1$. Then $|D_{A'}h(b)|\le C_H(V_j+\theta_j+e^{-\delta_ht(x_0)})$ with $\delta_h\defeq\delta_2/2$ and $C_H=C_H(d,L,a,\Wt,\gamma,c_V)$.
\end{lemma}
\begin{proof}
If $t(x_0)<6$ the claim is \cref{lem:a01:A5} ($h$ is the leg of the datum $\mathbf 1$). Let $t(x_0)\ge6$. (1) At $A=0$: by \cref{lem:a01:POI}(1), $D_0h^0=-D_0\Mop_0^{-1}V_j-D_0v_0$, so $|D_0h^0(b)|\le CV_j+Ce^{-\delta_h(t-5)}$ (\cref{lem:a01:A5} for the bounded source $V_j$, \cref{lem:a01:POI}). (2) $w\defeq h-h^0=\Mop(A')^{-1}S$ with $S\defeq a_j(Q(A')^*-Q(0)^*)\mathbf 1-(\Mop(A')-\Mop(0))h^0$. Pointwise $|S(z)|\le C\theta_j(1+|z-x_0|)^2$: the block terms by $|U(A'(\Gamma))-1|\le e_j|A'||\Gamma|\le C\theta_j(1+|z-x_0|)$ and $|h^0|\le C$; the Laplacian terms by pairing the two bonds of each direction, $\vth^{-2}[(1-U_+)+(1-U_-)]h^0(z)+\vth^{-2}\sum_\pm(1-U_\pm)(h^0(z_\pm)-h^0(z))$, where the first bracket is $-\vth e_jq(a_+-a_-)+O(\vth^2e_j^2|A'|^2)$, so its $\vth^{-2}$-multiple is $\le\theta_j+C\theta_j^2|z-x_0|^2$, and the second is $\le\sum e_j|A'||D_0h^0|\le C\theta_j|z-x_0|$ (\cref{lem:a01:A5} at $A=0$); at a first-layer site the normal pair is incomplete and the unpaired term is $\le\vth^{-1}e_j|A'|\cdot2C_{\mathrm{DV}}\vth\le C\theta_j|z-x_0|$ by \cref{lem:a01:DV}. (3) $|w(x)|\le C\theta_j(1+|x-x_0|)^2$ by \cref{lem:a01:A4}(c), shell by shell around $x_0$. (4) The interior estimate for $w$ on $B_{1/4}(x_0)$ (no killed site, $t\ge6$) with $-\Delta_{A'}w=S-V_jw-a_jP(A')w$, bounded by $C\theta_j$ on $B_{1/2}(x_0)$, gives $|D_{A'}w(b)|\le C\theta_j$. (5) $D_{A'}h(b)=D_0h^0(b)+(D_{A'}-D_0)h^0(b)+D_{A'}w(b)$, and the middle term vanishes since $A'=0$ on the bonds at $x_0$.
\end{proof}

\begin{proposition}[derivative clause for backgrounds of small-field data]\label{prop:a01:legsd}
Let $A$ be regular on $\Xl$, $\phi_j$ block data, $b$ a fine bond with base $x_0$ in the block $y_0$, $r_j\defeq R(1+\ln\ell_j^{-1})^{\varrho}$ Ba{\l}aban's region radius, and assume: (G) $|D_{A_j}\phi_j(\beta)|\le p_j$ on every block bond within block distance $2r_j$ of $y_0$; (S) $|\phi_j|\le M_j$; (F) King's face test $|\phi_j(y)|\le\sigma_j\le C_8p_j$ on the face blocks within $2r_j$ of $y_0$; and $M_j(V_j+\theta_j)\le p_j$, $M_je^{-\delta_*r_j}\le p_j$ with $\delta_*\defeq\min(\delta_3,\delta_h)/2$ (both hold for $\eta\le\eta_1$ and $R\ge R_*$: $M_j\theta_j\le Ce\lambda^{-1/4}L^{1/4}(1+\ln L)^pp_j^2\ell_j^{1/4}$). Then the background $\phi^{(j)}=a_j\Mop(\Xl,A)^{-1}Q(A)^*\phi_j$ satisfies $|D_A\phi^{(j)}(b)|\le C_{\mathrm{legs}}p_j$ at every fine bond of the cube, face bonds included, $C_{\mathrm{legs}}=C(1+C_8)$. On a sub-domain $\Omega$, under the same hypotheses for the data on $\Xl$, the domain background $\phi^{(j)}_\Omega\defeq a_j\Mop(\Omega,A)^{-1}Q^*(\phi_j|_\Omega)$ satisfies, at every bond of depth $\ge\frac12$,
\[
|D_A\phi^{(j)}_\Omega(b)|\ \le\ C_{\mathrm{legs}}p_j+CM_je^{-\delta'(\dep(b)-4)_+},\qquad\delta'\defeq\delta_3/2,
\]
hence $\le(C_{\mathrm{legs}}+1)p_j$ at depth $\ge\rho_j\defeq4+\delta'^{-1}\ln(CM_j/p_j)$, and $\rho_j\le r_j$ once $R\ge R'_*$.
\end{proposition}
\begin{proof}
Gauge to $A'$ as in \cref{lem:a01:H} (every quantity is gauge covariant). (1) For blocks $y$ with $|y-y_0|\le r_j$ let $U_y$ be the coarse transporter along the coordinate-ordered block path from $y_0$ to $y$ (inside the convex cube); telescoping (G) gives $\phi_j(y)=U_yc+\psi(y)$, $c\defeq\phi_j(y_0)$, $|\psi(y)|\le p_j|y-y_0|_1$; for $|y-y_0|>r_j$ put $\psi(y)\defeq\phi_j(y)$, $|\psi|\le M_j$. In the gauge $A'$, $|U_y-1|\le C\theta_j(1+|y-y_0|)^2$. (2) The $\psi$-part: $|D_{A'}[a_j\Mop^{-1}Q^*\psi](b)|\le\sum_y|D_{A'}l(b,y)||\psi(y)|\le C(p_j+M_je^{-\delta_3r_j})\le Cp_j$ by the decaying clause of \cref{lem:a01:A5}; the transport correction $|c||D_{A'}[a_j\Mop^{-1}Q^*((U-1)\mathbf 1_{|y-y_0|\le r_j})](b)|\le CM_j\theta_j\le Cp_j$; the truncation $|c||D_{A'}[a_j\Mop^{-1}Q^*\mathbf 1_{|y-y_0|>r_j}](b)|\le CM_je^{-\delta_3r_j}\le Cp_j$. (3) The constant part: $|c||D_{A'}h(b)|\le|c|C_H(V_j+\theta_j+e^{-\delta_ht(x_0)})$ by \cref{lem:a01:H}. (4) The cap: if $t(x_0)\le r_j$, transport along the normal block path from $y_0$ to the face layer (length $\le t(x_0)+1$) and (F) give $|c|\le\sigma_j+p_j(t(x_0)+1)$, so $|c|e^{-\delta_ht}\le C(1+C_8)p_j$; if $t>r_j$, $|c|e^{-\delta_ht}\le M_je^{-\delta_hr_j}\le p_j$; and $|c|C_H(V_j+\theta_j)\le Cp_j$. On $\Omega$: $w\defeq\phi^{(j)}_\Omega-\phi^{(j)}|_\Omega$ solves $\Mop(\Omega)w=0$ at depth $\ge1$ (same bonds and blocks there); with $\chi\defeq\psi(\tilde d_N)$, $\tilde d_N$ a mollified depth constant near block planes, $\Mop(\Omega)(\chi w)=[\Mop,\chi]w=:F$, supported at depth in $[3.25,3.75]$ plus one block, with $|F|\le C(|w|+|D_Aw|+\sup_{\text{block}}|w|)\le CM_j$ (\cref{lem:a01:A4} for both backgrounds, \cref{lem:a01:A5} at depth $\ge3$); \cref{lem:a01:A5} for bounded sources gives $|D_A(\chi w)(b)|\le CM_je^{-\delta_3(\dep(b)-4)}$, and $\chi=1$ near $b$ when $\dep(b)\ge4$; for $\frac12\le\dep(b)<4$ use \cref{lem:a01:A5} on both backgrounds.
\end{proof}

The face test is necessary: for the constant data $\phi_j\equiv c$ the background vanishes at the Dirichlet face and $|D_A\phi^{(j)}|$ is of order $|c|$ on the first block; this is why King imposes \cite[(5.41)]{K2} at each step.

\begin{lemma}[face-localised background of a covariantly constant datum, regular $A$]\label{lem:a01:FL}
Let $\Omega$ be a box of large blocks with Dirichlet conditions on its faces (the cube, or a sub-box), $y_0$ a block, $c\in\C$, $\mathsf T$ a tree of coordinate-ordered block paths from $y_0$, $\Theta_A(y)\defeq\overline{U(A(\Gamma_{\mathsf T}(y_0\to y)))}\,c$ the datum covariantly constant along $\mathsf T$, and $u_A\defeq a_j\Mop(\Omega,A)^{-1}Q(A)^*\Theta_A$. There are $\delta_1>0$, $C$ and $\theta_*>0$ such that for $A$ regular on $\Omega$ with $\theta_j\le\theta_*$, every bond $b$ of $\bar\Omega$ (face bonds included) and every block $y$: (a) $|D_{A,b}u_A|\le|c|[CV_j+Ce^{-\delta_1\dist(b,\DD)}+C\theta_j(1+|b-y_0|)^3]$; (b) $|\Theta_A(y)-(Q(A)u_A)(y)|\le|c|[CV_j+Ce^{-\delta_1\dist(B(y),\DD)}+C\theta_j(1+|y-y_0|)^3]$; (c) $\sup|u_A|\le C|c|$.
\end{lemma}
\begin{proof}
Step 1 ($A=0$, $\Theta=c$). The Dirichlet operator of a box of large blocks with King's weight is the odd restriction of the free operator on $\vth\Z^3$ with the periodic weight: the odd extension of a function vanishes on the Dirichlet planes, the weight couples no live site across a plane, and the free operator preserves the odd subspace, on which it is unitarily $\Mop(\Omega,0)$ (coercivity gives uniqueness). So $\Mop(\Omega,0)^{-1}F=(G^{\mathrm{free}}\Ebar_{\mathrm{odd}}F)|_{\bar\Omega}$ and $\Ebar_{\mathrm{odd}}Q^*c=Q^*(cs)$ with $s(y)=\pm1$ the coarse odd extension. Write $s=1-(1-s)$: by \cref{lem:a01:RC} on the free lattice, $a_jG^{\mathrm{free}}Q^*\mathbf 1=\mathbf 1-V_jG^{\mathrm{free}}\mathbf 1$, so $u_0=c[\mathbf 1-V_jG^{\mathrm{free}}\mathbf 1-a_jG^{\mathrm{free}}Q^*(1-s)]$ on $\bar\Omega$; $(1-s)$ is supported on the image boxes with $s=-1$, at distance $\ge\dist(b,\DD)$ from $b$, and bounded by $2$. With the weighted-$\ell^1$ bounds $\norm{DG^{\mathrm{free}}}_{\delta_1},\norm{G^{\mathrm{free}}}_{\delta_1}\le C$ of the free weighted operator (\cref{lem:a01:A23,lem:a01:A4} on $\vth\Z^3$), $|Du_0(b)|\le|c|[V_j\norm{DG^{\mathrm{free}}}_0+2a_j\Wt\norm{DG^{\mathrm{free}}}_{\delta_1}e^{-\delta_1\dist(b,\DD)}]$ at every bond, face bonds included (the intertwining holds there); the same computation with $Q$ in place of $D$ gives (b), and (c) is $\norm{G^{\mathrm{free}}}_0$. Step 2 (constant $B$): gauge covariance, $u_B=U_\lambda u_0(c')$ with $|c'|=|c|$, $\lambda=B\cdot x$; (a)--(c) hold with the $\theta$ term absent. Step 3 (regular $A$). Fix the bond $b$ and let $B$ be the constant field with the components of $A$ at reference bonds through $b$; $A'\defeq A-B$ obeys $e_j|A'(x)|\le d\theta_j(1+|x-b|_1)$ on the box (sum the first differences along a lattice path). With $w\defeq u_A-u_B$ (both vanish on $\DD$): $\Mop_Aw=F\defeq a_j[Q(A)^*\Theta_A-Q(B)^*\Theta_B]+(\Mop_B-\Mop_A)u_B$; the hopping difference is a first-order operator with coefficients $\le C\theta_j(1+|x-b|)$ and the block part has row sums $\le C\theta_j(1+|x-b|)$, the contour transports differ by $\le e_j\sum_\Gamma\vth|A'|$, and $|u_B|,|D_Bu_B|\le C|c|$ by Steps 1--2; so $|F(x)|\le C\theta_j|c|(1+|x-b|+|y_0-b|)^3$. Split $F=\sum_nF_n$ by shells $n\le|x-b|<n+1$; \cref{lem:a01:A5} in the bounded-source form (face bonds included) gives $|D_{A,b}\Mop_A^{-1}F_n|\le Ce^{-\delta'(n-1)}\norm{F_n}_\infty$ and the same for $|\Mop_A^{-1}F_n|$ within one block of $b$; summing, $|D_{A,b}w|+\sup_{|x-b|\le2}|w|\le C\theta_j|c|(1+|y_0-b|)^3$. Finally $D_{A,b}u_B=D_{B,b}u_B$ since $A$ and $B$ agree on $b$ (on a face bond both are the gauge-free $\vth^{-1}|u_B|$), which gives (a); (b) from the four-term split $\Theta_A-Q(A)u_A=[\Theta_B-Q(B)u_B]+[\Theta_A-\Theta_B]-[(Q(A)-Q(B))u_B]-Q(A)w$, and (c) from the sup clause with $\norm F_\infty=a_j\Wt|c|$.
\end{proof}

The face rate $\delta_1$ is the free rate of the weighted operator at $A=0$ and does not depend on Ba{\l}aban's large-block size $M$; this is what allows the near-layer constant $R_0$ of \cref{app:a01:upper} to be fixed before $M$.

\begin{lemma}[the block-mismatch identities]\label{lem:a01:EF}
\label{lem:a01:EB}
For every weight, domain, boundary condition and gauge field, with $\phi^{(j)}=a_j\Mop^{-1}Q(A)^*\phi_j$: (a) $\phi_j-Q(A)\phi^{(j)}=a_j^{-1}\Delta^{(j)}(A)\phi_j$; (b) with the \emph{uniform} transported average $Q^u(A)\psi(y)\defeq\sum_{x\in B(y)}\vth^3U(A(\Gamma_{y,x}))\psi(x)$, exactly $a_j(\phi_j-Q(A)\phi^{(j)})(y)=Q^u(A)[(-\Delta_A+V_j)\phi^{(j)}](y)$, whereas with $Q(A)$ in place of $Q^u(A)$ the left side acquires the factor $W_2$.
\end{lemma}
\begin{proof}
(a) $Q\phi^{(j)}=a_jQ\Mop^{-1}Q^*\phi_j$ and $\Delta^{(j)}=a_j-a_j^2Q\Mop^{-1}Q^*$. (b) The minimiser equation is $(-\Delta_A+V_j)\phi^{(j)}=a_jQ^*(\phi_j-Q\phi^{(j)})=:a_jQ^*m$, and pointwise $(Q^*m)(x)=w(x)U(\Gamma_{y,x})^*m(y)$ for $x\in B(y)$; applying $Q^u(A)$ gives $a_jm(y)\sum_x\vth^3w(x)=a_jm(y)$ by (W1) and $UU^*=1$; applying $Q(A)$ gives $a_jW_2m(y)$.
\end{proof}

Identity (b) is the form in which the block-mismatch clause of the lower bound is read on the cube; the naive summation by parts of the \emph{weighted} average $Q[-\Delta\phi](y)$ produces, besides the boundary flux of the block, interior terms $\sum_{\text{bonds}}\vth^{d-2}(w(x)-w(x'))(\phi(x)-\phi(x'))$ across the zero-weight layers, whose absolute bound $2d\Wt\sup|D\phi|/(L\vth)$ diverges as $\vth\to0$; the uniform transported average has no interior terms, and (a) needs no fine-level summation by parts at all.

% a01_H0rel part 4: Dirichlet-face positivity, the face family, the positivity (3.26) at the weight
\subsection{The Dirichlet-face positivity and the face family}\label{app:a01:face}

\begin{lemma}[Dirichlet-face positivity]\label{lem:a01:FPC}
In unit-$j$ units let $y'\in\Lone$ be a face block, $B(y')$ its unit cube, and $\Delta^{(j)}$ the block form of the cube at any gauge field $A$. Then for every block field $\phi$ and $\theta\in(0,1)$,
\[
\ip\phi{\Delta^{(j)}\phi}\ \ge\ (1-\theta)\ip\phi{\Delta^{(j)}\phi}+\theta\,\gamma_D\,|\phi(y')|^2,\qquad\gamma_D=\frac{a_j}{1+a_j\Wt^2}.
\]
The same holds for the composite multiscale form of \cite[(3.23)--(3.25)]{Ba2} restricted to the fine sites of $B^j(y')$, in the form
\[
\sum_{\text{bonds of }\overline{B}}\eps^{d}|D_A\phi_0|^2+\sum_{i<j}a\,\ell_{i+1}^{d-2}\sum_{B^{i+1}(y'')\subset B}|\phi_{i+1}(y'')-Q(A)\phi_i(y'')|^2\ \ge\ (\Wt+a_j^{-1})^{-1}\ell_j^{d-2}|\phi_j(y')|^2
\]
(physical units; the first sum over the bonds inside $B=B^j(y')$ and the face bonds from $B$ to the killed layer), for every gauge field.
\end{lemma}
\begin{proof}
By the minimiser identity, $\ip\phi{\Delta\phi}=\min_\Phi F(\Phi)$, $F=\ip\Phi{(-\Delta^D_A+V_j)\Phi}+a_j\norm{\phi-Q\Phi}^2$. In $\theta F$ keep only the gradient terms of bonds inside $B(y')$, the Dirichlet terms of the face bonds of $B(y')$ (each $\vth^{d-2}|\Phi(x)|^2$) and the block term of $y'$; all dropped terms are $\ge0$. With $g\defeq|\Phi|$, unitary transport gives $|D_A\Phi(b)|\ge|g(b_+)-g(b_-)|$, and the face terms are $\vth^{d-2}g(x)^2$. Discrete Friedrichs inequality on the unit cube with the killed layer behind one face: on each line normal to that face, $g(x)=\sum_{s\le m}[g(s)-g(s-1)]$ with $g(0)\defeq0$ at the killed site and $m\le\vth^{-1}$, so $g(x)^2\le\vth^{-1}\sum_{\text{line}}|\partial g|^2$; summing $\vth^dg^2$ over the line and over all lines, $\sum_{B(y')}\vth^dg^2\le E(g)$, the kept energy. Then $|(Q\Phi)(y')|\le\sum\vth^dw|\Phi|\le\Wt(\sum\vth^dg^2)^{1/2}\le\Wt E^{1/2}$. So $\theta F\ge\theta\min_{s\ge0}[\Wt^{-2}s^2+a_j(|\phi(y')|-s)_+^2]=\theta\gamma_D|\phi(y')|^2$ (springs in series), and minimising over $\Phi$ gives the first claim. Composite form: with $s_i\defeq(\ell_i^d\sum_B|\phi_i|^2)^{1/2}$, Friedrichs along the normal lines from the killed layer gives $s_0^2\le\ell_j^2E$ with $E$ the first sum; Cauchy--Schwarz with the weights gives $\norm{Q_1\phi_0}\le\Wt^{1/2}s_0$ at the weighted step and $\norm{Q\phi_i}\le s_i$ at the uniform steps $i\ge1$; the reverse triangle inequality gives $s_j\le\Wt^{1/2}s_0+\sum_it_i$, $t_i\defeq\norm{\phi_{i+1}-Q\phi_i}$, and Cauchy--Schwarz with the weights $(\Wt\ell_j^2,\ell_{i+1}^2/a)$ gives $s_j^2\le(\Wt\ell_j^2+\ell_j^2/a_j)(E+\sum_ia\ell_{i+1}^{-2}t_i^2)$, using $\sum_{i<j}\ell_{i+1}^2/a=\ell_j^2/a_j$.
\end{proof}

This is the extra term $\gamma_0\sum_{x\in\bar\Omega^{(k)}}|\phi(x)|^2$ of \cite[Prop.~5.4, (5.26)]{K2}, valid for every gauge field (King's printed proof \cite[(5.28)--(5.29)]{K2} expands around a constant field and needs $A$ regular). King's constant \cite[(5.38)--(5.40)]{K2} is consistent with it.

\begin{definition}[the face family]\label{def:a01:facefam}
At every scale $0\le i\le K_f$ of the upper bound's large-field decomposition, $i=0$ included, the \emph{face family} is the set of events $\{|\phi_i(y')|>\sigma_i\}$ at $y'\in\Lone^{(i)}\cap\Lambda_7^{(i-1)\prime}$ (with $\Lambda_7^{(-1)\prime}\defeq$ everything), where $\Lone^{(i)}$ is the set of $i$-blocks containing a site of $F_1$ (for $i\ge1$ the face blocks of both faces in every direction; for $i=0$ the first live layer $F_1$ itself, where $\phi_0$ is the fine field), and $\sigma_i\defeq(2/(\theta_2\gamma_D^{(i)}))^{1/2}p_i\le C_8p_i$ with $\gamma_D^{(0)}\defeq1$. At the stopping scale the same test is the indicator $\lambda_j$ of \cref{app:a01:weights} (it is \cite[(5.41)]{K2} with $C=\sigma_j/p_j$).
\end{definition}

At $i=0$ the payment is the $\theta_2$-share of the face-bond term itself: $\ip{\phi_0}{(-\Delta^D_A+V)\phi_0}\ge\sum_{\text{face bonds }b\ni x}|D_A\phi_0(b)|^2=n_F(x)|\phi_0(x)|^2$ (unit-$0$ units; no link on a face bond), with $n_F(x)\in\{1,2,3\}$ the number of face bonds at $x$. The scale-$0$ member must be stated on the first live layer (not as the set of boundary blocks, which at $i=0$ would be the killed plane on one face), because the inherited cap of \cref{lem:a01:IC} at $k=1$ needs a cap on the first live layer of both faces.

\begin{lemma}[payment, and the corrected near-face estimates]\label{lem:a01:F541}
(a) A violation of the face test at $y'$ at scale $i$ is paid by the factor $e^{-p_i^2/4}$ from the $\theta_2$-share of the scale-$i$ form (\cref{lem:a01:FPC}), and in the composite form of \cite[Prop.~3.1]{Ba2} by its composite clause; violating blocks of different scales are disjoint as fine sets by Ba{\l}aban's region rules, so the payments add. (b) For $x\in\Lambda_4^{(i)}$ with $\dist(x,\DD)\le r(\ell_i)$, a coordinate block path from $x$ to the face layer of length $\le r(\ell_i)$ lies in $\Lambda_3^{(i)}\subset\Lambda_0^{(i)}$ (since $\dist(\Lambda_4,\Lambda_3^c)>r$), so on the complement of the face family and of the gradient family $|\phi_i(x)|\le\sigma_i+r(\ell_i)p_{i-1}\le Cp_i^2$, and the fluctuation field obeys $|\phi_i'(x)|\le Cp_i^2$ there. (c) King's printed threshold ``$\dist(x,\partial T^{(k)})<p(L^k\eps)$'' in \cite[(5.43)--(5.45)]{K2} cannot be used: $p(\ell)=b_0(1+\ln\ell^{-1})^p$ with $2p\ge(d+1)\varrho$ (\cite[(3.54)]{Ba2}), so $p(\ell)/r(\ell)\to\infty$ and a contour of length up to $p(\ell)$ leaves $\Lambda_0^{(i)}$, where \cite[(5.41)]{K2} is imposed; the correct threshold is one corridor $r(\ell_i)$, with which the conclusions of (b) hold. (d) For $x$ at distance $>r(\ell_i)$ from $\DD$ the Dirichlet image effect on the background is $\le Ce^{-\delta_1r(\ell_i)}M_i\le p_i$ (\cref{lem:a01:FL,lem:a01:LOC}; $\varrho>1$), so \cite[(2.113)]{Ba2} holds there as printed. (e) Within distance $R_0$ of the face either a gradient violation or a face-family element is present on a shortest block path from a block of $R_s^{(i)}$ to the face layer, and at face blocks off the families the kept quartic is $\ge-C\lambda(\ell_i)p_i^8$ per block, an error $\le C\ell_i^\alpha$ per block; this is \cite[(5.46)--(5.49)]{K2} in event form.
\end{lemma}
\begin{proof}
(a) \Cref{lem:a01:FPC} with $\theta=\theta_2$ and $\sigma_i$ chosen so that $\frac12\theta_2\gamma_D\sigma_i^2/4=p_i^2/4$; the disjointness is the nesting of Ba{\l}aban's regions applied to a pointwise family (the face family is a condition on one block variable at one block of the index set, of the same type as $|\psi(y)-(Q\phi)(y)|>p$, so the rules (2.4)--(2.10) of \cite{Ba2} hold with it verbatim). (b) Triangle inequality along the path, the gradient clause $|D_{A_i}\phi_i|\le p_{i-1}$ on each of $\le r$ bonds and $|U|=1$; the fluctuation $\phi'=\phi_i-aL^{-2}C^{(i)}Q^*\phi_{i+1}$ has its second term bounded by $C\sup|\phi_{i+1}|$ on a neighbourhood (decay of $C^{(i)}$, \cref{lem:a01:A23}), which by (b) at scale $i+1$ is $\le Cp^2$. (c) Arithmetic: $p(\ell)/r(\ell)=(b_0/R)(1+\ln\ell^{-1})^{p-\varrho}$ with $p-\varrho=4$ at Ba{\l}aban's values $p=6$, $\varrho=2$. (d) \Cref{lem:a01:LOC} (sup form) times the block-sup bound $M_i$; $e^{-\delta R(1+\ln\ell^{-1})^2}M_i\le p(\ell)$ for $\ell\le\eta_1$. (e) Along a shortest cube-internal block path from $x\in R_s^{(i)}$ (where $|\phi_i(x)|\ge M_i$) to $\Lone^{(i)}$, of length $\le dR_0<s_b$ inside $\Lambda_7^{(i-1)\prime}$, either a bond carries a gradient violation, or the end block $y'$ has $|\phi_i(y')|\ge M_i-dR_0p_i\ge\frac12M_i\ge\sigma_i$ (for $\lambda(\ell_i)\le(2C_8)^{-4}$ and $dR_0\le\frac12\lambda(\ell_i)^{-1/4}$, conditions on $\eta_1$), a face-family element; the per-block error is $\lambda(\ell_i)(Cp_i^2)^4$ by (b).
\end{proof}

\begin{lemma}[Ba{\l}aban's per-scale lower bound at the weight]\label{lem:a01:329}
For a union $\Lambda$ of $k$-blocks with Neumann cuts and, where $\Lambda$ meets $\DD$, the Dirichlet condition of the cube, and $A$ regular on $\Lambda$ with $\theta_k\le\theta_*$,
\[
\ip\phi{\Delta^{(k)}(\Lambda,A)\phi}\ \ge\ \gamma_0^{w}\Big[\sum_{\text{block bonds }\langle y,y'\rangle\subset\Lambda}|U(A_k(\langle y,y'\rangle))\phi(y')-\phi(y)|^2+V_k\sum_y|\phi(y)|^2\Big]-C(\theta_k+e_kp_k)^2\sum_{y\in\Lambda}|\phi(y)|^2,
\]
with $\gamma_0^w>0$ depending on $(d,a,\Wt)$ only. On King's small-field set $|\phi|\le M_k'$ the last term is $\le Ce_k^2p_k^2M_k'^2|\Lambda|=O(\ell_k^{\kappa_0})|\Lambda|$; at $A=0$ it is absent. This is \cite[(3.29)]{Ba2} at the weight.
\end{lemma}
\begin{proof}
(1) Dirichlet faces only increase the form: the Dirichlet form on the live sites is the Neumann form plus the face terms $\vth^{d-2}|\phi(x)|^2\ge0$, so $\Delta^{(k)}_D\ge\Delta^{(k)}_N$ by the monotonicity of the minimiser identity. (2) $\ip\phi{\Delta\phi}=\min_\Phi[\ip\Phi{H\Phi}+a_k\norm{\phi-Q\Phi}^2]$. For a block bond $\langle y,y'\rangle$: $|\phi(y)-U\phi(y')|^2\le3(|\phi(y)-Q\Phi(y)|^2+|Q\Phi(y)-UQ\Phi(y')|^2+|\phi(y')-Q\Phi(y')|^2)$. On the two-block box $B(y)\cup B(y')$ (simply connected) gauge to the constant field and compare: the averages and covariant differences at $A$ and at the constant field differ by $\le C\theta_k\norm\Phi_{B\cup B'}$, and $\norm\Phi_B^2\le C(E_B(\Phi)+|\phi(y)|^2+|\phi(y)-Q\Phi(y)|^2)$ (block Poincar\'e with the weight); the block-bond transport $U(A_k)$ differs from the fine one by $Ce_kp_k$ (the photon block-versus-background step \cite[(A.2)--(A.6)]{K1}, which is $\kappa$-free and holds on the cube by \cref{lem:a01:PE,lem:a01:REG}). At $A=0$, $|\operatorname{avg}_wg-\operatorname{avg}_{w'}g|^2\le4\Wt C_P(d)E_{B\cup B'}(g)$ by Cauchy--Schwarz with the weights ($\sum\vth^dw^2=\Wt$ on a unit block) and the two-block Poincar\'e inequality. Mass: $V_k|\phi(y)|^2\le2V_k(\Wt\norm\Phi^2_{B(y)}+|\phi(y)-Q\Phi(y)|^2)$. Summing (each block term counted $\le2d$ times, each fine bond in $\le2$ boxes) and absorbing the $\theta_k$-parts gives the claim with $\gamma_0^w=[C(d)\max(4\Wt C_P,2\Wt,(6d+2)/a_k)]^{-1}(1-O(\theta_k))$.
\end{proof}

\begin{lemma}[the inherited face cap]\label{lem:a01:IC}
Let $k\ge1$ and fix a term of the step-$(k-1)$ large-field decomposition (a family of events with its bad union $S_{k-1}$). On the support of its characteristic functions (Ba{\l}aban's six families on $\Lambda_7^{(k-2)\prime}$ and the face family of scale $k-1$; at $k-1=0$ on the first live layer of fine sites), every face $k$-block $y\in\Lone^{(k)}$ with $B(y)\subset\Lambda_0^{(k-1)}$ satisfies, in unit-$k$ normalisation,
\[
|\phi_k(y)|\ \le\ L^{1/2}(\sigma_{k-1}+Lp_{k-1})\ \le\ C_{\mathrm{IC}}\,p_k,\qquad C_{\mathrm{IC}}\defeq L^{1/2}(1+\ln L)^p(C_8+L),
\]
free of $\eps$, $k$, $K$, $J$ and of $M$. By the nesting of the regions the cap holds at every face block of $\Lambda_2^{(k-1)\prime}$, in particular at every face block of the propagator domain $B^k(\Lambda_2^{(k-1)\prime})$, of the ring $\Lambda_6^{(k-1)\prime}\setminus\Lambda_7^{(k-1)\prime}$ and of $S_k$.
\end{lemma}
\begin{proof}
Every step-$(k-1)$ test located in $B(y)$ is imposed and passes: the face test at the $(k-1)$-blocks of $B(y)$ in $\Lone^{(k-1)}$ (their index set contains them since $\Lambda_0^{(k-1)}\subset\Lambda_7^{(k-2)\prime}$), the gradient test on the block bonds with both ends in $B(y)$ (a violating bond would put its endpoints in $S_{k-1}$), and the new-field test $|\phi_k(y)-(Q\phi_{k-1})(y)|\le p_{k-1}$ at $y$ ($y$ lies in the prime of $\Lambda_7^{(k-2)\prime}$ and $y\in P_s^{(k-1)}$ would give $B(y)\subset S_{k-1}$). Since $B(y)$ is a face block, some $(k-1)$-block $y''\subset B(y)$ is a face block with $|\phi_{k-1}(y'')|\le\sigma_{k-1}$; the gradient test along a block path inside $B(y)$ (length $\le L$) gives $|\phi_{k-1}|\le\sigma_{k-1}+Lp_{k-1}$ on all of $B(y)$, hence $|(Q\phi_{k-1})(y)|\le\sigma_{k-1}+Lp_{k-1}$ (Q5), and the new-field test gives $|\phi_k(y)|\le\sigma_{k-1}+(L+1)p_{k-1}$ in unit-$(k-1)$ normalisation, i.e.\ the claim after the conversion factor $L^{1/2}$ ($d=3$) and $p_{k-1}\le(1+\ln L)^pp_k$. At $k-1=0$ the step-$0$ face member is the cap on $F_1$ of \cref{def:a01:facefam}.
\end{proof}

Without the cap the fine-field class of \cite[(2.55)]{Ba2} fails at the collar cubes: for constant data of size $M_k$ the background's face gradient is of order $M_k$, not $p_k$. The cap is inherited from the previous step exactly as (2.55) is carried from step to step, as a characteristic function of the block field.

% a01_H0rel part 5: the perturbative input and the lower bound
\subsection{The perturbative input on the cube}\label{app:a01:TD}

Ba{\l}aban's perturbative bounds \cite[Prop.~1]{Ba3} are consumed, in the lower bound at every step and in the upper bound at every scale, in the field-dependent form: for a renormalised family $\mathcal R$ of order $\le\bar n=13$ of the expansion of the effective action, built from the lines (slices $G_{(i)}$ and $G_k$) and legs ($\phi^{(k)}$, photon backgrounds) of the step, with external fields in the class of \cref{lem:a01:legs} and possibly source vertices, the sum over orderings, indices and positions of the values of its members, with one vertex $x_0$ held, is bounded by $C\ell_k^{\beta_0}p(\ell_k)^{n_1}(1+K-k)^{p_0}e^{-\delta\dist(V_*)}$ times the external-field norms, with $V_*$ the set of vertices carrying external legs or sources and $\beta_0=\frac14-\beta$; the $R$-terms (orders $>\bar n$) are $O(1)\ell_k^{\kappa'}|\Lambda^{(k)}|$ with $\kappa'>d$. On the torus this is \cite[Prop.~1]{Ba3} with the hypotheses on the external fields supplied by \cite[Lemma~2.4]{Ba2}.

\begin{proposition}[the perturbative bounds on the cube]\label{prop:a01:TD}
On $\CK$ with King's flow, at every scale $k\le K_f$ and every regular $A$: (a) on the whole cube, the bound above holds for every renormalised family with the cube's lines and legs, with the counterterm King's free-boundary $\dm$ (\cite[p.~330]{K2}, field- and domain-independent), in the organisation of \cite[p.~330]{K2}: for a held vertex $x_0$ at face distance $t(x_0)\ge1$ (zone I) with the torus constants; for $t(x_0)<1$ (zone B) with the additional factor $\Lambda(x_0)\le C(1+\ln_+(1/t(x_0)))^{14}t(x_0)^{-(1-\beta)}e^{-\delta_\Lambda t(x_0)}$, which is summable: $\sum_{x_0\in\text{unit cube}}\vth^d\Lambda(x_0)\le C$; (b) the same on every sub-domain $\Omega$ of the cube (any union of $k$-blocks, meeting $\DD$ or not) for the families with the lines $G_{(i)}(\Omega,A)$, $G_k(\Omega,A)$ and the legs $\phi^{(k)}_\Omega$, and with one line replaced by a difference line $G_k(\Omega,A)-G_k(\Omega_2,A)$ of two Neumann domains as an external pair, provided every vertex cube lies at depth $\ge R_*$ in every propagator domain of the family; in the difference case the bound carries the extra factor $\exp[-\delta_c(\dist(\mathrm{cube}(v),\Xl\setminus\Omega_2)+\dist(\mathrm{cube}(v'),\Xl\setminus\Omega_2))]$ for the two ends $v,v'$ of the difference line; (c) the consequences \cite[(3.58)]{Ba1} hold with these inputs: the coefficients of the effective potential $V^{(k)}$ are $O(1)\ell_k^{\beta_0}$ and the remainders $O(1)\ell_k^{\kappa'}|\Lambda^{(k)}|$, uniformly in $k$, $K$, $\vth$ and the field.
\end{proposition}
\begin{proof}
(a) Zone I is the torus bound: at face distance $\ge1$ every line is, up to the face-localised difference of \cref{lem:a01:FLD} (exponentially small in the distance to the face and absorbed in the decay factors), the torus line at the same field, and the photon lines are the image sums of \cref{lem:a01:NL}; the class of the external legs is \cref{lem:a01:legs}. Zone B is the content of the boundary-remainder bounds established in the proof of \cref{prop:k01:Kpos} (the scalar boundary remainders at zero background and the photon face terms of \cite[pp.~331--332]{K2}, with the slice bounds $|D^{n_x}_xD^{n_y}_yG^\Omega_{(i)}(b,b')|\le C\ell_i^{-1-n_x-n_y}e^{-\kappa|x-y|/(2\ell_i)}$ and the cumulative line bound $|D^{n_x}_xD^{n_y}_yG^\Omega_{\le i_0}(b,b')|\le C|x-y|_\vth^{-1-n_x-n_y}e^{-\kappa|x-y|/(4\ell_{i_0})}$ at zero background on the cube and on King's localisation boxes, face bonds included, which follow from the odd-image representation of Step~1 of \cref{lem:a01:FL} and the free weighted bounds; and, at regular $A$, from \cref{lem:a01:A23,lem:a01:A4,lem:a01:A5} through the gauge-away of the constant part on each box): there every line is written as its free counterpart plus a boundary remainder localised at the face, and the remainder vertices are of higher degree, so that the degree bookkeeping of \cite[Sect.~2]{Ba3} is unchanged while the sum over the position of the held vertex produces $\Lambda(x_0)$. The photon-image part of the counterterm is the photon-image part of the same remainder families (\cref{app:a01:photon}). (b) is \cref{thm:a01:B2} below. (c) is \cite[p.~421]{Ba3}: the coefficient bounds use of the external scalar field exactly the sup clause, the gradient clause (P-grad) and the block clause (P-block) of \cref{lem:a01:PC}, at the current vertex \cite[(1.8)]{Ba3} and at the first two kernel vertices \cite[(1.18)]{Ba3}, every other family using the sup clause alone; on the cube the current vertex sits on live bonds only (a face bond carries no link) and the kernel vertices at weighted sites only (the weight vanishes at killed sites), so the two clauses are needed only there.
\end{proof}


% a01_H0rel part 5b: the lower bound
\subsection{The lower bound}\label{app:a01:lower}

Following \cite[Sect.~3]{Ba1} and \cite[Appendix]{K1}: $p_i\defeq p(\ell_i)$, $M_i\defeq p_i\lambda(\ell_i)^{-1/4}$, $M_i'\defeq p_{i-1}\lambda(\ell_{i-1})^{-1/4}=L^{1/4}(p_{i-1}/p_i)M_i$ the block-sup threshold of \cite[(3.2)]{K1}, $p_1(\ell)\defeq b_1(1+\ln\ell^{-1})^{p_1}$ the fluctuation cut-off with $p-1\ge p_1$.

\begin{definition}[the class carried through the induction]\label{def:a01:class}
In unit-$i$ units, with $\phi^{(i)}\defeq a_i\Mop(\Xl,A^{(i)})^{-1}Q(A^{(i)})^*\phi_i$, the class $(\mathrm{A.1})^{(i)}_C$ is defined by
\[
\sup_{x\in\Xl}|\phi^{(i)}(x)|\le C_1M_i,\qquad\sup_y|\phi_i(y)|\le C_1M_i,
\]
\[
(\mathrm G)_i:\ \sup_{b}|D_{A^{(i)}}\phi^{(i)}(b)|\le C_1'p_i\ \text{ over all bonds of the cube, face bonds included},\qquad(\mathrm E)_i:\ \sup_y|\phi_i(y)-Q(A^{(i)})\phi^{(i)}(y)|\le C_1''p_i,
\]
together with the photon clause $|\partial A^{(i)}|_\infty\le C'_\Gamma p_i$. On a face bond $(\mathrm G)_i$ reads $\vth^{-1}|\phi^{(i)}(x)|\le C_1'p_i$ at the adjacent live site: a fine-level linear vanishing with slope $O(p_i)$.
\end{definition}

The first two clauses are King's (A.1) with its Laplacian clause dropped; $(\mathrm G)$ and $(\mathrm E)$ are the two consequences of that clause which the lower bound consumes (\cref{lem:a01:PC}). On the cube the Laplacian clause itself holds as a consequence ($(-\Delta^D_A+V_i)\phi^{(i)}=a_iQ^*(\phi_i-Q\phi^{(i)})$ gives $|\Delta^D_A\phi^{(i)}|\le V_iC_1M_i+a_i\Wt C_1''p_i$), but it is not needed.

\begin{lemma}[compatibility, coefficient one]\label{lem:a01:GC}
Let $\psi\in(\mathrm{A.1})^{(i+1)}_C$, $|\phi_i'|_\infty\le p_1(\ell_i)$ and $\phi_i=aL^{-2}C^{(i)}Q_c^*\psi+\phi_i'$ at the gauge field of Ba{\l}aban's step ($A^{(i)}=A^{(i+1)}+A'^{(i)}$). Then $\phi_i\in(\mathrm{A.1})^{(i)}_C$ provided $L^{3/4}\ge3$, $C_1\ge3\max(C_A,1)\max(1,b_1/b_0)$, $C_1'\ge3C_D\max(1,b_1/b_0)$, $C_1''\ge\max\{2(1+C_{Q4})\max(1,b_1/b_0),C_1'\}$, and the three smallness conditions on $\eps_0^{\mathrm{Ba}}$ for the gauge-field remainders (below), where $C_A$, $C_D$ are the constants of \cref{lem:a01:A4}(b) and \cref{lem:a01:A5} for the leg of block data ($C_A,C_D\le Ca\Wt$) and $C_{Q4}\le C_A$ the constant in $|Q(A)u|_\infty\le C_{Q4}|b|_\infty$.
\end{lemma}
\begin{proof}
(1) \emph{Exact composition.} By the composition law of the kernels and the minimiser identity (the minimiser of the composed functional is the minimiser of the one-step functional with the composed average, the uniform coarse step having $Q_cQ_c^*=1$), $\phi^{(i)}(\phi_i)=\Phi_\psi+u'$ with $\Phi_\psi\defeq\phi^{(i+1)}(\psi)$ and $u'\defeq a_i\Mop^{-1}Q^*\phi_i'$, $|u'|_\infty\le C_A|\phi_i'|_\infty$. (2) \emph{Units.} A field in unit-$(i+1)$ units is $L^{1/2}$ times the field in unit-$i$ units and a length is $L$ times, so a fine-bond difference quotient converts by $L^{-3/2}$ and a block average by $L^{-1/2}$; $p_{i+1}\le p_i$, $\lambda(\ell_{i+1})=L\lambda(\ell_i)$. (3) \emph{Sup clauses.} $|\phi^{(i)}|_\infty\le|\Phi_\psi|_\infty+C_A|\phi_i'|_\infty\le C_1L^{-3/4}M_i+C_Ap_1(\ell_i)\le\frac13C_1M_i+\frac13C_1M_i$, the third share being left for the gauge-field remainder. Block clause: by the blockwise push-down identity $aL^{-2}C^{(i)}Q_c^*\psi=Q\Phi_\psi+\kappa'Q_c^*(\psi-Q_{i+1}\Phi_\psi)$ with $\kappa'=aL^{-2}/(a_i+aL^{-2})\le L^{-2}$ (minimise the composed functional over $\phi_i$ block by block: on each $(i+1)$-block the minimiser is $Q\Phi+\delta$ with $\delta$ constant on the block, and the one-dimensional minimisation gives $\delta=\kappa'(\psi-Q_{i+1}\Phi)$; this uses only that the coarse step is uniform), and $|Q\Phi|\le\sup|\Phi|$ (Q5), $|Q_c^*c|=|c|$: $|\phi_i|_\infty\le(1+2\kappa')C_1L^{-3/4}M_i+p_1(\ell_i)\le C_1M_i$. (4) $(\mathrm G)_i$ at live bonds: $|D\Phi_\psi(b)|\le L^{-3/2}C_1'p_{i+1}\le\frac13C_1'p_i$ by $(\mathrm G)_{i+1}$ at the same fine bond, and $|Du'(b)|\le C_D|\phi_i'|_\infty\le\frac13C_1'p_i$ by \cref{lem:a01:A5}; the third share is the gauge-field remainder. At face bonds the same three terms, each a face-bond derivative: $D\Phi_\psi$ by $(\mathrm G)_{i+1}$ on the same face bond, $Du'$ by the face-bond clause of \cref{lem:a01:A5}, and the remainder likewise; the difference $D_{B+A'}\Phi_\psi-D_B\Phi_\psi$ vanishes on a face bond (no link). (5) $(\mathrm E)_i$: by the push-down identity and the exact composition,
\[
\phi_i-Q(A)\phi^{(i)}(\phi_i)=\kappa'Q_c^*\big(\psi-Q_{i+1}\Phi_\psi\big)+\big(\phi_i'-Q(A)u'\big),
\]
and at every block $|\cdot|\le\kappa'L^{-1/2}C_1''p_{i+1}+(1+C_{Q4})p_1(\ell_i)\le L^{-5/2}C_1''p_i+\frac12C_1''p_i\le C_1''p_i$. (6) \emph{The gauge-field remainder.} The composition holds at a common gauge field; in Ba{\l}aban's step the push-down is taken at $B\defeq A^{(i+1)}$ while $\phi^{(i)}$ is built with $B+A'$. The background of $\psi$ at $B+A'$ differs from $\Phi_\psi$ by $\delta$ solving $\Mop(B+A')\delta=f$, $f=a_i[Q(B+A')^*-Q(B)^*]\omega_B-[\Mop(B+A')-\Mop(B)]\Phi_\psi$ with $\omega_B$ the pushed-down datum, $|\omega_B|_\infty\le C_1M_i$; expanding term by term, $f$ is a bounded function with $|f|\le Ce_ip_1(\ell_i)(C_1M_i+C_1'p_i)$ on Ba{\l}aban's fluctuation class ($|A'|,|\partial A'|\le C_\gamma p_1(\ell_i)$ by \cite[(2.18)]{K1} with \cite[Thm.~3.3]{K1} for the vector operators and the cut-off \cite[(3.50)]{Ba1}; on the cube by \cref{lem:a01:PE,lem:a01:REG}), involving only live bonds and contours inside blocks; by \cref{lem:a01:A4}(c) and \cref{lem:a01:A5} (bounded source), $|\delta|_\infty+\sup_b|D_A\delta(b)|\le C_{\mathrm{rem}}e\,p_1(\ell_i)(C_1\lambda^{-1/4}\ell_i^{1/4}+C_1'\ell_i^{1/2})p_i$, which is $\le\frac13C_1'p_i$ and $\le\frac13C_1M_i$ under $C_{\mathrm{rem}}e\,C_1\lambda^{-1/4}\beta^{1/4}p_1(\beta)\le\frac16C_1'$, $C_{\mathrm{rem}}e\,\beta^{1/2}p_1(\beta)\le\frac16$ and $C_{\mathrm{rem}}e\,p_1(\beta)(\lambda^{-1/4}\beta^{1/4}+(C_1'/C_1)\beta^{1/2})\le\frac13$ at $\beta=\eps_0^{\mathrm{Ba}}$ (monotone in $\beta$ on $(0,e^{1-4p_1}]$, so valid at every $\ell_i\le\eps_0^{\mathrm{Ba}}$); the block-averaged form of the remainder enters $(\mathrm E)_i$ with the sup clause of \cref{lem:a01:A4}(c) and is absorbed in the room $\frac12-L^{-5/2}\ge\frac13$ (here $C_1''\ge C_1'$ is used). The scale-$(i+1)$ class enters with coefficient one; no constant multiplies $p_{i+1}$, which is what \cite[(3.53)--(3.54)]{Ba1} requires, and the conditions (3.54) are unchanged.
\end{proof}

\begin{corollary}[King's face test as a consequence at every scale of the lower bound]\label{cor:a01:BR}
On $(\mathrm{A.1})^{(i+1)}_C$, at every face block $y$ of scale $i$, the push-down $b_i\defeq aL^{-2}C^{(i)}Q_c^*\psi$ obeys $|b_i(y)|\le(1+\vth)L^{-3/2}C_1'p_i+L^{-5/2}C_1''p_i$, hence $|\phi_i(y)|\le$ that $+p_1(\ell_i)$ on the induction set.
\end{corollary}
\begin{proof}
$b_i=Q\Phi_\psi+\kappa'Q_c^*(\psi-Q_{i+1}\Phi_\psi)$. $\Phi_\psi$ vanishes on the killed layer; along a normal path of length $\le1+\vth$ from the killed site to any point of $B(y)$ every bond is a live or face bond, so $|\Phi_\psi|\le(1+\vth)\sup|D\Phi_\psi|\le(1+\vth)L^{-3/2}C_1'p_i$ by $(\mathrm G)_{i+1}$ including its face part; $Q$ is an average; the second term is $(\mathrm E)_{i+1}$ with $\kappa'L^{-1/2}\le L^{-5/2}$.
\end{proof}

This is a consequence, not an inserted test: no block-field test is inserted at an intermediate scale of the lower bound (such a test would not be implied by the previous scale, the push-down norm having no numerical bound below one); King's remark that there is ``no change in the proof of the lower bound'' \cite[p.~343]{K2} is correct in this sense once King's Laplacian clause is replaced by $(\mathrm G)$, $(\mathrm E)$.

\begin{lemma}[the two properties consumed by the $n+n'=1$ vertices]\label{lem:a01:PC}
On $(\mathrm{A.1})^{(i)}_C$, at every block $y$, every live $x\in B(y)$ and every live bond $b$: (P-grad) $|D_{A^{(i)}}\phi^{(i)}(b)|\le C_1'p_i$; (P-block) $|\phi^{(i)}(x)-U(A^{(i)}(\Gamma_{x,y}))\phi_i(y)|\le C_Pp_i$ with $C_P\defeq C_1''+2dC_1'$.
\end{lemma}
\begin{proof}
(P-grad) is $(\mathrm G)_i$. (P-block): $Q\phi^{(i)}(y)=\sum_{x'}\omega(x')U(\Gamma_{y,x'})\phi^{(i)}(x')$ with $\omega\defeq\vth^3w\ge0$, $\sum\omega=1$; so $\phi^{(i)}(x)-U(\Gamma_{x,y})\phi_i(y)=U(\Gamma_{x,y})[Q\phi^{(i)}(y)-\phi_i(y)]+\sum_{x'}\omega(x')[\phi^{(i)}(x)-U(\Gamma_{x,y})U(\Gamma_{y,x'})\phi^{(i)}(x')]$; the first term is $(\mathrm E)_i$, and each term of the sum is the covariant difference along the path $\Gamma_{y,x'}$ followed by $\Gamma_{x,y}$, of length $\le2d$, the exact telescoping sum of $\vth D_A\phi^{(i)}$ over its bonds (live sites only, by the choice of contours), hence $\le2dC_1'p_i$ by $(\mathrm G)_i$.
\end{proof}

\begin{lemma}[final restriction at the stopping scale]\label{lem:a01:final}
Let $\chi_k\lambda_k$ be the small-field indicator of \cref{app:a01:weights} at the stopping scale (King's \cite[(3.2)]{K1} with the gradient clause on every block bond of the cube, times the face test). Then $\supp(\chi_k\lambda_k)\subset(\mathrm{A.1})^{(k)}_C$ once $C_1\ge\max(C_A,1)L^{1/4}(1+\ln L)^p$, $C_1'\ge C_G(1+\ln L)^p$, $C_1''\ge C_E(1+\ln L)^p$, $V_kM_k'\le p_k$ and $e_kM_k\le1$ (conditions on $\eta_1$), with $C_G,C_E$ depending on $(C_8,\delta_1,\delta',C_A,L,a,d)$ only. At the final scale $K_f$ the set $\{|\phi_{K_f}(y)|\le1\ \forall y\}$ lies in $(\mathrm{A.1})^{(K_f)}_C$.
\end{lemma}
\begin{proof}
Facts on the final set: (f1) $|\phi_k(y')|\le C_8p_k$ on $\Lone^{(k)}$; (f2) $|D_{A_k}\phi_k(\beta)|\le p_{k-1}$ on every block bond; (f3) $|\phi_k|\le M_k'$; hence (f5) $|\phi_k(y)|\le C_9(1+\dist(y,\Lone^{(k)}))p_{k-1}$, $C_9\defeq\max(C_8,1)$ (the cube is a box: every block is joined to the face layer by a cube-internal block path of that length), and (f6) $|\phi_k(y)-U(A_k(\Gamma^{\mathsf T}_{y,y_0}))\phi_k(y_0)|\le|y-y_0|_1p_{k-1}$ along a tree $\mathsf T$ of block paths. Sup clauses: (f3) and \cref{lem:a01:A4}(b), with $M_k'\le L^{1/4}(1+\ln L)^pM_k$. $(\mathrm G)_k$ at a bond $b$ (live or face) in the block $y_0$: write $\phi_k=\phi_k(y_0)\Theta^{\mathrm{blk}}+(\phi_k-\phi_k(y_0)\Theta^{\mathrm{blk}})$ with $\Theta^{\mathrm{blk}}$ the unit tree datum transported by $A_k$. Constant part: \cref{lem:a01:FL}(a) with $|b-y_0|\le1$, the tree datum transported with the block field $A_k$ differing from $\Theta_A$ by $\le Ce_kp_k|y-y_0||c|$ (the photon block-versus-background step \cite[(A.2)--(A.6)]{K1}, on the cube by \cref{lem:a01:PE,lem:a01:REG}): $\le|\phi_k(y_0)|[CV_k+Ce^{-\delta_1\dist(b,\DD)}+C\theta_k+Ce_kp_k]\le C[V_kM_k'+C_9\sup_t(1+t)e^{-\delta_1(t-1)}p_{k-1}+(\theta_k+e_kp_k)M_k']\le C'p_k$ by (f3), (f5), $\dist(b,\DD)\ge\dist(y_0,\Lone^{(k)})-1$, $V_kM_k'\le p_k$, $\theta_k\le Ce_kp_k$ and $e_kp_kM_k'\le C_1p_k$. Varying part: $\sum_y|D_{A,b}l(b,y)||\phi_k(y)-\phi_k(y_0)\Theta^{\mathrm{blk}}(y)|\le\sum_yCe^{-\delta'|b-y|}|y-y_0|_1p_{k-1}\le Cp_{k-1}$ by \cref{lem:a01:A5} in kernel form and (f6). $(\mathrm E)_k$: by \cref{lem:a01:EB}(a), $|\phi_k(y)-Q\phi^{(k)}(y)|=a_k^{-1}|\Delta^{(k)}\phi_k(y)|$; with the tree rooted at $y$, $\Delta^{(k)}\phi_k(y)=\sum_{y'}\Delta^{(k)}(y,y')(\phi_k(y')-\phi_k(y)\Theta^{\mathrm{blk}}(y'))+\phi_k(y)(\Delta^{(k)}\Theta^{\mathrm{blk}})(y)$; the kernel decays, $|\Delta^{(k)}(y,y')|\le Ce^{-\delta'|y-y'|}$ (\cref{lem:a01:A23}(b)), so the first sum is $\le Cp_{k-1}$ by (f6), and the second is $a_k|\phi_k(y)||\Theta^{\mathrm{blk}}(y)-Qu(y)|\le a_k|\phi_k(y)|[CV_k+Ce^{-\delta_1\dist(B(y),\DD)}+C(\theta_k+e_kp_k)]\le Cp_k$ by \cref{lem:a01:FL}(b), (f3), (f5). At the final scale the data are bounded by $1$ and $(\mathrm G)_{K_f}$, $(\mathrm E)_{K_f}$ follow from \cref{lem:a01:A5,lem:a01:A4}(b) with $|b|_\infty\le1$, since $p_{K_f}\ge b_0\ge1$.
\end{proof}

The face test is what places the face blocks in the class: without it, constant data of size $M_k'$ violate $(\mathrm G)_k$ near the face by the factor $M_k'/p_k$ (\cref{lem:a01:FL}(a) is sharp at the face).

\paragraph{The steps of the lower bound.} We follow \cite[Sect.~3]{Ba1} with King's modifications \cite[Appendix, p.~676]{K1}; $E_1$ is any field-independent constant (\cref{lem:a01:cnum}).
\begin{description}[leftmargin=1.5em,style=nextline]
\item[{B1 (restriction to small fields; \cite[(3.6)]{Ba1})}] The integrand of $\Zk(s)$ is positive, so inserting any indicator at any step gives a lower bound. We insert, at each step $i<k$, the indicator of the class $(\mathrm{A.1})^{(i)}_C$ (a class on the backgrounds, the block sup and the block mismatch, not King's six block-field tests), and at the stopping scale $k$ the indicator $\chi_k\lambda_k$, a subset of $(\mathrm{A.1})^{(k)}_C$ by \cref{lem:a01:final}. Nothing is paid in a lower bound.
\item[{B2 (inductive hypothesis; \cite[(3.26)--(3.32)]{Ba1})}] After $i$ steps, $\Zk\ge e^{-a_i}\int\chi_{(\mathrm{A.1})^{(i)}_C}e^{-S^{(i)}_k}$ with $a_i\defeq\sum_{i'<i}C\ell_{i'}^{\kappa'}|X^{(i')}|$ and $S^{(i)}_k$ King's effective action \cite[(3.14)]{K1} built on the cube; the identities for the normalisations and the representation of the fields after $i$ steps are \cite[(2.13)--(2.17)]{K1}, exact for every weight and boundary condition (the block-spin RG is a Gaussian identity for any positive covariance; the composition law uses only that the coarse steps are uniform).
\item[{B3 (the fluctuation step; \cite[(3.37)--(3.49), Props.~3.1--3.2]{Ba1})}] The Gaussian fluctuation integral is exact; the interaction is expanded to order $\bar n=13$ \cite[Appendix]{K1} with Cauchy estimates in the complexified couplings; the $R$-terms and the coefficient bounds are \cref{prop:a01:TD}(c) under the hypotheses supplied by \cref{def:a01:class,lem:a01:PC}; the one-loop determinants are bounded below by $\ln(1+x)\le x-x^2/2+\dots$ ($n$ odd) for the bounded operators $C^{(i)}$ of \cref{lem:a01:A23}(c).
\item[{B4 (cut-off removal and the consistency conditions; \cite[(3.50)--(3.55)]{Ba1})}] The cumulant bound $\int\chi e^Vd\mu\ge\exp[\sum_{n\le\bar n}E^T_n(V)/n!-R]$ with cut-off-free cumulants and $R=O(1)\ell_i^{\kappa'}|X^{(i)}_1|$, $\kappa'>d$, for a Gaussian measure with exponentially decaying precision (\cref{lem:a01:A1,lem:a01:A23}) and a real polynomial interaction with Ba{\l}aban's two-component cut-off: this is the lemma behind \cite[Prop.~3.2]{Ba1}, with its interaction input \cref{prop:a01:TD}(c); the compatibility inequality $\chi_{(\mathrm{A.1})^{(i+1)}_C}(\psi)\chi(|\phi_i'|\le p_1(\ell_i))\le\chi_{(\mathrm{A.1})^{(i)}_C}(\phi_i)$ is \cref{lem:a01:GC}, with coefficient one on the scale-$(i+1)$ class, the rooms $L^{3/4}$ (sup clauses), $L^{3/2}$ ($(\mathrm G)$), $L^{5/2}$ ($(\mathrm E)$), and the conditions \cite[(3.54)]{Ba1} unchanged.
\item[{B5 (King's Appendix: \cite[(3.2)]{K1} in place of \cite[(3.27)--(3.29)]{Ba1})}] King's Laplacian clause is replaced by $(\mathrm G)_i$, $(\mathrm E)_i$, which compose with coefficient one (\cref{lem:a01:GC}) and supply through \cref{lem:a01:PC} exactly the two properties of the external scalar field that \cite[p.~421]{Ba3} draws from \cite[Lemma~2.4]{Ba2} for his class; what \cite[Prop.~1]{Ba3} uses of the external legs otherwise is their norm, \cref{lem:a01:legs}(a). The final set $\chi_k\lambda_k$ is in the class by \cref{lem:a01:final}. The photon background's clauses \cite[(A.2)--(A.6)]{K1} are unchanged on the cube (\cref{app:a01:photon}).
\item[{B6 (the final step at $\ell_f=\ell_{K_f}$; \cite[(3.66)--(3.69)]{Ba1})}] Expand in the small $A^{(K_f)}$ (\cref{prop:a01:TD}(c)); $|V^{(K_f)}|\le O(1)\ell_f^{\beta_0}|X^{(K_f)}|$; the finite-dimensional integral over $\{|\phi_{K_f}(y)|\le1,|A_{K_f}|\le1\}\subset(\mathrm{A.1})^{(K_f)}_C$ (\cref{lem:a01:final}) of $\exp[-\frac12\ip{A_{K_f}}{\Delta^{(K_f)}A_{K_f}}-\frac12\ip{\phi_{K_f}}{\Delta^{(K_f)}\phi_{K_f}}]$ is $\ge\exp(-O(1)|X^{(K_f)}|)$, $|X^{(K_f)}|=\ell_f^{-d}|\CJ|$, using only the upper bounds $\Delta^{(K_f)}\le a_{K_f}\ell_f^{-2}$ on the integration region; and the normalisation identity \cite[(3.69)]{Ba1} gives $Z_{K_f}Z_{K_f}(0)e^{-E_0}\ge\exp[-\frac12a_{K_f}\ell_f^{d-2}(\Tr P_{K_f}G+\Tr P_{K_f}G^\gamma)]$ with $\Tr P_{K_f}G\le W_2\Tr(\Pi G)\le W_2N|X^{(K_f)}|\norm G$ (Q3). Hence $\Zk(0)\ge e^{-E_-|\CJ|}$ in the normalisation in which $E_1$ is the domain-adapted constant of \cite[(3.35)]{K1} on $\CJ$, with $E_-=O(\ell_f^{-d})$ depending on $\eps_0^{\mathrm{Ba}}$ and King's data only.
\item[B7 (the source)] $F(u_x)$ is King's plaquette source \cite[(2.7)]{K1}, whose vertices \cite[(3.20)]{K1} sit at $\ell^\infty$-distance $>1$ from $\R^3\setminus\CJ$ (zone I), so the face enters them only through propagators (\cref{prop:a01:TD}).
\end{description}

\begin{proposition}[lower line]\label{prop:a01:lower}
For $\eta\le\eta_1$ (so that \cite[(3.54)]{Ba1} holds at all scales $\le k$ and $j\le k-J_*$) and $j\ge J_*$: $\Zk(s)\ge e^{-a}\Ik j(s)$ with $a=\sum_{i<j}C\ell_i^{\kappa'}|X^{(i)}|=C\eta^\sigma|\CJ|$, $\sigma=\kappa'-d>0$; and $\Zk(0)\ge e^{-E_-|\CJ|}$ in the normalisation of B6.
\end{proposition}
\begin{proof}
B1--B7. The sum $\sum_{i<j}\ell_i^{\kappa'}|X^{(i)}|=|\CJ|\sum_{i<j}\ell_i^{\kappa'-d}$ is geometric and dominated by $i=j-1$, giving $a=C\eta^\sigma|\CJ|$ with $C$ independent of $j$ and $k$ (\cref{prop:a01:TD}, \cref{lem:a01:GC}, whose constants are free of $k$, $j$, $\vth$); the final restriction is a positivity step inside the class (\cref{lem:a01:final}), so the threshold $\sigma_k$ never enters a constant.
\end{proof}

% a01_H0rel part 6: the upper bound, the assembly, the clause (R_0)
\subsection{The upper bound}\label{app:a01:upper}

We follow \cite[Sects.~2--3]{Ba2} and \cite[Sect.~5.2]{K2}. Ba{\l}aban's six large-field families (the tests on $|B-QA|$, $|\partial A|$, $|A|$, $|\psi-Q(A)\phi|$, $|D_A\phi|$, $|\phi|$ of \cite[(2.2)]{Ba2}) are used at every scale $i\le K_f$, with the face family of \cref{def:a01:facefam} as a seventh family at every scale, $i=0$ included; at the stopping scale $j$ the small-field term carries $\chi_j\lambda_j$. The large-field decomposition, its regions $\Lambda_0^{(i)}\supset\Lambda_1^{(i)}\supset\dots\supset\Lambda_7^{(i)}$ (layers of thickness $r(\ell_i)$), the nesting rules and the inductive representation \cite[(2.43)--(2.52)]{Ba2} depend only on the set of event blocks and the number of event types, so they hold verbatim with the face family, which is a pointwise condition on one block variable at one block of the index set $\Lone^{(i)}\cap\Lambda_7^{(i-1)\prime}$. Four steps of Ba{\l}aban's printed proofs do not hold as printed at block factor $5$ and are used in a replaced form: the cube step of \cite[p.~433]{Ba3}, the cut-off $\zeta^{(k)}$ of \cite[(2.44)]{Ba2}, the boxes in the proof of \cite[Lemma~2.4]{Ba2}, and the cube partition \cite[(3.43)--(3.45)]{Ba2}, the last replaced by a product partition into boxes of large blocks with sides in $[r+s_b,2r+3s_b]$, which has every property used in \cite[Sect.~3C]{Ba2}; the replaced objects are defined from $(L,M,R,\ell_i,\bar n)$ alone and do not depend on the matter boundary condition, and the replacements change none of the printed statements. Their block-factor-$5$ form, with proofs, follows.
% a01_H0rel part 6x1: the four replaced steps of Balaban's proofs at block factor 5 (input from a01_H0rel_p6.tex)
\paragraph{The four replaced steps at block factor $5$.}
Scale-$i$ units (the block lattice of scale $i$ has spacing $1$); $u_i\defeq1+\ln\ell_i^{-1}$, so $u_{i-1}=u_i+\ln L$, $r_i\defeq r(\ell_i)=Ru_i^{\varrho}$, $p_i=p(\ell_i)=b_0u_i^{p}$ with $\varrho=2$, $p=6$; $s_b=LM$ is the large-block side; $\rho_i\defeq(r(\ell_{i-1})-s_b)/L$ is one previous-scale corridor in scale-$i$ units (\cite[p.~578]{Ba2}); $X'$ denotes the blocks of $X$ read as points of the next lattice \cite[(1.16)]{Ba1}; $\Omega\defeq B^i(\Lambda_2^{(i-1)\prime})$ is the propagator domain and $\Omega_1\defeq B^i(\Lambda_7^{(i-1)\prime})$ the vertex set of \cite[(1.1)--(1.5)]{Ba3}; $\delta_0$, $c_{\mathrm{Ba}}$ and the depth $R_{\mathrm{Ba}}$ are the constants of \cite[Props.~2.1--2.3]{Ba1}. Everything below is a statement about Ba{\l}aban's regions, cut-offs, boxes and partitions, which are built from $(L,M,R,\ell_i,\bar n)$ and the set of event blocks alone; nothing depends on the matter boundary condition, so each statement holds on the cube exactly as on the torus.

\begin{lemma}[the supply of depth, and the criterion]\label{lem:a01:X0}
(a) $\dist(\Lambda^{(i-1)}_{m+s},(\Lambda^{(i-1)}_m)^c)>s\,r(\ell_{i-1})-(s-1)$ for $s\ge1$ (scale-$(i-1)$ units), and in scale-$i$ units $\dist(B^i(\Lambda^{(i-1)\prime}_{m+n}),B^i(\Lambda^{(i-1)\prime}_m)^c)\ge n\,r(\ell_{i-1})/L-2$ for $n\le L+1$. The hypothesis of \cite[Prop.~1]{Ba3} asks depth $>2r_i$ of $\supp\tilde A$ in $\Omega$; at its use in \cite{Ba2}, where the domain is $B^i(\Lambda_0^{(i)})$, the regions supply $\dist(\supp\tilde A,\Omega^c)>2r_i-2-M_\theta-2\eta$ (the slice of $\theta_i$ and the fine blocks), which is all that the proof consumes. (b) Let $\alpha,\beta>0$, $c\ge0$. If $\alpha L\le\beta$, then $\beta\rho_i-\alpha r_i\ge\alpha R\varrho\ln L\,u_i^{\varrho-1}-\beta s_b/L$ for every $i$, so $\beta\rho_i\ge\alpha r_i+c$ at every scale once $R\ge(c+\beta s_b/L)/(\alpha\varrho\ln L)$. If $\alpha L>\beta$, then $\beta\rho_i<\alpha r_i$ for every $\ell_i<\exp\big(1-\ln L/((\alpha L/\beta)^{1/\varrho}-1)\big)$, whatever $R$ and $c$ are. If an object of size $\theta r_{i+1}$ at scale $i+1$ must sit at depth $\beta r_i$, the containment holds at every scale iff $L\theta\le\beta$, and for $L\theta>\beta$ it fails for every $\ell_i<\exp\big(1-\ln L/(1-(\beta/(L\theta))^{1/\varrho})\big)$ (now $r_{i+1}/r_i=(1-\ln L/u_i)^\varrho$ enters). (c) Where only exponential smallness across scales is used, $e^{-\delta\beta\rho_i}\le Ce^{-(\delta\beta/L)r_i}$, which is $O(\ell_i^{\kappa})$ for every $\kappa$ since $\varrho>1$.
\end{lemma}
\begin{proof}
(a) By \cite[(2.8)]{Ba2} $(\Lambda_{m+1})^c$ contains every large block within $r$ of $(\Lambda_m)^c$, so a point of $\Lambda_{m+1}$ is at distance $>r$ from $(\Lambda_m)^c$; along a lattice geodesic from $x\in\Lambda_{m+s}$ to $(\Lambda_m)^c$ each further layer adds $>r-1$. For the primed sets, points of $B^i(y)$, $B^i(y')$ with $y\in\Lambda_{m+n}$, $y'\in\Lambda_m^c$ differ coordinatewise by $(y-y')+\tau$, $|\tau|<L$, so $|x-x'|_\infty>n\,r(\ell_{i-1})-(n-1)-L$, and dividing by $L$ gives the claim for $n\le L+1$. (b) If $\beta/L\ge\alpha$ then $\beta\rho_i-\alpha r_i\ge\alpha(r_{i-1}-r_i)-\beta s_b/L$ and $r_{i-1}-r_i=R[(u_i+\ln L)^\varrho-u_i^\varrho]\ge R\varrho\ln L\,u_i^{\varrho-1}$. If $\alpha L>\beta$, $\beta\rho_i\le(\beta/L)r_{i-1}$ and $r_{i-1}/r_i=(1+\ln L/u_i)^\varrho<\alpha L/\beta$ exactly when $u_i>\ln L/((\alpha L/\beta)^{1/\varrho}-1)$. (c) $e^{-\delta\beta\rho_i}\le e^{\delta\beta s_b/L}e^{-(\delta\beta/L)r_{i-1}}$ and $e^{-cR(1+\ln\ell^{-1})^\varrho}\le\ell^\kappa$ for small $\ell$ because $\varrho>1$.
\end{proof}

Every place where $L$ enters \cite{Ba1,Ba2,Ba3} and \cite[Thm.~3.1]{K1} is a constant, a rate, a convergent sum, a condition on $(b_0,R,M,\eps_0^{\mathrm{Ba}})$, a divisibility, or a \emph{cross-scale containment} $\beta\rho_i\ge\alpha r_i+c$, and by \cref{lem:a01:X0}(b) a containment holds at every scale iff $\alpha L\le\beta$. At $L=5$ exactly four printed steps fail this test (\cref{thm:a06:BALA}): the cube step of \cite[p.~433]{Ba3} ($\alpha=\frac32$, $\beta=5$), the regularity of the vector background \cite[Lemma~2.3]{Ba2} ($\alpha=1$, $\beta=3$), the boxes of \cite[Lemma~2.4]{Ba2} ($\alpha=4$, $\beta=5$; this one fails at every $L\ge2$), and the split \cite[(2.89)]{Ba2} ($L\theta=L$, $\beta=2$); the partition \cite[(3.43)--(3.45)]{Ba2} is a divisibility entry which fails on some lattices at $L=5$. These are gaps of the printed proofs, not counterexamples to the theorems, whose statements never mention the objects concerned.

\subparagraph{(a) The cube step.} On \cite[p.~433]{Ba3} the localisation cubes of a renormalised class, which lie in a cube of side $\ell_{\bar n}$ depending on $\bar n$ only \cite[p.~432]{Ba3}, are enclosed in concentric cubes $\square_1\subset\square$ of sides $r_i$, $3r_i$, sums of large blocks; $G_i(\Omega,\tilde B)$ is written $G_i(\square,\tilde B)+\delta G_i(\Omega,\square,\tilde B)$ with the difference treated as an external field, the regular field $\tilde B$ is replaced on $\square$ by its constant part $B_0$, $B_0$ is gauged away, and $G_i(\square,0)$ is replaced by $G_i(0)$. The step presupposes $\square\subset\Omega$ (resp.\ $\square\subset\Omega_2$): the Neumann operator $G_i(\square,\tilde B)$ is built from the bonds of $\square$, where $\tilde B$ lives only if $\square\subset\Omega$, and \cite[(2.26)]{Ba1} is stated for nested domains. Every vertex of every class lies within sup-distance $1$ of $\Omega_1\cup\supp\tilde A$ (the vertices \cite[(1.6)--(1.7)]{Ba3} are summed over $\Omega_1$; those of \cite[(1.8)--(1.11)]{Ba3} carry $g_i$, $\supp g_i\subset\Omega_1$, or $\tilde A$; the pair vertices \cite[(1.12)--(1.15)]{Ba3} carry $A'$ or $\tilde A$ along a contour inside one block). For an old vertex $x_v$ of the class the containment reads
\begin{equation}\label{eq:a01:N}
\dist(x_v,\Omega^c)>\tfrac32r_\square+2s_b+\ell_{\bar n}\ \Longrightarrow\ \square\subset\Omega\ \text{and}\ \dist(\square_1,\square^c)\ge r_\square
\end{equation}
(cubes rounded outward to large blocks; triangle inequality). With the printed $r_\square=r_i$ and a vertex anchored at $\Omega_1$, whose depth is $\ge5r(\ell_{i-1})/L-2$ by \cref{lem:a01:X0}(a), this is a containment with $\alpha=\frac32$, $\beta=5$: it holds at every scale iff $L\le3$, and at $L=5$ it fails for all small $\ell_i$; the failing vertices ($\lambda$- and $\dm$-vertices summed over $\Omega_1$, $g_i$-localised photon legs near $\partial\Omega_1$) carry no $\tilde A$-leg, so no hypothesis on $\supp\tilde A$ helps.

\begin{proposition}[the cube step with a constant cube]\label{prop:a01:P12}
Let $\ell^\sharp\defeq3\ell_{\bar n}+4N_B$, where $N_B=n_1+n_2-1$ is the number of insertions per line in the expansion of \cite[p.~433]{Ba3} to order $\bar n$ (a function of $\bar n$ only), let $r^\flat$ be $\max\{3\ell^\sharp+R_{\mathrm{Ba}}+12,\,3R_{\mathrm{Ba}}+15,\,107\}$ rounded up to a multiple of $s_b$, and $D^\flat\defeq\frac32r^\flat+2s_b+\ell^\sharp$. Run the step with concentric cubes $\square_1,\square$ of sides $r^\flat$, $3r^\flat$ centred at the component cube. If every old vertex of the class has depth $>D^\flat$ in $\Omega$ (resp.\ $\Omega_2$), the proof of \cite[pp.~433--445]{Ba3} goes through and gives \cite[(1.33)]{Ba3} with a constant of the printed form, depending on $(\alpha_0,\bar n,d,a,M)$ only. If $R\ge R^\flat\defeq D^\flat+3$ and $R\ge M_\theta+2$, then at $L=5$ every class occurring in the use of \cite[Prop.~1]{Ba3} in \cite{Ba1,Ba2} has its old vertices at depth $>D^\flat$, so \cite[Prop.~1]{Ba3} holds at $L=5$ in its printed form.
\end{proposition}
\begin{proof}
The side $r_\square$ enters \cite[pp.~433--434]{Ba3} in five places. (1) \emph{Containment:} \eqref{eq:a01:N} with $r_\square=r^\flat$. (2) \emph{Old-vertex classes lie in $\square_1$:} the expansion in $\tilde B'$ uses \cite[(3.14), (3.44), (3.45)]{Ba1} to order $\bar n$; only lines between localisations whose boxes meet stay internal \cite[p.~432]{Ba3}, so along an expanded line a new vertex advances at most $2$, and the new vertices lie within $2N_B$ of the class-component cube; the boxes of all graphs of a connected component lie in a grid cube of side $3\ell_{\bar n}$ concentric with one member's cube (``localizations \dots\ in agreement'', \cite[p.~432]{Ba3}); so the class lies in a cube of side $\ell^\sharp$, and $3\ell^\sharp\le r^\flat$. (3) \emph{The two $\delta G$ external fields are $O(1)$:} $\delta G_i(\Omega,\square,\tilde B)$ is evaluated at points of $\square_1$, at distance $\ge r^\flat\ge R_{\mathrm{Ba}}+12$ from $\square^c$, with sources in $\square_1$, so \cite[(2.26)]{Ba1} for the pair $(\square,\Omega)$ gives $\norm{h\,\delta G_i\,h'}_{1,\alpha}\le c_{\mathrm{Ba}}e^{-\delta_0r^\flat}$; for $\delta G_i(\square,\eta\Z^d,0)$ take a torus $T'\supset\square$ with side a multiple of $2s_b$ and at least $\max\{3s_b,3r^\flat+2s_b+2\}$: at zero background $G_i(T',0)$ is the periodisation of $G_i(\eta\Z^d,0)$, so $\delta G_i(\square,\eta\Z^d,0)=\delta G_i(\square,T',0)+\sum_{n\ne0}G_i(\eta\Z^d,0)(\cdot,\cdot+n\ell')$, the first term obeying \cite[(2.26)]{Ba1} on $T'$ with constants independent of $\ell'$ (the constants of \cite{Ba4} depend on $d$, $a$, $M$, $L$ only) and the second vanishing as $\ell'\to\infty$ by \cite[(2.25)]{Ba1}. Both are external fields of dimension $0$ \cite[(2.1), p.~422]{Ba3}, so an $O(1)$ bound is all that is used. (4) \emph{The constant field:} for $x\in\square$ and $c$ the centre, $|\tilde B(x)-\tilde B(c)|\le d(\frac32r^\flat+s_b)\sup_\square|\partial^\eta\tilde B|=O(r^\flat p_i)=O(p_i)$ with $B_0\defeq\tilde B(c)$, in place of the printed $O(r_ip_i)\le O(p_i^2)$; every use of this bound is an upper bound monotone in the size of $\tilde B'$ (vertices weighted by $e(\ell_i)|\tilde B'|$, the remainder \cite[(1.16)]{Ba3} by $(e(\ell_i)^{1-\alpha})^{n+n'}$, the norm comparison between the norms \cite[(1.32)]{Ba3} at $B_0$ and at $\tilde B$ costing $(1+O(p_i))^{\#\mathrm{legs}}$, absorbed by the slack $\ell_i^{(4-d)/4-\beta}$ of the applications \cite[p.~421]{Ba3}); so a smaller cube only helps. (5) \emph{Gauge-out of $B_0$:} $B_0=\partial^\eta\lambda_0$ with $\lambda_0(x)=B_0\cdot(x-c)$ on $\square$; the Neumann covariant Laplacian on the bonds of $\square$ and the averaging $Q_i$ (contours inside blocks contained in $\square$) transform covariantly, $G_i(\square,B_0)=U(\lambda_0)G_i(\square,0)U(\lambda_0)^*$, and every class is invariant when the external scalar fields are transformed simultaneously \cite[(2.69)--(2.74)]{Ba2}; exact as long as $\square$ does not wrap around the torus (no condition on the cube). (6) \emph{$\Omega$-freeness after the split:} every old vertex lies in $\square_1\subset\square\subset\Omega$ and every new vertex in $\square$, so the characteristic function of $\Omega$ is $1$ wherever a vertex or a line end occurs; every object depending on $\Omega$ or on the cube ($\delta G_i(\Omega,\Omega_2,\tilde B)$, \cite[(1.16)]{Ba3}, the two $\delta G$'s, the gauge factors) has been moved into external fields, used on \cite[pp.~434--445]{Ba3} only through the norms \cite[(3.3)]{Ba3} at points of $\square_1$; the $\Omega_1$-dependent objects $g_i$, $1_{\Omega_1}$, $\dm(e',g_i,\lambda',\Omega_1,x)$ stay in the vertices, fixed before the step and entering only identities between a graph and its counterterm with the same restriction; on \cite[pp.~434--445]{Ba3} every propagator is $G_i(0)$, a slice, $G^\xi$, $C^\xi$ or $\Pi$, every sum runs over $\eta\Z^d$, $\xi\Z^d$, the cubes $\Delta(v)$, $\Omega_1$ or the torus identified with $\eta\Z^d$, Taylor's formula \cite[(3.10)]{Ba3} is applied to legs and to $g,g'$, and the $\lambda$-vertices are matched against counterterms with the same localisation \cite[(3.19), pp.~444--445]{Ba3}: none of this refers to the cube side or to $\dist(\cdot,\Omega^c)$. (7) \emph{Chains without old vertices} \cite[p.~433]{Ba3}: a class with no old vertex is a path $v_1$--$\cdots$--$v_m$ of new vertices of one product $(G_0V)^mG_0$, $G_0\defeq G_i(\square,B_0)$, with the cut lines at its ends as external scalar fields; the printed route through \cite[Prop.~2.1]{Ba3} is not used (its summation by parts \cite[(2.8)]{Ba3} is printed on the torus and on the Neumann box produces a face term where the $\tilde B'$-localisation does not vanish); instead, with $\norm\psi_E^2\defeq\ip\psi{G_0^{-1}\psi}$, $G_0^{-1}\ge D_{B_0}^*D_{B_0}$ and $G_0^{-1}\ge\gamma_\square$ (constant field, gauged out), each vertex of \cite[(1.8)--(1.11), (1.13)--(1.15)]{Ba3} is a bounded form in $\norm\cdot_E$: the form of \cite[(1.8)]{Ba3} is $\kappa\sum_{b\subset\square}\eta^d(D_{B_0}\phi)(b)\cdot M(b)\chi(b_-)$ with $|M(b)|\le C(\sup|\tilde B'|)^{n'}$, bounded through $\norm{D_{B_0}\phi}\le\norm\phi_E$ and $\norm\chi\le\gamma_\square^{-1/2}\norm\chi_E$; \cite[(1.9), (1.11)]{Ba3} likewise, with $|R_{\bar n+1}|\le1$; \cite[(1.10)]{Ba3} carries no derivative and uses $\gamma_\square^{-1/2}$ on both legs; a block pair \cite[(1.13)--(1.14)]{Ba3} uses $|\tilde B'(\Gamma)|\le2d\sup|\tilde B'|$ on its unit-block contours and $\sum_{x\in B(y)}\eta^d|\phi(x)|\le\norm\phi$; every form norm is $\le C(e(\ell_i)\sup|\tilde B'|)^{n'_v}$. Hence $E(H)=\ip{G_0^{1/2}K_1^*\Phi_1}{\hat K_2\cdots\hat K_{m-1}G_0^{1/2}K_m\Phi_m}$, $\hat K_v\defeq G_0^{1/2}K_vG_0^{1/2}$, is bounded by the product of the form norms times the end norms $\norm{h\Phi}_E$, where $h=1$ on the $2\eta$-neighbourhood of the end cube, $|\partial^\eta h|\le c_h$, and $\norm{h\Phi}_E^2=\norm{D_{B_0}(h\Phi)}^2+m^2\ell_i^2\norm{h\Phi}^2+a_i\norm{Q_i(B_0)h\Phi}^2\le C(1+c_h)^2\norm{h\Phi}_{1,0}^2$ by the Leibniz rule on a cube of side $\le4$; so $\norm{1_{\square^\eta(v)}\Phi}+\norm{1_{\square^\eta(v)}D_{B_0}\Phi}\le C\norm{h\Phi}_{1,0}$ is all that is used of an end field. The cut lines are bounded pointwise on $\square$, with their first differences in each slot, at pairs of boxes at distance $\ge1-2\eta\ge\frac35$, by \cref{lem:a01:RBK}: $|D^a_{B_0,x}D^b_{B_0,x'}G_0(x,x')|\le Ce^{-\delta|x-x'|}$ for $|x-x'|\ge\frac35$ and $a,b\in\{0,1\}$, at every pair of points of $\square$, its faces included, with no depth condition; and a cut line $K$ contracted between two functionals $A$, $B$ bounded by the sup-and-difference norms of the two end regions obeys $|A(x\mapsto B(K(x,\cdot)))|\le\norm A\norm B\max_{a,b\in\{0,1\}}\sup|D^a_xD^b_{x'}K(x,x')|$, because $x\mapsto B(K(x,\cdot))$ is linear in $K$ and its differences are $B(D_xK(x,\cdot))$ bond by bond, so no hypothesis $\dist\ge1$ on the two supports is needed and distance exactly $1$ is covered. A cut line between a chain end and an old-vertex class is an external scalar field of that class, read through its norms \cite[(3.3)]{Ba3} after localisation to the $\eta$-neighbourhood of the class's cube by a function $h'$ as in \cref{lem:a01:RBK}(c), the chain slot anywhere in $\square$: by the same linearity, with $B$ the class functional, bounded by Ba{\l}aban's norm \cite[(1.32)]{Ba3} of its external field, $|A(x\mapsto B(h'K(x,\cdot)))|\le\norm A\norm B\max_{a\in\{0,1\}}\sup_x\norm{h'D^a_xK(x,\cdot)}$ with the norm \cite[(1.32)]{Ba3} in the last factor, and \cref{lem:a01:RBK}(c) bounds that factor by $C(1+c_h)^2$ at every $x$ of the chain-end region, the faces of $\square$ included. The vertex forms are sums over the bonds of $\square$, so vertices at the faces need no separate treatment. (8) \emph{The small factors of \cite{Ba2} survive:} the factors $O(e^{-\delta_1r_i})$ that \cite[p.~422]{Ba3} produces at domain changes come from $\exp[-\delta_0d(\{\square(v)\})]$ in \cite[(1.33)]{Ba3}, i.e.\ from the lines removed before the cube step, whose bound depends on $\Omega$ only; the two-domain inputs enter as external fields, linearly. So the only place where $r_\square=r_i$ was used is (1), where $r^\flat$ serves by hypothesis; the constant acquires the factors of (3) and (4), both of the printed kind. For the last claim: a vertex anchored at $\Omega_1$ has depth $\ge5r(\ell_{i-1})/5-3\ge r_i-3\ge R-3\ge D^\flat$ (\cref{lem:a01:X0}(a), $r(\ell_{i-1})\ge r_i\ge R$); a vertex anchored at $\supp\tilde A$ has depth $>2r_i-2-M_\theta-2\eta\ge2R-2.4-M_\theta>D^\flat$ (\cref{lem:a01:X0}(a)); the other steps of the proof need depth at most $R_{\mathrm{Ba}}+13<D^\flat$ at old vertices or within $2\sqrt3+1$ of $\supp\tilde A$.
\end{proof}
% a01_H0rel part 6x1, continued: two-slot bounds of the box propagator at a constant field (input from a01_p6x1.tex after prop:a01:P12)
\begin{lemma}[two-slot bounds of the box propagator at a constant field]\label{lem:a01:RBK}
Let $i\ge1$, $\eta\defeq L^{-i}$ the fine spacing in scale-$i$ units, $\square$ a box of large blocks (a product of intervals, each a union of large blocks) that does not wrap around the torus, with Neumann conditions on its faces, $B_0$ a constant field, and $G_0\defeq G_i(\square,B_0)$ the propagator of step (7) of the proof of \cref{prop:a01:P12}, $G_0^{-1}=D_{B_0}^*D_{B_0}+m^2\ell_i^2+a_iQ_i(B_0)^*Q_i(B_0)$ on $\ell^2(\square\cap\eta\Z^d)$ with the uniform block average $Q_i$. Kernels are taken for the measure $\eta^d$, and $D^a_{B_0,x}$, $a\in\{0,1\}$, denotes the identity or a covariant forward difference on a bond of $\square$ at $x$. \textup{(a)} There are $C,\delta>0$ depending on $(d,L,a)$ and the mass range only, free of $i$, $\eta$, $\square$, $B_0$ and the positions, such that for all $x,x'\in\square$ and $a,b\in\{0,1\}$ ($d\ge3$; here $d=3$)
\[
|D^a_{B_0,x}D^b_{B_0,x'}G_0(x,x')|\ \le\ C\,(\eta+|x-x'|)^{-(d-2+a+b)}\,e^{-\delta|x-x'|},
\qquad\text{so}\qquad |D^a_{B_0,x}D^b_{B_0,x'}G_0(x,x')|\le5C\,e^{-\delta|x-x'|}\ \text{ for }|x-x'|\ge\tfrac35\ (d=3).
\]
\textup{(b)} No depth condition enters: the bound holds at every pair of points of $\square$, the bonds of its faces included.

\textup{(c)} \emph{One slot in Ba{\l}aban's norm.} Let $\Delta,\Delta'$ be two cubes of $\square$ at distance $\ge1$, let $h'$ be a localisation function on $\square$, $0\le h'\le1$, equal to $1$ on the $\eta$-neighbourhood of $\Delta'$ and vanishing at distance $\ge\frac12$ from $\Delta'$, with $|\partial^\eta h'|\le c_h$ and $|\partial^\eta\partial^\eta h'|\le c_h$, and let $x$ lie within $\eta$ of $\Delta$. Then for $a\in\{0,1\}$ the field $\Psi\defeq h'(\cdot)\,D^a_{B_0,x}G_0(x,\cdot)$ satisfies the three clauses of Ba{\l}aban's external-field norm \cite[(1.32)]{Ba3} (sup, first differences, $\alpha_0$-H\"older quotients of the first differences on unit balls) on all of $\square$, its faces included, with the bound $C(1+c_h)^2$, $C$ of the same kind as above.
\end{lemma}
\begin{proof}
(1) \emph{Gauge-out.} With $\lambda_0(x)\defeq B_0\cdot(x-c)$ on the non-wrapping $\square$, $G_i(\square,B_0)=U(\lambda_0)G_i(\square,0)U(\lambda_0)^*$ and $D_{B_0}(U(\lambda_0)\phi)=U(\lambda_0)D\phi$ bond by bond (step (5) of the proof of \cref{prop:a01:P12}; $|U|=1$), so every modulus in the claim equals the one at $B_0=0$. Let $B_0=0$.

(2) \emph{Even images.} Let $\square=\prod_\mu[\alpha_\mu,\beta_\mu)$ and let $T_\square$ be the torus of periods $2(\beta_\mu-\alpha_\mu)$, each a multiple of $2s_b$, so that the large blocks of $\square$ tile $T_\square$ periodically; let $\mathcal W_\square\cong(\Z_2)^d$ be the group generated by the reflections in the face planes $x_\mu=\alpha_\mu-\eta/2$ of $\square$ (planes between sites on block boundaries; on $T_\square$ the reflection in $x_\mu=\beta_\mu-\eta/2$ is the same element). They map sites to sites, bonds to bonds and blocks to blocks, the uniform average commutes with them, and the mass term is constant; every bond of $T_\square$ that crosses a face plane joins a site to its mirror image, where a $\mathcal W_\square$-even field has zero difference, and the remaining bonds of $T_\square$ inside $\square$ are the bonds of $\square$. Hence, exactly as in \cref{lem:a01:PE}(a) and \cref{lem:a01:NL}, the Neumann operator on $\square$ is the even restriction of the periodic one, $G_i(T_\square,0)\Ebar=\Ebar\,G_i(\square,0)$ with $\Ebar$ the even extension, and
\[
G_i(\square,0)(x,x')=\sum_{g\in\mathcal W_\square}G_i(T_\square,0)(x,gx')\qquad(x,x'\in\square),
\]
a difference in $x'$ along a bond of $\square$ acting on the image terms along the reflected bonds. Every image satisfies $\dist_{T_\square}(x,gx')\ge|x-x'|$: coordinatewise, for $x_\mu,x'_\mu\in[\alpha_\mu,\beta_\mu-\eta]$ and $r$ the reflection in $\beta_\mu-\eta/2$, $rx'_\mu-x_\mu=|x_\mu-x'_\mu|+2(\beta_\mu-\tfrac\eta2-\max(x_\mu,x'_\mu))\ge|x_\mu-x'_\mu|$, and the complementary distance around the period is $x_\mu+x'_\mu-2\alpha_\mu+\eta\ge|x_\mu-x'_\mu|$.

(3) \emph{Periodisation.} At zero field $G_i(T_\square,0)(x,y)=\sum_{n\in\Z^d}G_i(\eta\Z^d,0)(x,y+2\ell n)$, $\ell\defeq(\beta_\mu-\alpha_\mu)_\mu$ (the periodic operator is the periodisation of the free-lattice one, as in step (3) of the proof of \cref{prop:a01:P12}), every term having $|x-y-2\ell n|\ge\dist_{T_\square}(x,y)$, and at most $(2+k)^d$ values of $n$ have $|x-y-2\ell n|<k+1$, every period being $\ge1$. So, with the free-lattice bound (4) below at the rate $\delta'$ and $\rho\defeq|x-x'|$,
\[
|D^a_xD^b_{x'}G_i(\square,0)(x,x')|\ \le\ 2^dC'(\eta+\rho)^{-(d-2+a+b)}\sum_{k\ge\lfloor\rho\rfloor}(2+k)^de^{-\delta'k}\ \le\ C(\eta+\rho)^{-(d-2+a+b)}e^{-\delta'\rho/2}.
\]

(4) \emph{The free lattice.} On $\eta\Z^d$ at zero field, $G_i(\eta\Z^d,0)=G_1+\sum_{j=1}^{i-1}G_{(j)}$ with the slices $G_{(j)}\defeq G_{j+1}-G_j=a_jG_jQ^*C^{(j)}QG_j$ (the Schur identity behind \cref{lem:a01:legs}(b) and \cite[(2.17)]{K1}, with $C^{(j)}$ as in \cref{lem:a01:A23}(c)), so that the sum telescopes to $G_i$. In unit-$j$ units ($j$-blocks of side one, spacing $L^{-j}$) let $\ell_j(\cdot,u)\defeq a_jG_jQ^*e_u$ be the leg of the unit datum on the $j$-block $u$; since $G_j$ is real symmetric at zero field, $G_{(j)}(y,y')=a_j^{-1}\sum_{u,v}\ell_j(y,u)C^{(j)}(u,v)\ell_j(y',v)$, each variable sitting in its own leg, so
\[
D^a_yD^b_{y'}G_{(j)}(y,y')=a_j^{-1}\sum_{u,v}(D^a\ell_j)(y,u)\,C^{(j)}(u,v)\,(D^b\ell_j)(y',v).
\]
The legs obey $|D^a\ell_j(y,u)|\le Ce^{-\delta_2\dist(y,B(u))}$ at every site and every bond of the free lattice, $a\in\{0,1\}$ (the leg bounds \cite[(3.64)]{K1}, \cite[(2.12)]{Ba3}, at the uniform average $f=1$ in \cref{prop:a02:FL}; on the free lattice every bond is interior, so they are also \cref{lem:a01:A4}(b) and the interior case (4) of the proof of \cref{lem:a01:A5}), and $|C^{(j)}(u,v)|\le Ce^{-\delta_1|u-v|}$ (\cite[Thm.~3.3, (3.6)]{K1}; \cref{lem:a01:A23}(c)). With $|y-y'|\le\dist(y,B(u))+|u-v|+\dist(y',B(v))+\sqrt d$ and two geometric sums over $u,v\in\Z^d$, $|D^a_yD^b_{y'}G_{(j)}(y,y')|\le C_Se^{-\delta_S|y-y'|}$ at every pair of points, $\delta_S\defeq\frac12\min\{\delta_1,\delta_2\}$. In scale-$i$ units the slice $j$ has the kernel $s_j^{2-d}G_{(j)}(x/s_j,x'/s_j)$ with $s_j\defeq L^{j-i}$, and a scale-$i$ difference is $s_j^{-1}$ times a unit-$j$ one, so slice $j$ contributes at most $C_Ss_j^{-(d-2+a+b)}e^{-\delta_S|x-x'|/s_j}$; $G_1$, read in unit-$1$ units at the fixed spacing $L^{-1}$, obeys the same bound with $s\defeq L\eta$ and a constant depending on $L$ (its kernel is $L^d\ip{e_y}{G_1e_{y'}}$, bounded with the decay of \cref{lem:a01:A23}(a), and a difference at spacing $L^{-1}$ costs a factor $2L$). Finally, for $p\defeq d-2+a+b>0$ and $s$ ranging over the powers of $L$ in $[L\eta,L^{-1}]$: $e^{-\delta_S\rho/s}\le e^{-\delta_S\rho/2}e^{-\delta_S\rho/(2s)}$ as $s\le1$; the terms with $s\ge\max(\rho,\eta)$ sum to at most $(1-L^{-p})^{-1}\max(\rho,\eta)^{-p}$, and the terms with $s<\rho$ are $\rho^{-p}(\rho/s)^pe^{-\delta_S\rho/(2s)}$ with $\rho/s$ running over powers of $L$, summable to at most $C_p\rho^{-p}$. Hence $|D^a_xD^b_{x'}G_i(\eta\Z^d,0)(x,x')|\le C'(\eta+\rho)^{-p}e^{-\delta'\rho}$ with $\delta'\defeq\delta_S/2$, which is (4), and (3) gives the claim with $\delta\defeq\delta'/2$. The last clause of the statement is $(\eta+\rho)^{-p}\le(5/3)^3<5$ for $\rho\ge\frac35$.

(5) \emph{The clause (c).} By (1) every modulus is the one at $B_0=0$. The reflections of (2) and the translations of (3) are isometries of the lattice mapping bonds to bonds, so each clause of $x'\mapsto D^a_xG_i(\square,0)(x,x')$ at a pair $x',y'\in\square$ is bounded by the sum over the images of the same clause of $x'\mapsto D^a_xG_i(\eta\Z^d,0)(x,x')$ at the pair $(gx'+2\ell n,gy'+2\ell n)$, whose points are at least as far from $x$ as $x'$, $y'$ are, with the count of (3). Write $\Psi_0(x')\defeq D^a_xG_i(\eta\Z^d,0)(x,x')$ and $\rho'\defeq\min\{|x'-x|,|y'-x|\}$. \emph{Claim:} on the free lattice at zero field the sup and first-difference clauses of $\Psi_0$ at $x'$, and the H\"older quotient at a pair $x',y'$ with $|x'-y'|\le1$, are $\le Ce^{-\delta\rho'}$ whenever $\rho'\ge\frac1{20}$. Indeed, by the factorisation of (4), in unit-$j$ units $D^a_xG_{(j)}(x,\cdot)=\sum_v\ell_j(\cdot,v)\,a_j^{-1}\beta_j(v)$ is the leg of the block data $a_j^{-1}\beta_j$, $\beta_j(v)\defeq\sum_u(D^a\ell_j)(x,u)C^{(j)}(u,v)$, with $|\beta_j(v)|\le Ce^{-\delta_S\dist(x,B(v))}$ by the leg and $C^{(j)}$ bounds of (4) and one geometric sum. Split the data into the blocks with $\dist(x,B(v))\le\rho'_j/2$, $\rho'_j\defeq\rho'/s_j$ the distance in unit-$j$ units, and the rest. The leg of bounded block data obeys the sup bound of \cref{lem:a01:A4}(b), the first-difference bound of \cref{lem:a01:A5} and its H\"older bound on balls of radius $\frac1{16}$ with decay $e^{-\delta_3\dist(\cdot,S)}$ from the support $S$ of the data (the interior case, every bond of the free lattice being interior; printed: \cite[(3.63)--(3.65)]{K1}, \cite[(2.12)]{Ba3}); the far part has data of size $\le Ce^{-\delta_S\rho'_j/2}$, and the near part has its support within $\rho'_j/2$ of $x$, so at distance $\ge\rho'_j/2$ from $x'$ and $y'$; a pair at unit-$j$ distance $>\frac1{16}$ is bounded by $2\cdot16^{\alpha_0}$ times the first differences. So in unit-$j$ units every clause of slice $j$ is $\le Ce^{-\delta_6\rho'_j}$, $\delta_6\defeq\frac14\min\{\delta_S,\delta_3\}$. In scale-$i$ units the sup clause of slice $j$ scales as $s_j^{-(d-2+a)}$, the first differences as $s_j^{-(d-1+a)}$ and the H\"older quotients as $s_j^{-(d-1+a+\alpha_0)}$ (a scale-$i$ pair at distance $t\le1$ is a unit-$j$ pair at distance $t/s_j$; for $t>s_j/16$ use two first differences, $2\cdot16^{\alpha_0}s_j^{-(d-1+a)}t^{-\alpha_0}\le2\cdot16^{\alpha_0}s_j^{-(d-1+a+\alpha_0)}$); the slice $G_1$ at the fixed spacing contributes a bounded multiple of its first differences to the H\"older clause, every pair being at distance $\ge\eta$. The sum over $s$ is the geometric sum of (4) with $d-1+a+\alpha_0$ in place of $p$, giving $\le C(\eta+\rho')^{-(d-1+a+\alpha_0)}e^{-\delta'\rho'}\le Ce^{-\delta'\rho'}$ for $\rho'\ge\frac1{20}$; this is the claim, and with the image count it holds for $x'\mapsto D^a_xG_i(\square,0)(x,x')$ at pairs of $\square$ with $\rho'\ge\frac1{20}$. \emph{Localisation.} $\Psi=h'\Psi_\square$, $\Psi_\square\defeq D^a_xG_i(\square,0)(x,\cdot)$, and $D(h'\Psi_\square)=h'(\cdot+\eta e)D\Psi_\square+(\partial^\eta h')\Psi_\square$ bond by bond. Both $\Psi$ and $D\Psi$ vanish outside $S'\defeq\{\dist(\cdot,\Delta')\le\frac12+\eta\}$, whose points are at distance $\ge\frac12-2\eta\ge\frac1{10}$ from $x$ ($\eta\le\frac15$, $x$ within $\eta$ of $\Delta$, $\dist(\Delta,\Delta')\ge1$). Hence the sup and first-difference clauses of $\Psi$ are those of $\Psi_\square$ on $S'$ times $1+c_h$. For a H\"older pair $x',y'$ with both points at distance $\ge\frac1{20}$ from $x$, the Leibniz rule and the bounds $|h'(x')-h'(y')|\le dc_h|x'-y'|$, $|\partial^\eta h'(x')-\partial^\eta h'(y')|\le dc_h|x'-y'|$ give the quotient $\le C(1+c_h)^2$ from the three clauses of $\Psi_\square$; a pair with one point $y'$ within $\frac1{20}$ of $x$ has $D\Psi(y')=0$ and $|x'-y'|\ge\frac1{10}-\frac1{20}=\frac1{20}$ whenever $D\Psi(x')\ne0$, so its quotient is $\le20^{\alpha_0}$ times a first difference. This is (c).
\end{proof}

On the cube the same device, a constant box $Q_x$ of large blocks around each vertex cube at depth $\ge R_*$, is step (1) of the proof of \cref{thm:a01:B2}; the depth supply is \cref{lem:a01:X0}(a) there too.

% a01_H0rel part 6x2: the replaced steps (b)-(d) at block factor 5 (input from a01_H0rel_p6.tex after a01_p6x1)
\subparagraph{(b) The cut-off $\zeta^{(i)}$.} In \cite[(2.44)]{Ba2} $A^{(i)}=a_i\ell_i^{-2}\zeta^{(i)}G_iQ_i^*A$ with $\supp\zeta^{(i)}(\cdot,y)\subset\{|x-y|<r_i-2M\}$, $\zeta^{(i)}(x,y)=1$ for $|x-y|\le\frac12r_i$. The printed proof of \cite[(2.60)]{Ba2}, $\partial^\eta A^{(i)}=O(p_i)$ on $\Omega$, bounds the first two terms of \cite[(2.61)]{Ba2} by the restrictions \cite[(2.55)]{Ba2} at every coarse point $y'$ with $\zeta^{(i)}(x,y')\ne0$; \cite[(2.55)]{Ba2} is available on $\Lambda_{-1}^{(i-1)\prime}$, which contains $\Omega$ at depth $\ge3r(\ell_{i-1})/L-2$ (\cref{lem:a01:X0}(a), three layers). The containment $3r(\ell_{i-1})/L-2\ge r_i-2M$ has $\alpha=1$, $\beta=3$: it holds at every scale iff $L\le3$. The gap is real: by \cite[(2.9)]{Ba2} the restrictions survive the partial resummation only on $\Lambda'_{-1}$, $\Lambda^*_{-1}$, $\Lambda_{-1}$, so at a large block which stayed large at every finer scale $A(y')$ is unrestricted and the contribution $a_i\zeta^{(i)}(x,y')(G_iQ_i^*)(x,y')A(y')$ has no bound; hence \cite[(2.60)]{Ba2} as printed is not established in a collar of width $\approx r_i-3r(\ell_{i-1})/L$ along $\partial\Omega$, while \cite[p.~571]{Ba2} uses it on all of $\Omega$ and \cite[Prop.~1]{Ba3} assumes it on \cite[p.~412]{Ba3}. The same cut-off enters the split \cite[(2.89)]{Ba2} at the next scale.

\begin{proposition}[the shrunk cut-off]\label{prop:a01:Zth}
Replace $\zeta^{(i)}$ in \cite[(2.44)]{Ba2}, at every scale, by $\zeta^{(i)}_\theta$ with $\supp\zeta^{(i)}_\theta(\cdot,y)\subset\{|x-y|<\theta r_i-2M\}$, $\zeta^{(i)}_\theta(x,y)=1$ for $|x-y|\le\frac12\theta r_i$, $|\partial^\eta\zeta_\theta|\le1$, $|\partial^\eta\partial^\eta\zeta_\theta|\le2$ (a $C^{1,1}$ radial profile), with $\theta=\theta_5\defeq\frac25$; the first-step cut-off of \cite[p.~561]{Ba2} already has range $r/(2L)$ in the units of the next scale, so this returns to the first-step convention. Let $R\ge\max\{10(M+1),\,M_\theta+2,\,(M_\theta+7)/11.6\}$, where $\theta_{i+1}$ equals $1$ on $B^i(\Lambda_2^{(i)})$ and $0$ at distance $\ge M_\theta$ from it ($M_\theta<M$ by \cite[p.~567]{Ba2}, where $\theta_{i+1}$ ``varies smoothly'' on a slice of thickness $<M$; we take $M_\theta=2M/3$ and a $C^{1,1}$ profile). Then at $L=5$:
\begin{enumerate}[label=\textup{(Z\arabic*)},nosep,leftmargin=*]
\item for $y\in\Lambda_2^{(i-1)\prime}$, $x\in B^i(y)$, the $\zeta_\theta$-ball $\{|x-y'|<\theta r_i-2M\}$ lies in $\Lambda_{-1}^{(i-1)\prime}$, where \cite[(2.55)]{Ba2} holds;
\item \cite[(2.59)--(2.60), (2.64)]{Ba2} hold with $O(p_i)$ replaced by $O_L(p_i)$, the constant containing $p_{i-1}/p_i\le(1+\ln L)^p$;
\item the terms carrying $1-\zeta^{(i)}_\theta$ in \cite[(2.61), (2.64)]{Ba2}, the terms carrying $1-\zeta^{(i)}_\theta$, $1-\zeta^{(i+1)}_\theta$ in \cite[(2.88)--(2.90)]{Ba2}, and the second term of \cite[(2.89)]{Ba2} (the one with $\delta C^{(i)}_{\Lambda_0}$) are $O(\ell_i^\kappa)$ for every $\kappa$, with their derivatives and the $\alpha_0$-H\"older quotients of their derivatives;
\item the transition width $\frac12\theta r_i-2M\ge2$ required by the $C^{1,1}$ radial profile with $|\partial^\eta\zeta_\theta|\le1$ and $|\partial^\eta\partial^\eta\zeta_\theta|\le2$ holds;
\item the split \cite[(2.89)]{Ba2} holds \emph{as printed}, with the full $B$ in its fifth term, and the $\zeta^{(l)}_\theta$-balls around $\supp\theta_l$ lie in $\Lambda_{-1}^{(l-1)\prime}$ at every scale $l$.
\end{enumerate}
No $\zeta$ obeying the printed \cite[(2.44)]{Ba2} makes \cite[(2.89)]{Ba2} exact at $L=5$ for small $\ell_i$ (the minimal printed support $\frac12r$ gives $L\theta=\frac52>2$), so the shrunk cut-off is necessary at $L=5$ if \cite[(2.89)]{Ba2} is kept in Ba{\l}aban's form.
\end{proposition}
\begin{proof}
(Z1): $\theta r_i\le3r(\ell_{i-1})/L-2+2M$ is \cref{lem:a01:X0}(b) with $\alpha=\frac25$, $\beta=3$, $\alpha L=2\le3$, margin $\ge r(\ell_{i-1})/5-2\ge0$ for $R\ge10$. (Z2): in \cite[(2.61)]{Ba2} with $\zeta_\theta$ the first term sums over $y'$ in the $\zeta_\theta$-ball, where by (Z1) $|A(y')-A(y)|\le c_1p_{i-1}L|y'-y|$ (a path of $\le L|y'-y|$ bonds with $|\partial A|\le c_1p_{i-1}$) and \cite[Prop.~2.2]{Ba2} gives the kernel bound $c_{\mathrm{Ba}}e^{-\delta_0|x-y'|}$, so the term is $O(p_{i-1})=O_L(p_i)$; the third term is the exact \cite[(2.62)--(2.63)]{Ba2} ($G_iQ_i^*1$ is constant); the second is a tail of (Z3); \cite[(2.64)]{Ba2} is the same for the derivative, with $\partial^\eta(G_iQ_i^*1)=0$. (Z3): the $(1-\zeta_\theta)$ term of \cite[(2.61)]{Ba2} multiplies the single constant $A(y)$, $|A(y)|\le c_1p_{i-1}/(\lmu\ell_{i-1})$, with a kernel supported at $|x-y'|\ge\frac12\theta r_i$ and bounded by $c_{\mathrm{Ba}}e^{-\delta_0|x-y'|}$, so it is $\le Cr_i^de^{-\delta_0\theta r_i/2}p_{i-1}/(\lmu\ell_{i-1})\le\ell_i^\kappa$ for small $\ell_i$ (\cref{lem:a01:X0}(c)); in \cite[(2.89)]{Ba2} the third and fourth terms carry $1-\zeta^{(i)}_\theta$, $1-\zeta^{(i+1)}_\theta$ acting on $\Lambda_0^{(i)\prime}B$ with $|B|\le c_1p/(\lmu\ell)$ there, and the same bound gives them, their derivatives and H\"older quotients (kernel by kernel: \cite[(2.34), (2.36)]{Ba1} for $C^{(i)}$, $\delta C_{\Lambda_0}$, the composition identity of (Z5) for $G^\eta_{i+1}Q_{i+1}^*$, the $C^{1,1}$ profile at scales $i$, $i+1$); the second term $T_2(x)=a_iaL^{-2}\theta_{i+1}(x)\sum_y\zeta^{(i)}(x,y)(G_iQ_i^*)(x,y)(\delta C^{(i)}_{\Lambda_0}Q^*\Lambda_0^{(i)\prime}B)(y)$ has, for $x\in\supp\theta_{i+1}$ and $\zeta(x,y)\ne0$, $\dist(y,\Lambda_0^{(i)c})>r_i-M_\theta-2$ (two layers of \cite[(2.8)]{Ba2} minus the range of $\zeta$), and $\delta C_{\Lambda_0}(y,y')$ decays exponentially in $|y-y'|$ and in the distances of $y$, $y'$ to $\Lambda_0^{(i)c}$ \cite[(2.35)]{Ba1}, so $|T_2|+|\partial^\eta T_2|\le Cr_i^de^{-\delta'(r_i-M_\theta-2)}p/(\lmu\ell)\le C_\kappa\ell_i^\kappa$. (Z4): $\frac12\theta r_i-2M\ge2$ iff $r_i\ge(4+4M)/\theta=10(M+1)$, and $r_i\ge R$. (Z5): on the unit lattice of scale $i$, with $Q$ the one-step average, $P_i=Q_i^*Q_i$, $P_{i+1}=Q_i^*Q^*QQ_i$, $G_i=(H+a_iP_i)^{-1}$, $G^\eta_{i+1}=(H+a_{i+1}L^{-2}P_{i+1})^{-1}$, $\Delta^{(i)}=a_i-a_i^2Q_iG_iQ_i^*$, $C^{(i)}=(aL^{-2}Q^*Q+\Delta^{(i)})^{-1}$ and $a_{i+1}=aa_i/(a_i+aL^{-2})$ \cite[(2.13), (2.21), (2.31)]{Ba1}, one has the composition identity $a_iaL^{-2}G_iQ_i^*C^{(i)}Q^*=a_{i+1}L^{-2}G^\eta_{i+1}Q^*_{i+1}$: apply $H+a_{i+1}L^{-2}P_{i+1}$ to the left side, use $HG_i=1-a_iP_iG_i$, $1-a_iQ_iG_iQ_i^*=\Delta^{(i)}/a_i$, $\Delta^{(i)}C^{(i)}=1-aL^{-2}Q^*QC^{(i)}$ and $Q^*QQ^*=Q^*$; the coefficient of $Q^*QC^{(i)}Q^*$ vanishes iff $a_{i+1}(a_i+aL^{-2})=aa_i$, which is \cite[(2.13)]{Ba1}, and the remaining coefficient is $a_{i+1}L^{-2}$ by the same relation (this is \cite[(2.41)]{Ba1} in adjoint, rescaled form). By \cite[(2.80)]{Ba2} and the Dirichlet condition \cite[(2.32)]{Ba1}, $C^{(i)}_{\Lambda_0}Q^*B=C^{(i)}_{\Lambda_0}Q^*\Lambda_0^{(i)\prime}B$; writing $C_{\Lambda_0}=C^{(i)}+\delta C_{\Lambda_0}$ and $\zeta=1-(1-\zeta)$ and using the identity, $\theta_{i+1}A^{(i)}$ equals the five terms of \cite[(2.89)]{Ba2} with $\Lambda_0^{(i)\prime}B$ in the fifth term, identically in $A'$, $B$; the printed form differs by the defect $a_{i+1}L^{-2}\theta_{i+1}\zeta^{(i+1)}_\theta G^\eta_{i+1}Q^*_{i+1}\Lambda_0^{(i)\prime c}B$, which vanishes when $\zeta^{(i+1)}_\theta(x,y)=0$ for all $x$ with $\theta_{i+1}(x)\ne0$ and $y\notin\Lambda_0^{(i)\prime}$, i.e.\ when $L(\theta r_{i+1}-2M)\le2r_i-M_\theta-2-L$ (such $x$ lie within $M_\theta+1$ of $\Lambda_2^{(i)}$, such $y$ within $L$ of $\Lambda_0^{(i)c}$, and $\dist(\Lambda_2^{(i)},\Lambda_0^{(i)c})>2r_i-1$). With $L\theta_5=2$ this reads $2(r_i-r_{i+1})\ge M_\theta+7-10M$, and $r_i-r_{i+1}=R\ln5\,(2u_{i+1}+\ln5)\ge5.81R$. The balls around $\supp\theta_l$ reach at most $2r_l-10M+M_\theta+1$ beyond $\Lambda_2^{(l-1)}$ (scale-$(l-1)$ units), while $\Lambda_{-1}^{(l-1)}$ contains the $(3r(\ell_{l-1})-2)$-neighbourhood of $\Lambda_2^{(l-1)}$ (\cref{lem:a01:X0}(a)): $2r_l+M_\theta+8\le3r(\ell_{l-1})+10M$ holds since $r_l\le r(\ell_{l-1})$ and $r(\ell_{l-1})\ge R\ge M_\theta+2$. For the printed $\theta=1$ the defect condition has $L\theta=L$, $\beta=2$, so it holds at every scale iff $L\le2$; on the reading that absorbs the defect into $\tilde B^{(i)}$, the $\zeta^{(i+1)}$-balls must stay inside $\Lambda_{-1}^{(i)\prime}$, $\beta=3$, hence $L\le3$.
\end{proof}
The other uses of \cite[Lemma~2.3]{Ba2} are covered by (Z2): \cite[(2.68)]{Ba2} on the box of (c) below; \cite[(2.95)--(2.98)]{Ba2} at each scale $l$ together with (Z5); the regularity of $\tilde B$ needed by \cite[Prop.~1]{Ba3} at the background $\theta_iA^{(i)}$ on $\Omega$, where $\theta_i=1$; and the first-step use \cite[(2.25)]{Ba2}, whose cut-off has $\theta=1/(2L)<3/L$. On the consumer side, the all-small-field term of \cite[(3.4)]{K1} is affected only through the cut background, and the difference between the actions at $A^{(j)}$ and at $A_\theta$ is $O(\eta^{s'})|X^{(j)}|$ on $\supp\chi_j$, as used in the proof of \cref{prop:a01:upper}.

\subparagraph{(c) The boxes of \cite[Lemma~2.4]{Ba2}.} The printed proof \cite[p.~572]{Ba2} takes $\square_1,\square_2$ as the sums of large blocks of $\Lambda_7^{(i-1)\prime}$ within $2r_i$, $4r_i$ of $y$ and $\square=B^i(\square_2)$. For $y\in\Lambda_7$ at distance $O(1)$ from $\Lambda_7^c$ the step \cite[(2.67)]{Ba2} is not justified: the dropped source $\Lambda_6\phi-\square_1\phi$ contains $\phi$ on $\Lambda_6\setminus\Lambda_7$ at distance $O(1)$ from $x$, where $a_iG_iQ_i^*$ is $O(1)$, so this error is $O(\Upsilon_i)$, $\Upsilon_i\defeq c_1(\lambda\ell_{i-1})^{-1/4}p_{i-1}$, not $O(\ell_i^\kappa)$; the domain change $\Omega\to\square$ happens at distance $O(1)$ from $x$, so \cite[(2.26)]{Ba1} gives no smallness; and $x$ may lie within $R_{\mathrm{Ba}}$ of $\partial\square$, where \cite[(2.24)--(2.25)]{Ba1} need a parallelepiped with a constant field near its boundary, while $\Lambda_7\cap(\text{cube})$ has re-entrant edges. An untruncated box of radius $4r_i$ inside $\Omega$ would need $4r_i\le5\rho_i$, which fails for small $\ell_i$ at every $L\ge2$ ($\alpha=4$, $\beta=5$). The only consumer is the quartic convergence factor \cite[(2.78)--(2.79)]{Ba2}, a sum over $y\in\Lambda_7^{(i-1)\prime}\cap\Lambda_7^{(i)c}$, which contains every point of $\Lambda_7$ within $2r_i$ of $\Lambda_7^c$: the edge points are exactly the ones used.

\begin{lemma}[{\cite[Lemma~2.4]{Ba2}} with a complete proof, every $L$]\label{lem:a01:lem24}
Assume \cite[(2.55)]{Ba2}, and: $5(r(\eps_0^{\mathrm{Ba}})-s_b)/L\ge R_{\mathrm{Ba}}+c_d+2s_b$, $4r(\eps_0^{\mathrm{Ba}})\ge R_{\mathrm{Ba}}+c_d$, $R\ge10s_b+10$, with $c_d\defeq\sqrt d\,s_b$; $\eps_0^{\mathrm{Ba}}$ small enough that $r(\ell_{i-1})/r(\ell_i)\le2L$ for all $i$ (at $L=5$, $\varrho=2$ this is $(1+\ln5)^2\le10$, true for all $\eps_0^{\mathrm{Ba}}\le1$); the mass condition $\Upsilon_im^2\ell_i^2(G_i(\square,0)Q_i^*1)(x)=O(p_i)$ (Ba{\l}aban's own use of \cite[(2.75)]{Ba2}); and $Ca_i\Upsilon_i(1+r_i)^{d-1}\exp(-\delta_0[(r(\ell_{i-1})-s_b)/L-c_d])\le\ell_i^\kappa$, which holds for all small $\eps_0^{\mathrm{Ba}}$ at every $R>0$ because $\varrho>1$. Then \cite[(2.65)--(2.66)]{Ba2} hold for every $y\in\Lambda_7^{(i-1)\prime}$, $x\in B^i(y)$, $b\subset B^i(\Lambda_7^{(i-1)\prime})$, with constants independent of $\eps$, $i$, the lattice and the position of $y$ relative to $\partial\Lambda_7$.
\end{lemma}
\begin{proof}
Put $s_i\defeq\min\{4r_i,5\rho_i-2s_b\}$, let $\square_2$ be the cube of large blocks of sup-radius $s_i$ around the large block of $y$, $\square\defeq B^i(\square_2)$, and $\square_1\phi\defeq1_{Q(y,s_i/2)}\Lambda_6\phi$ with $Q(y,s)$ the sup-ball. Since $\dist(\Lambda_7,\Lambda_2^c)\ge5\rho_i$ (\cref{lem:a01:X0}(a)), $\square\subset\Omega$; $\square$ is a rectangular parallelepiped and a union of large blocks, and $s_i-c_d\ge R_{\mathrm{Ba}}$ at every $i$ by the hypotheses. (1) \emph{Source cut:} $\Lambda_6\phi-\square_1\phi$ is supported at distance $\ge s_i/2-c_d$ from $x$, and $\dist(x,\Omega^c)\ge5\rho_i-c_d\ge R_{\mathrm{Ba}}$, so \cite[(2.25)]{Ba1} on $\Omega$ with the sup bound $\Upsilon_i$ on $\phi$ gives $|a_i(G_i(\Omega,A^{(i)})Q_i^*(\Lambda_6\phi-\square_1\phi))(x)|\le Ca_ie^{-\delta_0(s_i/2-c_d)}\Upsilon_i$. (2) \emph{Domain change:} $G_i(\square,A)=G_i(\Omega,A)+\delta G_i(\square,\Omega,A)$, and by \cite[(2.26)]{Ba1} the second term applied to $Q_i^*\square_1\phi$ at $x$ carries $e^{-\delta_0(3s_i/2-2c_d)}$; the same for $D^\eta$ by \cite[(2.58)]{Ba2}. Hence $\phi^{(i)}(x)=a_i(G_i(\square,A^{(i)})Q_i^*(A^{(i)})\square_1\phi)(x)+O(a_i\Upsilon_ie^{-\delta_0(s_i/2-c_d)})$, and likewise for the covariant derivative: this is \cite[(2.67)]{Ba2} with the boxes chosen above. (3) \emph{Constant background:} $A_0\defeq A^{(i)}(y)$; joining points of the convex box by straight lattice paths and using \cite[(2.60)]{Ba2} (\cref{prop:a01:Zth}(Z2)) on $B^i(\Lambda_2^{(i-1)\prime})\supset\square$, $|A^{(i)}(z)-A_0|\le O(p)\sqrt d\,s_i\le O(pr_i)$ on $\square$; the expansion \cite[(3.44)]{Ba1} and \cite[Prop.~2.2]{Ba1} on the parallelepiped $\square$ replace $A^{(i)}$ by $A_0$ with error $O(\ell_i^{\kappa_0})$, as in \cite[(2.68)]{Ba2}. (4) \emph{Gauge-out:} $A_0=\partial^\eta\lambda_0$ with $\lambda_0(z)=A_0\cdot(z-y)$ on $\square$, so $U(A_0(\Gamma_{z,y}))=e^{ieA_0\cdot(z-y)}$ for every contour inside $\square$, $G_i(\square,\partial\lambda_0)=U(\lambda_0)G_i(\square,0)U(\lambda_0)^*$ and $Q_i(\partial\lambda_0)=U(\lambda_0)_{\mathrm{coarse}}Q_i(0)U(\lambda_0)^*$: \cite[(2.69)--(2.74)]{Ba2}; legitimate because $\square$, of diameter $\le8r_i+s_b$, does not wrap around (on the cube, no condition). Write $\phi'\defeq U(\lambda_0)^{-1}\phi$. (5) \emph{Reproduction of constants:} by \cite[(2.75)]{Ba2}, $a_iG_i(\square,0)Q_i^*1=1-m^2\ell_i^2G_i(\square,0)Q_i^*1$, hence $a_i(G_i(\square,0)Q_i^*\square_1\phi')(x)-\phi'(y)=a_i(G_i(\square,0)Q_i^*[\square_1(\phi'-\phi'(y))])(x)-\phi'(y)a_i(G_i(\square,0)Q_i^*1_{\square\setminus\square_1})(x)-\phi'(y)m^2\ell_i^2(G_i(\square,0)Q_i^*1)(x)$. The first term is $O(p)$: on the convex cube $Q(y,s_i/2)\subset\Lambda_{-1}$ the restriction $|D\phi|\le c_1p$ gives $|\phi'(z)-\phi'(y)|\le O(p)|z-y|$ and the kernel decays exponentially. The second: $\square\setminus\square_1\subset(Q(y,s_i)\setminus Q(y,s_i/2))\cup(Q(y,s_i/2)\setminus\Lambda_6)$ lies at distance $\ge\min\{s_i/2,\rho_i\}-c_d$ from $x$, so it is $\le c_{\mathrm{Ba}}e^{-\delta_0(\min\{s_i/2,\rho_i\}-c_d)}\Upsilon_i$ by \cite[(2.25)]{Ba1} on the parallelepiped. The third is the mass-defect term of the printed proof, $O(p_i)$ under the mass condition. (6) \emph{Derivatives:} a bond $b\subset B^i(y)$ lies at distance $\ge s_i-c_d\ge R_{\mathrm{Ba}}$ inside $\square$ and $\ge5\rho_i-c_d$ inside $\Omega$, so the printed derivative bounds apply at $b$ and (1)--(5) give \cite[(2.66)]{Ba2} in the form \cite[(2.77)]{Ba2}; no re-entrant edge is approached. \emph{Errors:} $s_i/2\ge\rho_i$ iff $\rho_i\le2r_i$ and $3\rho_i\ge2s_b$, the first by $r(\ell_{i-1})/r_i\le2L$, the second by $R\ge10s_b+10$; so every error is $\le Ca_i\Upsilon_i(1+r_i)^{d-1}e^{-\delta_0(\rho_i-c_d)}$, which is $\le\ell_i^\kappa$ by the last hypothesis; the rest of the printed proof, \cite[(2.69)--(2.77)]{Ba2}, runs on $\square$. Nothing depends on $L$ except through the named constants.
\end{proof}
At $L=5$ the hypotheses read $r(\eps_0^{\mathrm{Ba}})\ge R_{\mathrm{Ba}}+(3+\sqrt3)s_b$, $r(\ell_{i-1})/r_i\le6.81\le10$, $R\ge10s_b+10$: lower bounds on $R$ and caps on $\eps_0^{\mathrm{Ba}}$, none restricting $L$. On the cube the same device (a box of large blocks inside the propagator domain about the collar cubes, with the cap inherited from the previous step) is the King box of D8.

% a01_H0rel part 6x3: the replaced step (d) and the parameter conditions (input from a01_H0rel_p6.tex after a01_p6x2)
\subparagraph{(d) The partition \cite[(3.43)--(3.45)]{Ba2}.} The printed ``regular partition into a lattice of cubes, each a sum of large blocks, with side in $(r,2r)$'' needs, per coordinate, a divisor $c$ of the number $N$ of large blocks with $r/s_b<c<2r/s_b$; at $L=2$ a power of $2$ always works, while at $L=5$ the divisors of $N_i=2MJ/(\ell_is_b)$, a number of the form $2\cdot5^{b}J$, have gaps of ratio up to $\frac52$, so the window can contain no divisor, and with integer sides and even $N$ the printed partition fails exactly when $r_i/s_b$ is an integer. The partition is used in \cite[Sect.~3C]{Ba2} only through four properties: the covering \cite[(3.43)]{Ba2}; the neighbour count $3^d$ in \cite[(3.44)--(3.45)]{Ba2}; the volume bounds $(2r)^d$ in \cite[(3.48)--(3.49)]{Ba2}; and the corridor step of \cite[p.~593]{Ba2}.

\begin{lemma}[the product partition and the cube count]\label{lem:a01:PP}
Let $q_i\defeq\lceil(r_i+s_b)/s_b\rceil$ and let every coordinate carry $N\ge q_i$ large blocks (on King's cube, $N=2MJ/(\ell_is_b)$ large blocks per side in scale-$i$ units, which exceeds $q_i$ for $J\ge J_0$; on a torus, $N$ per period). Cut each coordinate interval into consecutive arcs of between $q_i$ and $2q_i-1$ large blocks (write $N=nq_i+\varrho'$, $0\le\varrho'<q_i$, and give the $\varrho'$ extra blocks to the $n$ arcs, at most $q_i-1$ each), and let $\mathcal Q_i$ be the product partition into boxes. Call $Q\sim Q'$ if their arc indices differ by at most one in every coordinate. Then: \textup{(Q1)} every box is a union of large blocks; \textup{(Q2)} every box is a product of intervals of at most $2r_i+3s_b\le(2+3s_b/R)r_i$ points; \textup{(Q3a)} every box has at most $3^d$ boxes $Q'\sim Q$; \textup{(Q3b)} if $\dist(z,x)\le r_i+s_b$ then $Q(z)\sim Q(x)$. Consequently, with $I_i\defeq\{z:\dist(z,(\Lambda_7^{(i-1)\prime})^c)\ge s_b\}$ and $\mathcal C_i\defeq\{Q\in\mathcal Q_i:Q\cap\Lambda_0^{(i)c}\cap I_i\ne\emptyset\}$, every $Q\in\mathcal C_i$ satisfies $Q\sim Q(x_e)$ for the charge point $x_e$ of some element $e$ of the family $F_i$ of large-field events at scale $i$, hence $|\mathcal C_i|\le3^d|F_i|$; the volume bound $(2r)^d$ becomes $(2r_i+3s_b)^d$ and the corridor bound $22r$ becomes $22r_i+3s_b$, both $O(1)r_i^d$.
\end{lemma}
\begin{proof}
(Q1) Arcs are unions of large-block intervals. (Q2) An arc has at most $(2q_i-1)s_b<2(r_i+s_b)+s_b$ points, and $s_b\le(s_b/R)r_i$. (Q3a) The boxes $\sim Q$ are the $3^d$ neighbours of $Q$ in the product structure. (Q3b) Every arc has $\ge q_is_b\ge r_i+s_b$ points, so if $|z-x|_\infty\le r_i+s_b$ the two arcs in each coordinate are equal or consecutive (otherwise a whole arc lies strictly between and the coordinate distance is $\ge q_is_b+1$). For the count: let $Q\in\mathcal C_i$ and $z\in Q\cap\Lambda_0^{(i)c}\cap I_i$; its large block $\beta$ lies in $Q$ by (Q1) and, having sup-diameter $<s_b$, in $\Lambda_7^{(i-1)\prime}$; since $z\notin\Lambda_0^{(i)}$, $\beta$ is not one of the blocks of \cite[p.~570]{Ba2}, so $\dist(\beta,S_i)\le r_i$ where $S_i$ is the point set of $F_i$; take $z'\in\beta$, $y\in S_i$ with $\dist(z',y)\le r_i$, $y$ a point of an element $e$ with $\dist(y,x_e)\le L-1<s_b$; then $\dist(z',x_e)\le r_i+s_b$ and $Q=Q(z')\sim Q(x_e)$ by (Q3b); count with (Q3a). The printed \cite[(3.45)]{Ba2} and the first inequality of \cite[(3.46)]{Ba2} do not hold as printed (a large block straddling $\partial\Lambda_7^{(i-1)\prime}$ lies in $\Lambda_0^{(i)c}$ with $F_i$ possibly empty); the count above, on the interior $I_i$, becomes the count of D14 over the paid elements through \cref{lem:a01:PP5}; in D14 the boxes of $\mathcal C_{i'}$ for $i'<i$ are carried along by the corridor construction of \cite[(3.48)--(3.50)]{Ba2} with the thickness $7(r_{i-1}+s_b)+Ls_b$ inside Ba{\l}aban's $(9r,10r)$ once $(7+L)s_b<2r_{i-1}$.
\end{proof}
% a01_H0rel part 6x3, continued: the count over the paid elements (input from a01_p6x3.tex after lem:a01:PP)
\begin{lemma}[the count over the paid elements]\label{lem:a01:PP5}
Let the boxes $\mathcal Q_i$, the relation $\sim$, the arcs of $q_i$ to $2q_i-1$ large blocks, the set $\mathcal C_i$ and the family $F_i$ be those of \cref{lem:a01:PP}, with the charge point $x_e$ of an element $e$ as in its proof (every point of $e$ lies within $\dist\le L-1$ of $x_e$). Let $F_i^{\mathrm{paid}}\subset F_i$, and suppose that every $e\in F_i\setminus F_i^{\mathrm{paid}}$ has a \emph{witness} $w(e)\in F_i^{\mathrm{paid}}$ with $\dist(x_e,x_{w(e)})\le r_i+s_b$. Then
\[
|\mathcal C_i|\ \le\ 5^d\,|F_i^{\mathrm{paid}}|.
\]
In D14 the unpaid elements are the near $R_s$ elements, the blocks $y\in R_s^{(i)}$ within $\ell^\infty$-distance $R_0$ of the face layer, and $F_i^{\mathrm{paid}}$ consists of Ba{\l}aban's six families, $R_s$ restricted to its far elements, and the face family; the witness of D7(b) (\cref{lem:a01:F541}(e)) has a point within distance $dR_0+1$ of a point of $e$, so $\dist(x_e,x_{w(e)})\le dR_0+1+2(L-1)\le s_b+L-2<r_i+s_b$ by $s_b\ge dR_0+L+1$ (D7(c)) and $r_i\ge R>L$, and the hypothesis holds.
\end{lemma}
\begin{proof}
Let $Q\in\mathcal C_i$. The proof of \cref{lem:a01:PP} gives $z'\in Q$ and $e\in F_i$ with $\dist(z',x_e)\le r_i+s_b$. Put $e'\defeq e$ if $e\in F_i^{\mathrm{paid}}$ and $e'\defeq w(e)$ otherwise; in both cases $e'\in F_i^{\mathrm{paid}}$ and $\dist(z',x_{e'})\le2(r_i+s_b)\le2q_is_b$, since $q_is_b\ge r_i+s_b$. In each coordinate the arc indices of $Q=Q(z')$ and of $Q(x_{e'})$ differ by at most $2$: if they differed by $3$ or more, two whole arcs, of at least $q_is_b$ points each, would lie strictly between the two coordinates, whose distance would then be $\ge2q_is_b+1$. A box has at most $5^d$ boxes whose arc indices differ from its own by at most $2$ in every coordinate (fewer on a torus with fewer than five arcs in some coordinate). So $Q\mapsto e'$ maps $\mathcal C_i$ into $F_i^{\mathrm{paid}}$ with fibres of at most $5^d$ boxes, which is the claim. For the application: a near $R_s$ element $e$ at the block $y$ has $\dist(y,x_e)\le L-1$; \cref{lem:a01:F541}(e) gives, on a block path of length $\le dR_0$ from $y$ to the face layer, a bond of the gradient family or a block of the face family, that is an element $w(e)$ of a paid family one of whose points lies within distance $dR_0+1$ of $y$ and within $L-1$ of $x_{w(e)}$; the triangle inequality gives $\dist(x_e,x_{w(e)})\le dR_0+2L-1$, and $dR_0+2L-1\le s_b+L-2$ is $s_b\ge dR_0+L+1$.
\end{proof}


The replaced objects of (a)--(d), $r^\flat$, $D^\flat$, $\zeta^{(i)}_\theta$, the box $\square$ of \cref{lem:a01:lem24} and the partition $\mathcal Q_i$, are built from $(L,M,R,\ell_i,\bar n)$ and the set of event blocks alone, so they are the same on the cube as on the torus (\cref{thm:a06:BALA}); they are auxiliary objects of the proofs, not of the statements, and under (b) the all-small-field action changes by $O(\ell_i^\kappa)|X^{(i)}|$, which the statements absorb. The parameter conditions they add are lower bounds on $R$ ($R\ge R^\flat$; $R\ge10(M+1)$, $R\ge M_\theta+2$; $R\ge10s_b+10$; $R\ge6L\sqrt d$ for the covering of D14) and caps on $\eps_0^{\mathrm{Ba}}$, all fixed in Ba{\l}aban's order of choices ($d$, $a$, $\bar n$, $\alpha_0$, $L$, $M$, $R$, $\varrho$, $p$, $b_0$, then $\eps_0^{\mathrm{Ba}}$), so the order is acyclic.


\begin{theorem}[the perturbative bounds on sub-domains]\label{thm:a01:B2}
Fix a scale $k\le K_f$, a union $\Omega$ of $k$-blocks of the cube (any shape, meeting $\DD$ or not), $A$ regular on $\Omega$ with $\theta_k\le\theta_0$, and a renormalised family $\mathcal R^\Omega$ of order $\le13$ of Ba{\l}aban's expansion with the lines $G_{(i)}(\Omega,A)$, $G_k(\Omega,A)$, the legs $\phi^{(k)}_\Omega$, photon legs of the class of \cite[p.~412]{Ba3}, possibly source vertices, possibly one difference line $\delta G_k(\Omega,\Omega_2;A)$ of two Neumann domains as an external pair, with at least one external field or source vertex, and all vertex cubes at depth $\ge R_*\defeq48$ in every propagator domain. Then for one held vertex $x_0$ (or none),
\[
\sum_{\text{orderings, indices, positions}}|E^{(k)}(G)|\ \le\ C\ell_k^{\beta_0}p(\ell_k)^{n_1}(1+K-k)^{p_0}e^{-\delta\dist(V_*)}\Lambda(x_0)\,\Pi_b,
\]
with $\Lambda(x_0)$ as in \cref{prop:a01:TD}(a), $\Pi_b=1$ without a difference line and $\Pi_b=\exp[-\delta_c(\dist(\mathrm{cube}(v),\Xl\setminus\Omega_2)+\dist(\mathrm{cube}(v'),\Xl\setminus\Omega_2))]$ for the ends $v,v'$ of the difference line; the constants do not depend on $k$, $K$, $\vth$, $J$, $\Omega$, $A$ or $x_0$. The same holds for a finite set of such difference lines and difference legs, each contributing its own factor, with the constant multiplied by $C_1^{|S|}$.
\end{theorem}
\begin{proof}
The proof of \cref{prop:a01:TD}(a) is run on $\mathcal R^\Omega$ with the following changes. (1) \emph{Localisation.} Every line and leg of the family is localised on a Dirichlet box $Q_x$ of large blocks around its vertex cube, inside $\Omega$ at depth $\ge R_*$: the difference between the operator on $\Omega$ and the operator on $Q_x$ is a cut-off and commutator argument of the type of \cref{lem:a01:FLD} (a smooth cut-off at the Neumann faces of $Q_x$, the commutator bounded by \cref{lem:a01:A5}, the result bounded by \cref{lem:a01:A4}(c)), exponentially small in the depth, and the box-vs-whole-cube difference near $\DD$ is \cref{lem:a01:LOC}; so every line at a vertex of depth $\ge R_*$ is the line of the box up to $Ce^{-\delta R_*}$ times the box line, with the Combes--Thomas rate, free of Ba{\l}aban's $M$. (2) \emph{Zero-background identities on the box.} The five identities used in the organisation of \cite[p.~330]{K2} (the free counterterm, the bulk rewriting of the renormalised families, the constant-datum identity \cref{lem:a01:RC}, the slice identity, and the Ward-type cancellation of the zeroth photon moment, which on a box is exact since a constant gauge field is pure gauge there) hold on the Dirichlet box by the odd-image representation of Step~1 of \cref{lem:a01:FL}, and at a regular background by gauging away the constant part of $A$ on the box (gauge covariance, \cref{app:a01:weights}) and expanding to first order in the remainder, whose coefficients are $O(\theta_k(1+|x-x_0|))$ (\cref{lem:a01:FL}, Step~3); the first-order photon tadpole at a background is bounded by the same expansion. (3) \emph{Remainder vertices.} Each boundary-remainder vertex is a layered vertex of the class of \cref{prop:a01:TD}(a) conjugated by the gauge transformation of (2), plus an ordinary vertex of degree raised by $D'+2\ge1$, plus an any-degree term with the small factor of (1); conjugation does not change absolute values, so the degree bookkeeping of \cite[Sect.~2]{Ba3} is unchanged and the extra vertices only raise the containing stages. (4) \emph{Lines near the face, difference lines.} The decay along paths of the lines on $\Omega$ is \cref{lem:a01:A23,lem:a01:A4}; the H\"older clause of a difference line near $\DD$ comes from the proof of \cref{lem:a01:A5}, cases (3)--(4), applied to $w=\Mop(\Omega_2)^{-1}g$ with $g$ on the cut shell. The difference line $\delta G_k(\Omega,\Omega_2)=G_k(\Omega)-G_k(\Omega_2)$ is $G_k(\Omega)Y G_k(\Omega_2)$ with $Y$ the bond terms of the cut, so its kernel between two points carries the factor $e^{-\delta_c(\dist(\cdot,\Xl\setminus\Omega_2)+\dist(\cdot,\Xl\setminus\Omega_2))}$ by Combes--Thomas on both sides, which is $\Pi_b$. (5) \emph{Telescoping.} A family some of whose lines are written $o=o'-\Delta o$ is the finite sum $\sum_S\mathcal R^{(S)}$ over the sets $S$ of decomposed objects, $\mathcal R^{(S)}$ carrying $\prod_S(-\Delta o)$ and $\prod_{O\setminus S}o'$; $\mathcal R^{(\emptyset)}$ is the family with every object replaced by its counterpart on $\Omega$, with the same counterterms, and for $S\ne\emptyset$, $\mathcal R^{(S)}$ is the renormalised family of the cut graphs $G\setminus S$, with every subtraction of the divergent subgraphs of $G\setminus S$ and the same counterterm values, the divergent subgraphs containing a line of $S$ carrying no subtraction: this is the regrouping \cref{hyp:k01:RENorb}, and no subtraction is lost. Each $\mathcal R^{(S)}$ is then a family with $|S|$ admissible external pairs, bounded by (1)--(4) with the factor $\Pi_S=\prod_{\sigma\in S}C_\sigma e^{-\delta_\sigma(\dots)}$; the constant per pair is $C_1=C_1(d,L,\bar n)$. (6) \emph{Depth.} At every consumer (D5, D6, D7, D10 below) the vertex cubes lie in a $\Lambda_7$-level set and the propagator domains are $\Lambda_2$-level, five nested layers apart (\cite[p.~411]{Ba3}, \cite[(2.53), (2.56), Prop.~2.6]{Ba2}), so their depth in every propagator domain is $\ge\lfloor R\rfloor-1\ge51>R_*$ once $R\ge\max(52,\lceil R_{\mathrm{Ba}}\rceil+1)$, where $R_{\mathrm{Ba}}$ is Ba{\l}aban's depth constant of \cite[(2.19)]{Ba4}.
\end{proof}

\begin{proposition}[changes of propagator domain]\label{prop:a01:dGD}
At every scale $k\le K_f$ of the upper bound, the domain-change statements of \cite[(2.18)--(2.50), (2.108)--(2.115)]{Ba2} hold on the cube with the weighted flow: (i) the comparison of a Dirichlet-conditioned covariance $C_\Lambda$ with the covariance of the Neumann domain $B^k(\Lambda_2^{(k)})$, $C_\Lambda-C|_{\Lambda\Lambda}=C_\Lambda X_{\Lambda\Lambda^c}C_{\Lambda^c\Lambda}$, is $O(M_{k+1}e^{-cr(\ell_k)})$ in the external-field norm; (ii) the scalar translation \cite[(2.23)]{Ba2} and its replacement \cite[(2.25)]{Ba2} differ by $O(M_1e^{-\delta r(\eps)})=O(\eps^{\kappa'})$ for every $\kappa'$; (iii) the change \cite[(2.108)]{Ba2} of the block operators $\Delta^{(k)}$ between $B^k(\Lambda_2^{(k-1)\prime})$, $B^k(\Lambda_2^{(k-1)\prime}\cap\Lambda_5^{(k)c})$ and $B^k(\Lambda_2^{(k)})$ obeys \cite[(2.109)]{Ba2}, $|h_k(x,x')|\le Ce^{-c\min(r(\ell_k),4r(\ell_{k-1})/L)}e^{-\delta|x-x'|}$; (iv) the Neumann cuts along $\partial\Lambda_6$, $\partial\Lambda_7$ inside the Dirichlet-restricted covariances \cite[(2.38)--(2.40)]{Ba2} cost $O(\eps^{\kappa'})|\partial\Lambda_5|$, resp.\ $O(\eps^{\kappa'})|\Lambda_5|$; (v) the exact recursion \cite[(2.41)]{Ba1} on the Neumann domain $B^k(\Lambda_2^{(k)})$ holds (it uses only the uniform coarse step); (vi) the restricted fluctuation covariance $C^{(k)}_{\Lambda_5^{(k)}}$ of \cite[(2.115)]{Ba2} and the Neumann hole of \cite[(2.46)--(2.47)]{Ba2} cost $O(\ell_{k-1}^{\kappa'})|\Lambda_6^{(k-1)}|$; (vii) the photon statements are those of the torus by \cref{lem:a01:PE}(c). The constants depend on $(d,L,a,\Wt,c_V)$ and King's and Ba{\l}aban's constants only, not on $J$, $K$, $k$, $\eps$, the domains (any unions of blocks in the classes where the objects are defined, the face met or not, re-entrant edges allowed) or the field.
\end{proposition}
\begin{proof}
Each statement is an identity of resolvents (the bond cut $C_\Lambda-C'_\Lambda=-C_\Lambda Y_{\Lambda\Lambda}C'_\Lambda$, the restriction identity above, the Schur complement for the hole) followed by the Combes--Thomas decay of \cref{lem:a01:A23}(a) on each factor and the coercivity \cref{lem:a01:A1} on every domain (which is what replaces Ba{\l}aban's $\gamma_0'$; the step-$0$ lower bound for every field is the Neumann gap with the weight, as in the proof of \cref{lem:a01:A1}), the sup bounds of \cref{lem:a01:A4} and the legs' norms of \cref{lem:a01:legs} on the domains involved; the changes relative to the torus are constants only (the factor $\Wt$ from (Q5), (Q6) wherever a block average is bounded, and $\gamma_0''$ of \cref{lem:a01:A23}(c) in place of Ba{\l}aban's $\gamma_0'$). The distances $r(\ell_k)$, $4r(\ell_{k-1})/L$ between the sets compared are Ba{\l}aban's region geometry, unchanged. The photon part is \cref{lem:a01:PE}(c).
\end{proof}


% a01_H0rel part 6b: the steps of the upper bound
\paragraph{The steps of the upper bound.}
\begin{description}[leftmargin=1.5em,style=nextline]
\item[{D1 (large/small split; \cite[(2.2)--(2.3)]{Ba2})}] At each step the fields $B,A,\psi,\phi$ are tested; each test is tied to a positive term of the action and a violation yields $e^{-c_0^{\mathrm{Ba}}p(\ell_i)^2}$; with the Dirichlet condition the matter Gaussian only gains the positive face terms, which lower the integrand, and every printed large-field factor is monotone in them.
\item[D2 (the face family)] At every scale $0\le i\le K_f$ the face test with threshold $\sigma_i$ is added to the six families (\cref{def:a01:facefam}); a violation is paid $e^{-p_i^2/4}$ by \cref{lem:a01:F541}(a) from the $\theta_2$-share of the scale-$i$ form, the $(1-\theta_2)$-share remaining for D12. The family is forced: without it the background of face data of size $M_i$ has a face gradient of order $M_i$ (\cref{prop:a01:legsd}, remark), which violates the fine-field class \cite[(2.55)]{Ba2}.
\item[{D3 (region geometry; \cite[(2.4)--(2.10), (2.43)]{Ba2})}] As said above, verbatim with seven families. The pair $B^i(y)\cup B^i(y')$ of \cref{lem:a01:FPC}'s payment is one block here, so no clean-pair condition arises; tested blocks are clean (no Neumann cut through $B^i(y')$) by the nesting rules.
\item[{D4 (scalar large fields; \cite[(2.11)--(2.13)]{Ba2})}] $\lambda(\eps)|\phi|^4+\frac12\dm\eps^2|\phi|^2\ge-O(\eps^{\beta_0})$ and $\ge\frac12\lambda(\eps)|\phi|^4$ for $|\phi|^2\ge O(1+\ln\eps^{-1})$: unchanged, $\dm$ being King's free-boundary counterterm, a vertex, weight-independent.
\item[{D5 (translation and splitting; \cite[(2.18)--(2.19)]{Ba2})}] The photon translation with the Dirichlet-conditioned $C^{(0)}_{\Lambda_0}$ is \cref{lem:a01:PE}(c); the scalar analogue uses the operators on sub-regions with \cref{lem:a01:A1,lem:a01:A23,lem:a01:A4,lem:a01:DV,lem:a01:DVloc} at every point and \cref{lem:a01:A5,lem:a01:legs} in the restricted form (depth $\ge R_{\mathrm{Ba}}$, which is where Ba{\l}aban's splitting evaluates them: his vertex sets are $\Lambda_7$-level and his propagator domains $\Lambda_2$-level), \cref{thm:a01:B2} for the perturbative terms on these sub-domains and \cref{prop:a01:dGD}(i), (ii), (iv), (vii) for the changes of domain and of boundary condition. No face hypothesis on the splitting domains is needed on the cube.
\item[{D6 (expansion to order $\bar n$, cumulants, removal of the interaction from $\Lambda_7^c$; \cite[(2.22)--(2.28), Prop.~2.6]{Ba2})}] \Cref{prop:a01:TD}(c) and \cref{thm:a01:B2} with the difference lines of the pair of propagator domains of \cite[Prop.~2.6]{Ba2} as external pairs (the telescoped families are finite signed sums of members with $s\ge1$ difference objects, each admissible with the gain distance $\dist(\cdot,\Xl\setminus\Omega_2)$, so the extra factor is $e^{-5\delta_cr(\ell_k)}=O(\ell_k^{\kappa})$ at the margin $>5r(\ell_k)$ of Ba{\l}aban's Proposition~2.6: this is the interaction-term bound \cref{hyp:k01:BALint} on the cube, with the printed $e^{-5\delta_0r(\ell_k)}$ replaced by $e^{-5\delta_cr(\ell_k)}$); the replacements \cite[(2.108)--(2.115)]{Ba2} by \cref{prop:a01:dGD}(iii), (v), (vi).
\item[{D7 (the collar quartic; \cite[Prop.~2.1, (2.57), (2.79)]{Ba2})}] \Cref{thm:a01:B2} on the propagator domain of \cite[(2.56)]{Ba2} keeps $-\lambda(\ell_i)\sum_{x\in B^i(\mathrm{collar})}\vth^d|\phi^{(i)}(x)|^4$ exactly. The step \cite[(2.79)]{Ba2} rests on \cite[(2.65)]{Ba2}, i.e.\ on $|\phi^{(i)}(x)|\ge\frac12|\phi_i(y)|$ for $x\in B^i(y)$, $y\in R_s^{(i)}$. (b) \emph{The near layers:} for $y\in R_s^{(i)}$ within $\ell^\infty$-distance $R_0$ of the face layer, \cref{lem:a01:F541}(e) supplies a witness (a gradient violation or a face-family element) within distance $dR_0+1$, paid by its own family. (c) \emph{All other blocks:} for $y$ at distance $>R_0$ from $\DD$, the background on $B^i(y)$ differs from the bulk background by the image effect, bounded by $C'e^{-\delta_1(R_0-2)}M_i\le\frac14|\phi_i(y)|$ on the box $\Omega_b\defeq\Xl\cap(\text{the }\ell^\infty\text{ block ball of radius }\lfloor\dep(x)/2\rfloor\text{ about }y)$ (\cref{lem:a01:FL}, Step~1's box-versus-free factor at $A=0$ and Step~3 at regular $A$) plus the box-versus-domain difference (a cut-off and commutator argument at the Neumann faces of $\Omega_b$, as in \cref{thm:a01:B2}(1), of size $CM_ie^{-(\delta_2/2)(\dep(x)/2-8)}\le p_i$ at the D7 points, which lie several $r$-layers inside the propagator domain), with $R_0$ the smallest integer $\ge3$ with $C'e^{-\delta_1(R_0-2)}\le\frac14$; $R_0$ depends on $(d,L,a,\Wt,c_V)$ only and is fixed before Ba{\l}aban's $M$ ($s_b=LM\ge dR_0+L+1$). Under the caps $(C_B+3)\lambda(\ell_i)^{1/4}\le\frac18$, $C'(V_i+\theta_i)\le\frac18$, this gives $|\phi^{(i)}(x)|\ge\frac12|\phi_i(y)|$, hence \cite[(2.79)]{Ba2} with $\frac12\to\frac1{16}$, and the exponent of $R_s$ in \cite[(3.41)]{Ba2} changes only in its constant.
\item[{D8 (the backgrounds in the regular and fine-field classes; \cite[Lemmas~2.3--2.4, Prop.~2.2]{Ba2})}] $A^{(i)}$: \cref{lem:a01:PE}(a), (d) and \cref{lem:a01:REG}. $\phi^{(i)}$: \cite[Prop.~2.2]{Ba2} is \cref{lem:a01:A5} (every bond, face bonds included); the fine-field class \cite[(2.55)]{Ba2}, $|D_{\tilde A}\phi^{(i)}(b)|\le c_1'p_{i-1}$ on $B^i(\Lambda_7^{(i)})$, is \cref{prop:a01:legsd} with the cap $|\phi_i(y_0)|\le\sigma_i+\dist(y_0,\Lone^{(i)})p_i$ for $\dist\le r_i$ (transport along a block path to the face layer inside $\Lambda_7^{(i-1)\prime}\setminus S_i$, since $\dist(\Lambda_7^{(i)},X_i)>7r_i$) and $e^{-\delta_1r_i}M_i\le p_i$ for $\dist>r_i$; the product $(\sigma_i+tp_i)e^{-\delta_1t}\le Cp_i$. At the collar cubes $\square\subset B^i(\Lambda_7^{(i-1)\prime})$ of \cite[Prop.~2.1]{Ba2}, where the data of \cite[(2.56)]{Ba2} are cut one previous-scale layer away and the scale-$i$ tests are known only off the bad set, the cap is inherited from step $i-1$ (\cref{lem:a01:IC}), inserted as a characteristic function of $\phi_i$ as \cite[(2.55)]{Ba2} is carried to step $i$; with it, for every fine bond $b$ within $1$ of $\square$: $|D_A\phi^{(i)}_\Omega(b)|\le C_Gp_i$ and $|\phi_i(y)-(Q\phi^{(i)}_\Omega)(y)|\le C_Kp_i$ at the block $y$ of $\square$, by \cref{prop:a01:legsd} on the King box of radius $\lfloor m_i/2\rfloor$ about $\square$ inside the propagator domain ($m_i\defeq r(\ell_{i-1})/L\ge70$), the truncation entering as the tail $Cc_1M_ie^{-\delta_cr(\ell_{i-1})/L}\le p_i$ and the box-versus-domain difference as in D7(c); the block clause from \cref{lem:a01:EB}. Every use of \cite[(2.55)]{Ba2} in \cite{Ba2} is polynomial in $p$ and paid by factors $e^{-cp^2}$.
\item[{D9 (fluctuation fields small; \cite[Lemmas~2.5, 2.7]{Ba2})}] $|A'|\le O(1)p$ \cite[(2.86)]{Ba2} by \cref{lem:a01:PE}; $|\phi'|\le O(1)p$ \cite[(2.114)]{Ba2} from the block tests and the covariance bounds (\cref{lem:a01:A23}(c), whose lower bound at the weight is \cref{lem:a01:329} plus the coarse block Poincar\'e inequality); near the face \cref{lem:a01:F541}(b)--(d) with the corrected threshold $r(\ell_i)$ (the contour of \cite[(5.44)]{K2} must stay where the gradient test holds, which a path of length $\le r_i$ from $x\in\Lambda_4^{(i)}$ does, since $\dist(\Lambda_4^{(i)},\Lambda_0^{(i)c})\ge4r_i$). \Cref{lem:a01:F541}(d) supplies \cite[(2.113)]{Ba2} at distance $>r_i$ from the face.
\item[{D10 (the $Z$-factors and the tail removal; \cite[(2.92)--(2.109)]{Ba2})}] $\det(C^{(i)})^{\pm1/2}$ with $C^{(i)}$ bounded (\cref{lem:a01:A23}); the removal of the tail field $\tilde B^{(i)}$ of \cite[(2.89)--(2.90)]{Ba2} from the $Z$-factors, the lines and legs and the quadratic form costs $C\ell_i^{\kappa'}|\Lambda^{(i)}|$: the matter form by the envelope argument at the regular backgrounds, the determinant by the exact block split $\ln\det\Mop=\ln\det(-\Delta_A+V)+\ln\det(1+a_iQ(-\Delta_A+V)^{-1}Q^*)$ with the coarse determinant bounded in summed form and the block-free one by its closed-loop representation at the two regular backgrounds, and the perturbative part by \cref{thm:a01:B2} with the tail as an external photon field; on sub-domain histories the tail lies at distance $>r(\ell_i)/2$ from every place where the domain operators differ from the cube operators, so the domain remainders are $O(\ell_i^\kappa)e^{-cr(\ell_i)}$ times the volume.
\item[{D11 (the last step; \cite[(3.11)--(3.12)]{Ba2})}] $G'_k\ge\gamma_1\lmu^2\ell_k^2$ (photon, \cref{lem:a01:PE}(e)); $G''_k\ge\gamma_1m^2\ell_k^2$ on $\Lambda_5^{(k)}$ from $V\ge m^2$ and the $\ell^2$ bounds of \cref{lem:a01:A23}, the Dirichlet terms only adding; $\norm{H'_k},\norm{H''_k}\le O(1)e^{-\delta_1r(\ell_k)}$ requires only $\ell^2$ bounds (\cref{lem:a01:A23}(a)).
\item[{D12 (positivity; \cite[Prop.~3.1, (3.26)]{Ba2})}] For every configuration in the support of the characteristic functions of \cite[(3.21)]{Ba2} (the fine-field class on each collar), in the multiscale form \cite[(3.23)--(3.25)]{Ba2}: \cite[(3.26)]{Ba2} on the $(1-\theta_2)$-share follows from the Neumann separation \cite[(3.27)--(3.28)]{Ba2} (printed algebra, weight-free) and \cref{lem:a01:329} (the per-scale lower bound \cite[(3.29)]{Ba2} at the weight, whose $\phi$-independent error $O(\ell_k^{\kappa_0})|\Lambda_k|$ is carried on that support exactly as in print), with the Dirichlet faces only adding; the photon exponents are unchanged (\cref{lem:a01:PE}(e)); the face family is paid on the $\theta_2$-share by the composite clause of \cref{lem:a01:FPC}, on blocks of the index set, and payments of different scales are on disjoint fine sets and add. The convergence factors \cite[(3.31)--(3.34)]{Ba2} follow with $(1-\theta_2)$ on the two matter positivity exponents, the photon exponents unchanged, plus $\exp(-\frac14p_i^2|\text{face violations at scale }i|)$ at every scale.
\item[{D13 (integration of the block fields; \cite[(3.32)--(3.38)]{Ba2})}] The normalisation identities are the minimiser identity; the exceptional integrals over $|\phi_{i+1}(y)-(Q\phi_i)(y)|>p$ as printed; after the rescaling, \cite[(3.37)]{Ba2} is $\int d\phi_0\exp(-\frac12\ip{\phi_0}{(-\Delta^D_{\tilde A}+V)\phi_0})$ with the Dirichlet Laplacian and $V\ge m^2>0$, and the diamagnetic inequality of \cite[Sect.~3.1, (3.4)--(3.7)]{BFS1} in its site-diagonal form (the resolvent Neumann series with the site-diagonal $\vth d_x+\vth^3V(x)+$ the face terms, every walk weight non-negative and $\norm{U(\tilde A)}=1$, so $|(-\Delta^D_{\tilde A}+V)^{-1}(x,y)|\le(-\Delta^D+V)^{-1}(x,y)$, and the interpolation of \cite[Sect.~3.3, (3.16)--(3.17)]{BFS1}) gives $\le\int\exp(-\frac12\ip\phi{(-\Delta^D+V)\phi})=e^{E_{0,s}}$.
\item[{D14 (combinatorics; \cite[Sect.~3C, (3.51)--(3.54)]{Ba2})}] Purely combinatorial: $|\Lambda_0^{(i)c}|\le O(1)r(\ell_i)^d(|\mathcal C_0|+\dots+|\mathcal C_i|)$ and $-\sum_i(c_0'p_i^2-O(1)r(\ell_i)^{d+1})|\mathcal C_i|\le0$ if $2p>(d+1)\varrho$ (here $12>8$) and $b_0$ is large, with the product partition into boxes of large blocks in place of \cite[(3.43)--(3.45)]{Ba2} and the cube count $|\mathcal C_i|\le5^d|F_i^{\mathrm{paid}}|$ of \cref{lem:a01:PP5}, the paid part of the family comprising Ba{\l}aban's six families ($R_s$ restricted to its far elements) and the face family (the near $R_s$ elements are charged to their witnesses of D7(b) and not counted); the constant becomes $c_0'\defeq5^{-d}\min\{\frac18(1-\theta_2)a,\frac14(1-\theta_2)\gamma_0^w,\frac14,\frac1{16}\}$, the region entropy $2^{7|\Lambda_0^{(i)c}|}$ is absorbed in the $O(1)(1+\ln\ell_i^{-1})|\Lambda_0^{(i)c}|$ of \cite[(3.47)]{Ba2}. We use half of each payment for \cite[(3.54)]{Ba2} and keep $e^{-\frac12c_0'p_i^2|\mathcal C_i|}$.
\item[{D15 (final step; \cite[(2.116)--(2.118)]{Ba2})}] $P^{(K_f)}\le O(\ell_f^{\beta_0})|\Lambda_7^{(K_f)}|$, $\sum_{i\le K_f}O(\ell_i^{\kappa'})|X^{(i)}|\le O(1)|\CJ|$ (geometric, $\kappa'>d$), and the final integral with the region entropy gives $\exp(O(1)|X^{(K_f)}|)$, $|X^{(K_f)}|=\ell_f^{-d}|\CJ|$: \cref{prop:a01:TD}(c).
\end{description}

\begin{proposition}[upper line]\label{prop:a01:upper}
Under the hypotheses of \cref{prop:a01:lower}, $\Zk(s)\le e^{a}\Ik j(s)+e^{-\frac12c_0'p(\eta)^2+E_+|\CJ|}$, with $a$ as there, $c_0'$ the constant of D14 and $E_+$ depending on $\eps_0^{\mathrm{Ba}}$ and King's data only.
\end{proposition}
\begin{proof}
In the inductive formula \cite[(2.43)]{Ba2} at scale $j$ (seven families at every scale $i\le j$, the gradient family at scale $j$ imposed on every block bond of the cube), the term with no large-field block at any scale $i\le j$ is $\int\chi_j\lambda_j\,e^{-S^{(j)}_k[A_\theta]+\sum_{i<j}O(1)\ell_i^{\kappa'}|X^{(i)}|}$ (D1--D10), with Ba{\l}aban's cut background $A_\theta=a_j\ell_j^{-2}\zeta^{(j)}_\theta G_jQ^*A_j$ in every scalar expression \cite[(2.44), (2.50)]{Ba2}; the difference $|S^{(j)}_k[A^{(j)}]-S^{(j)}_k[A_\theta]|\le C\eta^{s'}|X^{(j)}|$ on $\supp\chi_j$ (the matter form by the envelope argument between two regular backgrounds; the determinant by the block split of D10 at the two backgrounds; the perturbative part by \cref{prop:a01:TD} with the background-type photon field $A^{(j)}-A_\theta$), so the term is $\le e^{a}\Ik j(s)$ after enlarging $C$ in $a$. Every other term has a large-field block at some scale $i\le j$ and carries the kept factor $e^{-\frac12c_0'p_i^2}\le e^{-\frac12c_0'p(\eta)^2}$ ($p$ is decreasing in the scale). Continue the RG for these terms to the final scale $\ell_f\asymp\eps_0^{\mathrm{Ba}}$ (D11--D15, with the face family imposed and paid at each step, D2); the sum over all admissible histories with the remaining half of the payments and \cite[(3.54)]{Ba2} is $\le e^{O(1)|X^{(K_f)}|}=e^{E_+|\CJ|}$, $E_+=O(\ell_f^{-d})$ (\cite[(2.117), (3.42)--(3.47)]{Ba2} with (3.54)). The constants are those of D1--D15: independent of $k$, $j$ and $\vth$, each step's constants being so (the payments $e^{-p_i^2/4}$ are exact for every scale by the choice of $\sigma_i$; at every scale the lattice regime is the only regime on the cube, and \cref{lem:a01:DV,lem:a01:A5} cover it uniformly in $\vth\le1$).
\end{proof}


% a01_H0rel part 7: assembly, normalisation-free form, the clause (R_0)
\subsection{Assembly, the normalisation-free form, and the clause \Rzero}\label{app:a01:assembly}

\begin{lemma}[field-independent constants cancel]\label{lem:a01:cnum}
Let $Z'(s)$, $I'(s)$ be $\Zk(s)$ and $\Ik j(s)$ computed with the effective action shifted by a field- and source-independent constant $c$ (a different $E_1$, a different normalisation of the measure, the domain-adapted $E_0$, $E_1$ of \cite[(3.35)]{K1}). If $e^{-a}I'\le Z'\le e^aI'+bZ'(0)$, then $e^{-a}\Ik j\le\Zk\le e^a\Ik j+b\Zk(0)$ with the same $a,b$.
\end{lemma}
\begin{proof}
$Z'=e^{-c}\Zk$ and $I'=e^{-c}\Ik j$ for every $s$; divide.
\end{proof}

\begin{proof}[Proof of \cref{prop:k01:H0rel}]
Fix $\eta_1\le\min(5^{-J_*},\eps_0^{\mathrm{Ba}})$ and smaller where the conditions on $\eta_1$ collected above require it (\cref{lem:a01:GC,lem:a01:final,prop:a01:legsd}, D7, D14), so that \cite[(3.54)]{Ba1} holds at every scale $\le j$ for $b_0\ge b_0^*$; the condition $j\ge J_*$ is Ba{\l}aban's $\eps\le\eps_0(\eta)=\eta5^{-J_*}$ and $j\le k-J_*$ is King's $\eta\le5^{-J_*}$, so that $j\le K_f\le k$. \Cref{prop:a01:lower,prop:a01:upper} give, in the normalisation of B6,
\[
e^{-a}\Ik j(s)\le\Zk(s)\le e^{a}\Ik j(s)+e^{-\frac12c_0'p(\eta)^2+E_+|\CJ|},\qquad\Zk(0)\ge e^{-E_-|\CJ|},
\]
hence the remainder is $\le e^{-\frac12c_0'p(\eta)^2+(E_++E_-)|\CJ|}\Zk(0)$. By \cref{lem:a01:cnum} the same holds for $\Zk$ and $\Ik j$ in the normalisation of \cref{ssec:k01:setting}. This is \eqref{eq:k01:H0rel} with $C\eta^\sigma$ in place of $C\eta^\sigma|\CJ|$ (the volume $|\CJ|=(2MJ)^3$ is absorbed in $C$, which may depend on $J$), $c=\frac12c_0'$ and $\rho(\eta)=e^{-cp(\eta)^2+C}$ with $C\ge(E_++E_-)|\CJ|$. The constants are independent of $k$ and $j$: of $j$ (i.e.\ of $\eta\le\eta_1$) because every step's constants are, and of $k$ (i.e.\ of $\vth$) because no step depends on the spacing except through the lattice regime, which \cref{lem:a01:DV,lem:a01:DVloc,lem:a01:A5} cover with the same constants; the dependence on $x$ enters through the source vertices only (B7, \cref{prop:a01:TD}). The clause \Rzero\ is \cref{thm:a01:R0} below.
\end{proof}

\begin{theorem}[the remainder is the large-field part of Ba{\l}aban's decomposition]\label{thm:a01:R0}
In the setting of \cref{prop:k01:H0rel}, let $P\defeq\prod_{i\le j}P^{(i)}$, where $P^{(i)}$ is the product, over all points, bonds and blocks of the scale-$i$ index sets of the all-small history (the whole cube's lattices at every scale), of the small-field indicators of the families of scale $i$ (Ba{\l}aban's six families, the face family, and at $i=j$ King's $\chi_j$ and $\lambda_j$). Then: (a) $P$ is a function on the joint variable space of Ba{\l}aban's block-spin representation with values in $\{0,1\}$, the same for every source $s$; (b) $\int P\,d\mu_{k,s}$ is the exact contribution of the all-small history at scale $j$, and $\int P\,d\mu_{k,0}\le e^{C\eta^\sigma}\Ik j(0)$; (c) $\int(1-P)\,d\mu_{k,0}\le e^{-\frac12c_0'p(\eta)^2+E_+|\CJ|}\le\rho(\eta)\Zk(0)$.
\end{theorem}
\begin{proof}
\emph{The joint measure.} Each block-spin step inserts $1=N_i^{-1}\int dB\,d\psi\exp(-\frac a2|B-QA|^2-\frac a2|\psi-Q(A)\phi|^2)$ \cite[(2.1)]{Ba2}, with $N_i$ field-independent; so $\Zk(s)=N^{-1}\int d\mu_{k,s}$ with $d\mu_{k,s}\defeq e^{-S}e^{X(s)}\prod_i[\text{kernels}]\,d(\text{all variables})$ a positive measure, and $N$ independent of $s$. \emph{The tested quantities are source-free:} they are fields, block fields and fluctuation fields (differences between a field and a background, \cite[(3.16)]{K1}), and backgrounds are linear maps of block fields built from the source-free quadratic forms; the source enters only the vertices \cite[(3.20)]{K1} and the factor $e^{X(s)}$. This gives (a). \emph{Decomposition at one scale.} Expanding $\prod[\chi(\text{small})+\chi(\text{large})]$ over the points, bonds and blocks of scale $i$ gives \cite[(2.4)--(2.6)]{Ba2}: $1=$ a sum over large-field tuples of products of indicators; the regrouping by $\Lambda_0$ and the partial resummation \cite[(2.9)--(2.10)]{Ba2} replace, for fixed $\Lambda_0\ne\Xl$, the indicator of ``the true tuple generates $\Lambda_0$'' by a sum over minimal admissible tuples of non-negative indicator products, which is $\ge$ the indicator it replaces; for $\Lambda_0=\Xl$ the only admissible tuple is the empty one and its term is exactly the product of all small-field indicators of that scale. \emph{Induction over scales.} In each history the next scale's decomposition is inserted where the RG continues; inside a term the translations are changes of variables with unit Jacobian, the Gaussian integrations are exact, and the removal of interactions from large-field regions and the large-field estimates \cite[(2.11)--(2.14)]{Ba2} are termwise upper bounds of a positive integrand; this is why \cite[(2.43)]{Ba2} is an inequality bounding $\Zk$ itself. \emph{Telescoping}, exact and pointwise, with $P_i\defeq\prod_{i'\le i}P^{(i')}$, $P_{-1}\defeq1$: $1-P=\sum_{i=0}^{j}P_{i-1}(1-P^{(i)})$, the $i$-th summand being the indicator of ``all tests of scales $<i$ pass everywhere, some test of scale $i$ fails''; by the identity-or-over-count at each scale, $P_{i-1}(1-P^{(i)})\le$ the sum of the weights $w_H$ of the histories $H$ whose first large-field block is at scale $i$, and summing over $i$, $1-P\le\sum_{H\ne H_0}w_H$, $H_0$ the all-small history (weight $P$). Integrate against the positive measure $d\mu_{k,0}$: $N^{-1}\int(1-P)d\mu_{k,0}\le\sum_{H\ne H_0}[\text{the (2.43)-term of }H\text{ continued to the final scale}]$, each history's exact contribution being bounded by its term because every operation after the insertion of the indicators is an identity or a termwise upper bound. The proof of \cref{prop:a01:upper} bounds exactly this sum by $e^{-\frac12c_0'p(\eta)^2+E_+|\CJ|}$, and $\Zk(0)\ge e^{-E_-|\CJ|}$ in the same normalisation (B6), which gives (c). (b): the all-small history's exact contribution is $\int P\,d\mu_{k,s}$, which the proof of \cref{prop:a01:upper} bounds by $e^{a}\Ik j(s)$. The normalisation-free form follows by \cref{lem:a01:cnum}. Positivity of $(1-P)\,d\mu_{k,0}$ holds because $P\le1$; whether \cite[(2.9)]{Ba2} is an identity or an over-count does not matter, only $1-P\le\sum_{H\ne H_0}w_H$ is used.
\end{proof}

